\documentclass[11pt,a4paper]{article}

\usepackage{fontspec}
\setmainfont{Libertinus Serif}
\setsansfont{Libertinus Sans}
\setmonofont{Noto Sans Mono}[
  Scale=MatchLowercase,
  Weight=400,
  BoldFont={Noto Sans Mono},
  BoldFeatures={Weight=700}
]
\usepackage[english]{babel}
\usepackage[a4paper,margin=2.4cm]{geometry}
\usepackage[skip=6pt plus 1pt,indent=0pt]{parskip}
\setlength{\emergencystretch}{1em}
\usepackage{xcolor}
\usepackage{listings}
\usepackage{hyperref}
\usepackage{longtable}
\usepackage{array}
\usepackage{enumitem}
\setcounter{tocdepth}{2}

\definecolor{abclblue}{HTML}{163A5F}
\definecolor{abclgray}{HTML}{F4F4F4}
\hypersetup{
  colorlinks=true,linkcolor=abclblue,urlcolor=abclblue,citecolor=abclblue,
  pdftitle={Programming in ABCL/R2: A Guided Introduction to Objects, Concurrency, and Reflection},
  pdfsubject={Objects, messages, concurrency, and reflection in ABCL/R2},
  pdflang={en}
}
\lstset{
  basicstyle=\ttfamily\small,
  backgroundcolor=\color{abclgray},
  frame=single,
  rulecolor=\color{black!20},
  columns=fullflexible,
  keepspaces=true,
  showstringspaces=false,
  breaklines=true,
  aboveskip=\bigskipamount,
  belowskip=\medskipamount
}
\newcommand{\code}[1]{\texttt{#1}}
\newcommand{\filepath}[1]{\path{#1}}
% Source excerpts are prepared by the external SBCL build step, not by TeX.
\input{build/tutorial-build/listings.tex}
\newcommand{\sourcelisting}[2][]{%
  \edef\preparedlisting{\csname listing:\detokenize{#2|#1}\endcsname}%
  \lstinputlisting[#1]{\preparedlisting}}
% Keep short listings together; put spacing outside the minipage because
% ordinary vertical glue at the top of a minipage is discarded. The
% non-discardable inner padding also leaves room for the listing's frame.
\newenvironment{codeblock}{%
  \par\addvspace{\medskipamount}\noindent
  \begin{minipage}{\linewidth}%
  \vspace*{\medskipamount}%
  \lstset{aboveskip=0pt,belowskip=0pt}%
}{%
  \end{minipage}\par\addvspace{\medskipamount}}
\newcommand{\wholelisting}[2][]{%
  \begin{codeblock}%
  \sourcelisting[#1]{#2}%
  \end{codeblock}}
\newcolumntype{L}[1]{>{\raggedright\arraybackslash}p{#1}}

\title{Programming in ABCL/R2\\[0.5em]
  \large A Guided Introduction to Objects,\\
  Concurrency, and Reflection}
\author{For the September 4, 1998 distribution}
\date{\today}

\begin{document}
\maketitle
\begin{abstract}
This book takes you from your first interactive session to concurrent,
reflective objects in ABCL/R2. Each major concept is introduced through a
small example, with exercises to help you apply what you have learned.
You need experience with object-oriented programming and some Lisp, but
not Common Lisp, actors, or reflection. The goal is to read and explain
the original programs in the accompanying distribution, not just run
the additional teaching examples.
\end{abstract}
\paragraph{How to use this book.}
Read the opening chapters in order, including the Lisp reading guide.
Then work through futures, reflection, and the original-source reading projects.
The original examples form the thread connecting concepts to complete programs.
They live in \filepath{examples/}; experiment on copies, not the reference sources.
Executable programs live in \filepath{tutorial/examples/}.
The later sections contain worked exercises and a complete job-service solution.
Start from the repository root unless a command says otherwise.
\begingroup
\linespread{0.90}
\small
\setlength{\parskip}{0pt}
\tableofcontents
\endgroup
\newpage

\section{What Is ABCL/R2?}

ABCL/R2 is an object-oriented, concurrent, reflective programming language.
A program consists primarily of objects that receive messages and execute
scripts. Concurrency is part of the language: messages are placed in object
queues, and computations in different objects can make progress independently.

The accompanying distribution
(documented as version 0.60, September 1998) implements this model using
partial compilation. The examples here use its SBCL implementation.
We will unpack these implementation terms in the reflection chapter.
Masuhara's 1992 thesis supplies the historical model~\cite{masuhara}.

\paragraph{Thinking in objects.}
Start from the familiar idea of an object with state and behavior:
\begin{enumerate}[nosep]
  \item It owns state.
  \item It receives messages and matches them against patterns.
  \item A matching script is selected for execution.
  \item The script can compute, change state, create objects, and send
        further messages.
\end{enumerate}

\subsection{From Method Calls to Message Protocols}

In a familiar OOP language, a call such as \code{account.balance()}
selects a method, executes its body, and returns a value to the caller.
ABCL/R2 separates these jobs. A sender constructs a data value, sends it
to an object, and may arrange a destination for a later reply.
The receiving object matches that data against its scripts.
A \emph{script} is a handler: a message pattern followed by expressions
to execute when the pattern matches.

\begin{longtable}{L{0.28\textwidth}L{0.62\textwidth}}
\textbf{Familiar idea} & \textbf{ABCL/R2 counterpart and difference}\\\hline
Object fields & State bindings belonging to an object; outer lexical
bindings can also be shared.\\
Method selection & First matching message pattern, not necessarily a
method name plus a fixed argument signature.\\
Method arguments & Data inside the message; pattern variables extract
the parts the handler needs.\\
Return value & An explicit reply message, not the handler's final value.\\
Inheritance & Can be implemented by a reflective protocol; the original
bird/emu example does not declare a superclass.\\
\end{longtable}

A \emph{protocol} specifies the messages an object accepts and what its
clients may expect: replies, state changes, or further messages.
For an account, \code{[:deposit 10]} and \code{:balance} can be two
protocol messages. The colon is part of a Lisp keyword, not special
method-call punctuation; the next chapter explains keywords.

\subsection{An Actor-Like Mental Model}

An \emph{actor} is a useful introductory model of an independently
addressable participant that reacts to messages and can send more messages.
For an ordinary ABCL/R2 object, imagine an inbox and a handler processing
one request at a time. Other objects have their own work.
\emph{Concurrent} means their activities can be outstanding together;
it does not require simultaneous execution on different CPUs.

The distinction becomes concrete in this sequence:
\begin{enumerate}[nosep]
  \item A client sends a deposit request to an account.
  \item The request enters the ordinary receiving object's message machinery.
  \item The account selects a matching handler and updates its balance.
  \item If the protocol requires it, the handler explicitly sends a reply.
  \item The client uses that reply, or continues without waiting if it sent
        asynchronously.
\end{enumerate}
Sending is therefore not the same event as finishing the requested work.
An asynchronous sender does not automatically receive a completion notice.
Conversely, a waiting sender cannot finish its request unless someone
sends the required reply.

This is an actor-like explanation, not a claim that ABCL/R2 enforces pure
actor isolation. Its lexical bindings can be shared between objects.
\emph{Reflection} means that a program can inspect or control parts of
its own execution machinery. It can therefore change the ordinary
selection and execution of handlers.
Learn the ordinary model first; later chapters state where it is extended.

\subsection{Three Questions to Ask While Reading}

For every send in an original program, ask:
\begin{enumerate}[nosep]
  \item Which object receives this message, and which pattern matches it?
  \item Does the sender wait, and where will any reply be sent?
  \item What establishes completion: one reply, a count, or a terminal message?
\end{enumerate}
In \filepath{examples/misc/future-ex.abclr2}, these questions distinguish
the tree walker, the object consuming its replies, and the printer.
In \filepath{examples/misc/inherit.abclr2}, they distinguish an application
object from the metaobject running its dispatch protocol.
You do not need to understand those files yet; these are our first
original-source reading goals.

\paragraph{A route through the book.}
Work in stages rather than trying to absorb every library at once:
\begin{enumerate}[nosep]
  \item \emph{Build the ordinary model.} Follow
        Sections~\ref{sec:first-session}--\ref{sec:files}, trying the
        small sessions and exercises. At the bank account
        (Section~\ref{sec:account}), explain which pattern handles each
        request and why the caller receives a reply.
  \item \emph{Follow results over time.} Study futures in
        Section~\ref{sec:futures}, then trace the original tree walker
        in Section~\ref{sec:original-future}. Explain why its terminal
        message, not an empty buffer or a printed line, marks completion.
        Express messages are an advanced extension; you may return to
        that experiment after learning reflection.
  \item \emph{Open the execution machinery.} Read
        Section~\ref{sec:reflection}, then the original lazy-control,
        inheritance, and interpreter projects in
        Section~\ref{sec:originals}. Keep the application reply separate
        from the interpreter's execution-completion message.
  \item \emph{Combine the mechanisms.} Work through the library and
        debugging projects, then Pcontrol and Time Warp
        (Sections~\ref{sec:pcontrol} and \ref{sec:timewarp}).
        These are advanced reading projects, not prerequisites for
        ordinary object programming. Return to the file map in
        Appendix~\ref{sec:source-map} when following their helper objects.
\end{enumerate}
The job-service project (Section~\ref{sec:project}) is a parallel practice
track: it combines ordinary objects and reply protocols, so you can try
it after the foundations without first mastering the reflective projects.
Use the worked exercise answers to check your reasoning, not just the
final values. At each stage, being able to explain a message's recipient,
reply destination, and completion condition is the checkpoint for moving on.


\section{Installation and Your First Session}
\label{sec:first-session}

Build the interpreter and start it from the project directory:

\begin{lstlisting}[language=bash]
make
./src/saved_abclr2
\end{lstlisting}

The prompt accepts ABCL/R2 expressions. Type \code{:q} to quit or
\code{:reset} to restart the session. To make experiments reproducible,
save your program files in a separate directory of your choice.
For example, load a program named \filepath{demo.abclr2} from the session:

\begin{lstlisting}
(inc "demo" :use-source t)
\end{lstlisting}

The Makefile uses SBCL and writes \filepath{src/saved_abclr2}.
Running the saved image does not require installing it system-wide.

Omit the \code{.abclr2} extension when using \code{inc}.
\code{:use-source t} forces translation from the ABCL/R2 source rather
than reuse of an existing translation. Keep this option while trying
the small sessions. Section~\ref{sec:files} covers compilation, libraries,
modules, and loading contexts once objects and messages are familiar.

\subsection{A Lisp Reading Guide}
\label{sec:lisp-reading}

This section teaches the Lisp notation used by the original examples.
Try the small expressions at the ABCL/R2 prompt. Comments beginning with
a semicolon explain the expected result; do not type the prompt itself.
The guide is deliberately selective: it is not a Common Lisp course.

\subsubsection{Expressions, Symbols, and Literal Data}

An expression such as \code{(+ 2 3)} puts the operation first.
For an ordinary function call, argument expressions are evaluated before
the function is called. Numbers and strings are literal values.
An unquoted ordinary symbol in an expression names a binding:
\begin{lstlisting}
(+ 2 (* 3 4))                  ; => 14
(let ((amount 10)) (+ amount 2)) ; => 12
\end{lstlisting}
The inner multiplication finishes before the outer addition.
\code{let} introduces a local binding; its body can refer to \code{amount}.
A \emph{binding} associates a name with a value. Its \emph{lexical scope}
is determined by the surrounding source structure: the body of this
\code{let} is where its local name is available.

A symbol is not a string. \code{amount} is a name, \code{'amount} is that
name as data, and \code{"amount"} is a string of characters.
The quote character is shorthand for \code{quote}:
\begin{lstlisting}
'amount                       ; => AMOUNT
'(a b c)                      ; => (A B C)
(list 'a 'b 'c)               ; => (A B C)
\end{lstlisting}
Quote prevents evaluation of its operand. Thus the symbols inside
\code{'(a b c)} are not looked up as variables.
Unescaped symbol names are normally printed in uppercase; \code{n} and
\code{N} read as the same symbol in the same package.

A \emph{keyword} begins with a colon and evaluates to itself.
It is convenient as a protocol tag: \code{:deposit}, \code{:new},
and \code{:finish} are distinct symbols. No variable definition is
needed to use them.
\begin{lstlisting}
(list :deposit 10)            ; => (:DEPOSIT 10)
(let ((x 7)) [:point x (+ x 1)]) ; => (:POINT 7 8)
\end{lstlisting}
Here square brackets construct a list of evaluated elements.
They also introduce ABCL/R2 operations such as object creation or sending;
the recognized operator determines which interpretation applies.
In a \emph{pattern}, however, a symbol such as \code{x} binds a part
of the incoming value. Construction and matching are different uses
of the same readable list notation.

\subsubsection{Truth, Tests, and Control Flow}

\code{nil} means false and is also the empty list. \code{t} is true.
Every value other than \code{nil}, including zero, counts as true.
\begin{lstlisting}
(if nil :yes :no)             ; => :NO
(if 0 :yes :no)               ; => :YES
(null '())                    ; => T
(not t)                       ; => NIL
\end{lstlisting}
\code{if} is a \emph{special form}: it does not eagerly evaluate both
branches. It evaluates the test and only the selected branch.
\code{progn} evaluates several expressions in sequence, with the final
expression providing its value. \code{when} is a conditional sequence;
\code{cond} chooses the first clause whose test succeeds.
\code{and} stops at the first false operand; \code{or} stops at the
first true operand. These are control operations, not ordinary eager calls.

Use \code{=} to compare numbers. Use \code{eq} for symbol or object
identity, as in testing a terminal keyword or checking a self-transfer.
Use \code{equal} for structural comparison of lists and strings.
Do not assume that \code{eq} compares two independently constructed
lists by their contents.

\subsubsection{Lists, Cons Cells, and Dotted Tails}

A \emph{cons cell} has two parts: its \code{car} and its \code{cdr}.
A proper list chains cons cells and ends in \code{nil}.
\code{first} and \code{rest} are readable names for the head and tail
of a list; \code{car} and \code{cdr} expose the underlying pair.
\begin{lstlisting}
(cons 'a '(b c))              ; => (A B C)
(car '(a b c))                ; => A
(cdr '(a b c))                ; => (B C)
(cons 'd 'e)                  ; => (D . E)
(cdr '(d . e))                ; => E
\end{lstlisting}
The dot means that the final tail is the value after the dot.
\code{(d . e)} is a pair, not the two-element proper list \code{(d e)}.
Its tail is the symbol \code{e}, not the list \code{(e)}.
\code{(a b)} is shorthand for \code{(a . (b . nil))}.

In a bracket pattern, \code{[A . B]} binds the head to \code{A} and
the entire tail to \code{B}. \code{[[A] . B]} additionally requires
the head to be a one-element proper list. These distinctions drive the
original tree walker; do not replace dotted pairs by ordinary lists.

\code{append} concatenates lists, copying all but the final argument.
\code{nconc} performs destructive concatenation.
\code{length} counts elements of a proper sequence; \code{nth} selects
a zero-based element of a list. Later original programs also use arrays:
\code{make-array} creates an array and \code{aref} accesses an element.
Do not pass an arbitrary dotted pair where a proper list is required.

\subsubsection{Bindings and Updating Values}

\code{let} evaluates its initializers before its new bindings are available
to the body. \code{let*} introduces bindings one at a time, so a later
initializer can refer to an earlier binding.
\begin{lstlisting}
(let* ((x 3) (y (+ x 1))) (* x y)) ; => 12
(let ((x 3)) [x := (+ x 1)] x)    ; => 4
\end{lstlisting}
\code{[x := value]} changes an ABCL/R2 binding.
\code{setq} is also available, but this parser permits only one assignment
pair in a \code{setq} form.
\code{setf} is a host macro for updating a \emph{place}, such as the
\code{cdr} of a pair or an array element. Do not confuse changing a binding
with changing an existing data structure.

An \emph{association list} stores bindings as pairs:
\begin{lstlisting}
(let ((bindings '((n . 3) (balance . 20))))
  (cdr (assoc 'balance bindings))) ; => 20
\end{lstlisting}
\code{assoc} finds the pair whose key matches.
\code{pairlis} pairs a list of names with a list of values.
The original environment and script examples use association lists to
describe bindings; those lists are data, whereas an environment is an
object that implements lookup and update messages.
A local set of bindings for one scope is called a \emph{frame}; an
environment can link its frame to an enclosing environment.

\subsubsection{Reading a Loop}

ABCL/R2 implements \code{do} directly. Read its three parts as local
bindings, a termination clause, and a body:
\begin{lstlisting}
(do ((n 0 (+ n 1)))
    ((= n 3) n)
  (format t "~A~%" n))
; Prints 0, 1, 2 on separate lines; returns 3.
\end{lstlisting}
\code{n} starts at zero. Before each iteration the termination test is
checked. If false, the body runs and the step expression updates \code{n}.
If true, the result expressions of the termination clause run.
With several step expressions, their new values are computed before the
loop variables are updated. A variable with no step expression keeps its
value unless the body changes it.

The tree walker uses \code{(do ((tree received-tree)) ((null tree)) ...)}.
There is no step expression: its \code{match} clauses explicitly update
\code{tree}. Its loop result is \code{nil}, but the handler subsequently
sends \code{!:finish}. The loop result is not that reply.

\subsubsection{Functions, Macros, and Reader Notation}

\code{defun} defines a host Lisp function; its parameter list describes
arguments and its body computes a result.
\code{\#'car} is shorthand for \code{(function car)} and denotes a function,
not the symbol \code{'car}. \code{funcall} calls a function value;
\code{mapcar} applies a function to corresponding list elements:
\begin{lstlisting}
(mapcar #'car '((a . 1) (b . 2))) ; => (A B)
\end{lstlisting}
\code{apply} calls a function using a list for its final arguments.
The priority queues use \code{(apply \#'sort item-list predicate keys)}:
when \code{keys} is \code{(:key \#'first)}, this is a call to
\code{sort} with that keyword pair appended to its arguments.
\code{sort} can rearrange its input list destructively; use its returned
list rather than assuming that the old head still names the result.
The host function is called with values, not ABCL forms to execute later.

\begin{lstlisting}
(apply #'+ '(1 2 3))             ; => 6
(let ((items (list :a :b)))
  (push :new items)
  (list (pop items) items))       ; => (:NEW (:A :B))
\end{lstlisting}

\code{push} changes a place to hold a new cons cell; \code{pop} returns
its first element and changes the place to the remaining tail.
\code{pushnew} inserts only when the item is not already present
under its equality test. These macros update bindings or places;
they are not just pure list constructors. \code{caar}, \code{cadr},
and \code{cddr} abbreviate \code{car} of \code{car}, \code{car} of
\code{cdr}, and \code{cdr} of \code{cdr}, respectively.

\code{lambda} describes an anonymous host function. As explained below,
it does not automatically capture ABCL/R2 state.
A \emph{closure} is a function value together with the lexical bindings
it retains. Host Lisp closures retain host Lisp bindings, not arbitrary
bindings in the ABCL/R2 interpreter's environment objects.

A macro transforms source expressions before the transformed program
is translated or executed. Unlike a function, it receives argument forms,
not their computed values. \code{defmacro} defines such a transformation.
Backquote builds a source template; comma inserts a form into it,
and comma-at splices a list of forms into the surrounding list.
The original inheritance macro is a compact example:
\wholelisting[firstline=14,lastline=15]{examples/misc/inherit.abclr2}
\code{\&rest args} collects all argument forms into a list.
Thus \code{(inherit Any Reply Sender bird)} becomes
\code{(reify-args-to-meta :inherit Any Reply Sender bird)}.
The macro inserts the tag and preserves the argument forms; it does not
itself perform inheritance.
You only need to recognize the transformation here;
Section~\ref{sec:reflection} explains what the reflective operation does.

Other parameter markers you will meet are \code{\&optional} for optional
positional arguments and \code{\&key} for keyword arguments.
\code{\&aux} introduces initialized local variables in a host function
or macro's parameter list. The macros in Copy and Report use it with
\code{gensym}, which creates fresh symbols so generated temporary
bindings do not accidentally capture a name in the caller's source.
In \code{(inc "demo" :use-source t)}, the keyword selects an option
and \code{t} supplies its value.
\code{r2sys::create-primary-eval} names an internal symbol in the runtime
package. A package groups symbol names; the double colon exposes an
internal name rather than promising a portable public API.
\code{\#.} in a source file requests evaluation while reading;
it is different from an ordinary expression executed later.
Do not treat a source file containing it as inert quoted data.
\code{\#\textbackslash{}>} is a character literal, distinct from the
one-character string \code{">"}; Copy's display-name helpers inspect
such characters in strings without changing the object's identity.

\code{\#| ... |\#} encloses a block comment; many original files put
sample sessions there, so loading the file does not run those commands.
\code{\#+feature} includes the next form only if that reader feature is
enabled, and \code{\#-feature} does the opposite. The original sources
use \code{\#+COMMENT} for inactive alternatives in the normal build.
This is a reader condition, not a block comment or an extra runtime
branch: do not combine such alternatives with the active implementation.

\code{format} uses a control string: \code{\~{}A} prints a value readably
for a person, \code{\~{}S} prints Lisp notation, and \code{\~{}\%} emits
a newline. The first argument \code{t} selects standard output.
\code{\~{}D} prints a decimal integer and \code{\~{}\&} begins a fresh
line when needed. An open output stream can replace \code{t}; the
monitoring examples use \code{open} to create their log streams and
\code{close} when recording ends.
These explanations are enough to read the printer in the original
future example without assuming familiarity with Common Lisp formatting.

\paragraph{Checkpoint, with answers.}
What are the values of \code{(cdr '(a b))} and \code{(cdr '(a . b))}?
The first is the list \code{(B)}; the second is the symbol \code{B}.
Why does \code{'n} work as an environment lookup key?
Because the environment needs the variable's name as data, not its value.
Why is the inheritance macro not an ordinary function?
Because it constructs a reflective expression from unevaluated source forms.


\subsection{The Boundary with Common Lisp}

ABCL/R2 uses a Common Lisp host, but its scripts are not arbitrary
Common Lisp programs. The following is a compatibility boundary, not a
list of prerequisites you must already know:

Supported special forms include \code{declare}, \code{function}, \code{if},
\code{let}, \code{let*}, \code{progn}, \code{quote}, and \code{setq}.
The \code{do} form is also implemented as a special form.
This version does not provide the entire Common Lisp language:
\code{tagbody/go}, \code{block/return}, and multiple-value facilities are
unavailable. ABCL forms such as \code{wait-for} and \code{match-loop}
are also unsupported.
The original guide additionally excludes \code{wait-for-loop},
routine calls, the \code{suicide} form, and the brace-delimited form
for sending several synchronous requests together~\cite{guide}.
This does not exclude an ordinary now-type send: \code{<==} still
waits for one request's reply. Use explicit reply protocols or futures
when coordinating several requests.
The contents of \code{declare} are ignored by this implementation;
do not rely on a declaration to validate data.

\paragraph{A boundary with host Lisp.}
A \code{function} form can embed a host Lisp lambda, but that lambda
cannot capture the surrounding ABCL/R2 state or lexical environment.
Its body cannot contain ABCL/R2 message sends or object definitions.
Bindings introduced inside the lambda are ordinary host Lisp bindings.
Use an object script, rather than a host lambda, when you need access
to ABCL/R2 state or communication~\cite{guide}.

\section{Objects, Scripts, and Messages}
\label{sec:objects}

\subsection{A Small Useful Object}

Define an object using \code{[object ...]}. Its behavior is specified in
a \code{script} block. The arrow \code{=>} connects a message pattern
to the body that handles a matching message.

\begin{lstlisting}
[object greeter
  (script
    (=> [:hello name]
        (format t "Hello, ~A!~%" name)))]
\end{lstlisting}

The pattern \code{[:hello name]} binds \code{name} to the argument of
a message tagged \code{:hello}. The object is bound to the name
\code{greeter}. To trigger its greeting, send a message:

\begin{lstlisting}
[greeter <= [:hello "Ada"]]
; Prints: Hello, Ada!
\end{lstlisting}

\subsection{Replying with \code{!}}

The expression \code{!exp} sends the value of \code{exp} to the reply
destination of the current message. Here is a simple calculator:

\begin{lstlisting}
[object calculator
  (script
    (=> [:add x y] !(+ x y))
    (=> [:square x] !(* x x)))]

[calculator <== [:add 2 3]]       ; => 5
[calculator <== [:square 7]]      ; => 49
\end{lstlisting}

The exclamation mark denotes a reply in the message protocol.
Use ordinary Common Lisp functions such as \code{format} for side effects
such as logging or console output.
An ordinary final expression does not automatically reply to the sender:
use \code{!} or send explicitly to the bound reply destination.
Conversely, \code{!} sends a message; it does not end the handler.
Later expressions can still execute, and several \code{!}
forms can send several replies. Use exactly one reply for an ordinary
now-type request; a future stream can deliberately accept several.

\subsection{Message Patterns}

Patterns can contain variables, literals, and nested lists:

\begin{lstlisting}
[object router
  (script
    (=> [:route :urgent text]
        (format t "URGENT: ~A~%" text))
    (=> [:route destination text]
        (format t "to ~A: ~A~%" destination text)))]
\end{lstlisting}

The first script matches the literal \code{:urgent}; the second binds
\code{destination} and \code{text}. In this distribution, the
\code{match} form does not support \code{where} constraints.
For constraints on \emph{script patterns}, load the \code{heavy-meta}
library before defining the object. These are distinct facilities:
loading the library does not add constraints to \code{match}.
Section~\ref{sec:guards} follows the original counter and explains the
required metaobject setup, including explicit groups.

The ordinary metaobject selects the first matching script in its script
list. Put specific patterns before catch-all patterns, as the original
bird/emu example in \filepath{examples/misc/inherit.abclr2} does.
A message with no matching script is warned about and discarded;
it is not retained until a suitable script appears.
A now-type caller receives no automatic failure reply in that case.
Check message tags and shapes before assuming that a wait will finish.

A header can also bind the sender separately from the payload and reply
destination: \code{(=> pattern @ R from S body)}.
Here \code{R} is where the reply should go and \code{S} is who sent the
request; forwarding can make those different objects.
The original inheritance example uses both fields.

\subsection{Parentheses, Brackets, and Quoting}

Parentheses normally express a Lisp function call:
\code{(+ 2 3)} calls addition. Brackets have several roles.
A recognized operator makes them an ABCL form, such as object creation,
assignment, or message transmission. Otherwise, they construct a list
of evaluated values. Thus \code{[:add 2 3]} constructs the data
\code{(:ADD 2 3)}; it does not call an addition method by itself.
The receiver's script gives that data meaning.

\begin{lstlisting}
(let ((x 7))
  [:point x (+ x 1)])     ; => (:POINT 7 8)
'(red green blue)        ; Literal symbols, not variable lookups
[:label "red"]           ; A keyword and a string
\end{lstlisting}

A keyword such as \code{:point} evaluates to itself. An ordinary symbol
such as \code{x} denotes a variable in an expression, but becomes a
binding variable in a pattern. The same bracket syntax therefore serves
as construction at the sending site and decomposition at the receiving site.
Semicolons introduce comments to the end of a line.

\subsection{Matching Data Inside a Script}

Use \code{match} when an already selected handler needs to inspect a value.
Clauses use \code{is}, and a final \code{otherwise} clause provides a fallback.
The wildcard \code{\_} matches a value without introducing a named binding.
A dotted pattern separates a list's head from its remaining tail.

\sourcelisting[firstline=14,lastline=21]{tutorial/examples/01-basics.abclr2}

\begin{lstlisting}
[classifier <== [:classify [:point 2 3]]]
; => (:POINT 2 3)
[classifier <== [:classify [1 2 3]]]
; => (:NONEMPTY 1 (2 3))
[classifier <== [:classify nil]]   ; => :EMPTY
[classifier <== [:classify 9]]     ; => :ATOM
\end{lstlisting}

The point case must precede the broad head/tail case in this \code{match}:
a point is also a nonempty list. Pattern matching tests structure;
it does not validate that coordinates are numbers. Use \code{if} in
the selected body for such a check.

\subsection{Temporary Variables and Loops}

An object's \code{state} survives successive messages.
A \code{temporary} declaration at the beginning of a handler instead
creates local variables for that invocation. Here the accumulator starts
at zero every time:

\sourcelisting[firstline=23,lastline=30]{tutorial/examples/01-basics.abclr2}

\begin{lstlisting}
[list-sum <== [:sum [2 3 4]]]   ; => 9
[list-sum <== [:sum nil]]       ; => 0
\end{lstlisting}

The \code{do} loop binds \code{rest}, tests for termination, runs its body,
and updates \code{rest} using \code{cdr}. You can also use \code{let}
for a local accumulator. This distribution provides \code{dotimes}
and \code{dolist} macros that expand into supported \code{do} forms.
Do not assume that every host Lisp macro expands into forms the ABCL
parser understands.

\paragraph{Exercise 1.}
Define \code{sum-of-squares} with a \code{[:sum values]} handler. Reuse the loop structure,
but add the square of each number. Check that an empty list returns zero.
Compare your reasoning with Section~\ref{sol:squares}.


\section{State and Assignment}
\label{sec:state}

A state entry has the form \code{(state [variable := initial-value])}.
Assignment uses \code{[variable := expression]}.

\begin{lstlisting}
[object counter
  (state [n := 0])
  (script
    (=> :inc
        [n := (+ n 1)]
        !n)
    (=> :get !n)
    (=> [:add amount]
        [n := (+ n amount)]
        !n))]

[counter <== :inc]         ; => 1
[counter <== [:add 4]]     ; => 5
[counter <== :get]         ; => 5
\end{lstlisting}

State variables are directly accessible inside a script.
Expressions in a script body execute in sequence.
The example also shows that a message pattern can consist of a single keyword.

\subsection{Lexical Variables}

ABCL/R2 shares lexical variables among objects created within their scope.
This differs from ABCL/1: a change to a shared variable can become visible
to other objects. With concurrent accesses, do not assume a particular
execution order. Encapsulate shared state in an object responsible for
handling access to it.

\section{Two Kinds of Messages}
\label{sec:messages}

The syntax distinguishes two basic ways to send messages:

\begin{longtable}{L{0.28\textwidth}L{0.29\textwidth}L{0.33\textwidth}}
\textbf{Form} & \textbf{Meaning} & \textbf{Typical use}\\\hline
\code{target <== message} & Now-type message: the caller waits for a reply. &
Queries or computations whose result is needed immediately.\\
\code{target <= message} & Past-type message: delivery is asynchronous. &
Events, starting work, and actor-style pipelines.\\
\code{target <= message @ reply} & Past-type message with an explicit
reply destination. & Sending a later reply to a specified object.\\
\end{longtable}

\begin{lstlisting}
; Query and wait for the result:
[calculator <== [:square 9]]

; Start work asynchronously:
[router <= [:route :urgent "done"]]
\end{lstlisting}

A past-type message lets the sending computation continue without waiting
for a reply. This raises practical questions: When is the work complete?
In what order do messages arrive? How is shared state accessed?
The protocol between your objects should answer these questions.

\subsection{Explicit Reply Destinations and Continuations}

A reply destination is itself an object.
A continuation object is convenient for a short reply handler:

\begin{lstlisting}
[counter <= :get @ [cont value
                  (format t "The counter is ~A~%" value)]]
\end{lstlisting}

The form \code{[cont pattern body]} is an efficient shorthand for an object
with one script. Its body should finish quickly and predictably.
Loops or additional now-type messages in continuation bodies may lead
to deadlock.
Its body is compiled as non-reifying code, so \code{reify-args} and
\code{reify-args-to-meta} cannot be used there~\cite{guide}.
These forms expose expressions and their execution context to a program;
Section~\ref{sec:reflection} introduces that reflective access.
For now, use continuations for short, ordinary reply processing.

\subsection{Forwarding a Reply Without Waiting}

The \code{@ R} field in a script header binds the current reply destination.
It is different from an argument in the message payload.
This forwarding object keeps the original caller's destination:

\sourcelisting[firstline=32]{tutorial/examples/01-basics.abclr2}

\begin{lstlisting}
[forwarder <== [:square 8]]     ; => 64
\end{lstlisting}

The caller waits on the forwarder, but the forwarder's handler does not
wait on the calculator. It sends a past-type message with \code{@ R}.
The calculator's reply then goes directly to the original caller.
Do not add an extra \code{!} to this handler: that would send a second reply.
This style is useful when an object coordinates work but does not need
the result itself.

\subsection{Object Identity and the Current Object}

Inside a script, \code{Me} denotes the current object.
A handler can send a message to itself using \code{[Me <= message]},
usually to arrange work after the current handler finishes.
Under the ordinary metaobject, a handler is processed to completion before
another normal handler for that object begins. This remains true when
the current handler is suspended waiting on a now-type reply.

Consequently, a now-type message to the same ordinary object can deadlock:
the current handler waits for a reply from a handler that cannot begin.
The same problem arises in a cycle: object A waits on B while B waits on A.
Use forwarding or a later past-type phase to break that wait cycle.
Ordinary queue serialization protects the object's handlers from overlapping;
it does not make operations on several objects into one transaction.
A \emph{join} waits for the notifications required by an application
protocol, such as both children's answers, before continuing. It does
not mean that every object or host thread has stopped.

\subsection{Concurrency and Parallel Hardware}

Independent objects can be logically concurrent even when the host executes
one computation at a time. The supplied runtime has a scheduling queue
in \filepath{src/process.lsp}; its existence does not imply that an object
is mapped to a dedicated operating-system thread or CPU.
Long computations in a host Lisp function can also delay other work.
Use the word \emph{concurrency} for the protocol model, and measure the
particular runtime before making claims about parallel speedup.


\section{A Complete Example: A Bank Account}
\label{sec:account}

The following object combines state, patterns, replies, and a simple
invariant. It rejects withdrawals that would make the balance negative.

\begin{codeblock}
\begin{lstlisting}
[object account
  (state [balance := 0])
  (script
    (=> [:deposit amount]
        (if (> amount 0)
            (progn [balance := (+ balance amount)] !balance)
            !:invalid-amount))
    (=> [:withdraw amount]
        (if (and (> amount 0) (>= balance amount))
            (progn [balance := (- balance amount)] !balance)
            !:rejected))
    (=> :balance !balance))]

[account <== [:deposit 100]]     ; => 100
[account <== [:withdraw 30]]     ; => 70
[account <== :balance]           ; => 70
\end{lstlisting}
\end{codeblock}

\paragraph{Exercise 2.}
Add a \code{[:transfer other amount]} handler. First check whether the
withdrawal is possible, then send a deposit to \code{other}.
Consider whether to use a now-type or past-type message, what should
happen if the destination rejects the deposit, and when the transfer
can be reported as complete.
Compare your protocol with Section~\ref{sol:transfer}.

\section{Concurrency as a Programming Model}

A producer can send work to a consumer without waiting for an individual
reply to each item:

\begin{codeblock}
\begin{lstlisting}
[object consumer
  (script
    (=> [:item x]
        (format t "Processed: ~A~%" (* x x))))]

[object producer
  (script
    (=> :run
        (progn
          [consumer <= [:item 2]]
          [consumer <= [:item 3]]
          [consumer <= [:item 4]]
          !:sent)))]

[producer <== :run]
\end{lstlisting}
\end{codeblock}

The three past-type messages are sent to \code{consumer}.
The reply \code{:sent} confirms that the producer has sent the requests;
it does not confirm that the consumer has completed them.
Avoid relying on an assumed global execution order, especially with
multiple senders. A clear protocol specifies both permitted messages
and how completion is acknowledged.

\paragraph{Exercise 3.}
Give the consumer a state counter \code{n} that increases after each
\code{:item}. Add a \code{:count} message that returns its current value.
Experiment with querying the count while work is still outstanding.
Then design an acknowledgment that tells the producer when all its
items have been processed.
Compare your completion condition with Section~\ref{sol:count}.

\section{Creating and Organizing Objects}

Objects can be created inside expressions. Start with an ordinary local object:

\begin{lstlisting}
(let ((prefix "[log]"))
  [object local-logger
    (script
      (=> [:write text]
          (format t "~A ~A~%" prefix text)))]
  [local-logger <= [:write "start"]])
\end{lstlisting}

The object name is a binding in the surrounding environment.
The object can use lexical variables from that environment.
Because these variables are shared, take care when mutating them.

A \emph{group} supplies an object-creation and execution policy for its
members. It is not a class or a shared state container.
An object normally uses its creation context's group; an explicit
\code{(group manager)} selects another group manager. Groups and custom
metaobject generators become important in Section~\ref{sec:reflection}.

\begin{lstlisting}
[group logging-group]
[object logger-a (group logging-group) (script (=> :ping !:ok))]
[object logger-b (group logging-group) (script (=> :ping !:ok))]
(eq [group-of logger-a] logging-group)          ; => T
(eq [group-of logger-a] [group-of logger-b])     ; => T
(eq logger-a logger-b)                         ; => NIL
\end{lstlisting}
Here \code{group-of} returns the ordinary object's group manager.
Sharing that manager does not make two objects identical or combine
their state. Custom generators in the original projects must be read
on their own terms; selecting a generator does not certify every
introspection operation it might support.

\section{Files, Libraries, and Tools}
\label{sec:files}
\label{sec:loading}

Loading a file is part of program execution in ABCL/R2. Its expressions
run in an environment, using a group and evaluator; the same translated
file can therefore behave differently when loaded by different callers.
First separate preparing a file from executing it:

\par\medskip\noindent
\begin{minipage}{\linewidth}
\begin{tabular}{L{0.32\textwidth}L{0.58\textwidth}}
\textbf{Expression} & \textbf{Purpose}\\\hline
\code{(inc "name")} & Translate if needed, then load and evaluate in the
caller's environment, group, and evaluator context.\\
\code{(comp "name")} & Translate and compile if needed; does not run the
ordinary ABCL application expressions.\\
\code{(inc-lib "name")} & Load and evaluate in the global environment,
top-level group, and current top-level evaluator.\\
\code{(require 'name)} & Search for and load an unregistered module through
\code{inc-lib}.\\
\code{(provide 'name)} & Record a module name in \code{*modules*}; does not
load a file or validate its definitions.\\
\code{(pwd)} / \code{(cd "dir")} & Historical directory commands;
not implemented for SBCL in this port.\\
\code{:q} / \code{:reset} & Quit the interpreter or reset the system.\\
\end{tabular}
\end{minipage}\par\medskip

Use a file's \emph{base pathname}, without its type:
\code{(inc "demo" :use-source t)} loads \filepath{demo.abclr2}.
Translation writes a neighboring \filepath{demo.lsp}; compilation writes
a host-specific compiled file, normally \filepath{demo.fasl} on SBCL.
Do not pass the \code{.abclr2} type as part of the base name: the loader
derives all three paths from that argument.
The directory must permit writing the generated files. Use scratch
copies when exploring the original examples, as the checks do.

Relative paths are resolved against the process's current directory,
not automatically against the containing source file's directory.
Start the interpreter in the required directory or use absolute paths.
The host-specific implementations in \filepath{src/topfunc.lsp} have
no SBCL branch for changing that directory. Calling \code{inc-lib} does
not make its relative path relative to \filepath{lib/}.

\subsection{Options and File Selection}

The user's guide specifies the public options~\cite{guide}; the active
implementations are in \filepath{src/inc.lsp} and \filepath{src/user.lsp}.
The defaults are:

\noindent\begin{minipage}{\linewidth}
\begin{longtable}{L{0.20\textwidth}L{0.15\textwidth}L{0.55\textwidth}}
\textbf{Option} & \textbf{Default} & \textbf{Accepted by and effect}\\\hline
\code{:use-source} & \code{nil} & \code{inc}, \code{inc-lib}, \code{require},
\code{comp}. Non-nil forces translation from ABCL source. On a load it
also bypasses the compiled-file branch; on \code{comp} it forces compilation.\\
\code{:verbose} & \code{t} & All four commands. Controls translation and
the passed-through load/compile verbosity, not every warning or diagnostic.\\
\code{:compile} & \code{nil} & \code{inc}, \code{inc-lib}, \code{require}.
Non-nil runs \code{comp} before loading. \code{comp} itself has no
\code{:compile} option.\\
\end{longtable}
\end{minipage}

\code{inc} and \code{inc-lib} can prompt for a filename when it is
omitted. \code{comp} requires a name, and \code{require} requires a
module name. In scripts and reproducible sessions, always supply it
explicitly. \code{provide} takes only its module-name argument.

\paragraph{Follow the three-stage choice.}
For an explicit \code{inc} or \code{inc-lib}, read the loader as follows:

\begin{enumerate}
  \item If \code{:compile} is true, first perform the compilation command
        with the selected \code{:use-source} and \code{:verbose} values.
  \item Require the \code{.abclr2} source to exist. Translate it if forced,
        if there is no \code{.lsp}, or if the source is newer than the
        translation. Even a cached compiled file does not remove this
        source-file requirement.
  \item Load the compiled file only when \code{:use-source} is false,
        that file exists, and it is \emph{strictly newer} than the Lisp
        translation. Otherwise load the translation. Execute its
        expressions in the selected loading context.
\end{enumerate}

\code{comp} uses the same translation test, then compiles if forced,
if the compiled file is absent, or if the translation is newer than
the compiled file. Otherwise it reports that the compiled file is up
to date. It does not execute an ordinary ABCL object definition or
application send merely by compiling it. Translation still processes
macros; do not equate compilation with an absence of all host-side effects.

Timestamp comparisons are not content comparisons. Equal timestamps
do not satisfy ``strictly newer'' in the load test, so loading may select
the Lisp file even after compilation. The combination
\code{:compile t :use-source t} compiles first but still follows the
forced-source load branch. Compilation is not a promise that this
particular load executes the FASL. This is usually a performance/file
selection distinction, not a different intended program result.

For day-to-day editing, prefer the explicit source path:

\begin{codeblock}
\begin{lstlisting}
(inc "demo" :use-source t)          ; retranslate, load, and run
(comp "demo" :use-source t)         ; prepare compiled output; do not run
(inc "demo" :compile t)             ; update as needed, then load and run
\end{lstlisting}
\end{codeblock}

These are three independent commands, not a required sequence. Do not
repeat a loading command solely because a translation was cached:
\code{inc} and \code{inc-lib} execute the file's expressions on each call,
including assignments, object creation, and application sends.
\code{:verbose nil} can reduce output while keeping those effects.

\subsection{Execution Context and Module Loading}

\paragraph{An included fragment sees its caller's environment.}
The supplemental \filepath{tutorial/examples/15-loaded-fragment.abclr2}
contains just these two expressions:

\wholelisting{tutorial/examples/15-loaded-fragment.abclr2}

It deliberately depends on \code{load-token}; it is a fragment to load
in a prepared environment, not a self-contained program. Try:

\begin{codeblock}
\begin{lstlisting}
[load-token := :global]
[loaded-value := :unchanged]
[loaded-object := nil]
(let ((load-token :local) (loaded-value nil) (loaded-object nil))
  (inc "tutorial/examples/15-loaded-fragment" :use-source t)
  (list loaded-value [loaded-object <== :value]))
; => (:LOCAL :LOCAL)
loaded-value                       ; => :UNCHANGED
\end{lstlisting}
\end{codeblock}

The fragment updates the local \code{loaded-value}, binds its new object
locally, and that object's script retains the local \code{load-token}.
The same fragment loaded by \code{inc-lib} instead uses the global
environment:

\begin{codeblock}
\begin{lstlisting}
(let ((load-token :local) (loaded-value nil) (loaded-object nil))
  (inc-lib "tutorial/examples/15-loaded-fragment" :use-source t)
  (list loaded-value loaded-object)) ; => (NIL NIL)
loaded-value                        ; => :GLOBAL
[loaded-object <== :value]           ; => :GLOBAL
\end{lstlisting}
\end{codeblock}

Group context matters too. Ordinary objects created by \code{inc} inside
a client's script inherit that caller's group; \code{inc-lib} evaluates
the file in the top-level group using the current \code{*eval*}.
An explicit \code{group} or \code{meta-gen} in the loaded definitions
can still choose different machinery. Loading a library does not itself
replace existing objects' evaluator references or compose its policy
with every custom group. This is why the original examples' setup and
loading order must be read as part of their protocols.

\paragraph{Require a module name, not an arbitrary file path.}
\label{sec:library-path}
\code{*library-directories*} contains the directories searched by the
ABCL/R2 \code{require} macro. It selects the first directory containing
the named \code{.abclr2} source, then uses \code{inc-lib} to load it.
The default value is \code{(list "")}: only the current working directory
is searched. The distribution's \filepath{lib/} directory is not added
automatically. Before running examples that require its modules, add it
at the ABCL/R2 prompt, starting the interpreter at the repository root:

\begin{codeblock}
\begin{lstlisting}
[*library-directories* :=
  (cons #P"lib/" *library-directories*)]
(require 'future)
\end{lstlisting}
\end{codeblock}

Here \code{\#P"lib/"} is Common Lisp's pathname literal for a directory.
Keep the trailing slash: the search uses the directory component of
each entry. Relative entries are resolved against the working directory,
not the executable's location or the directory of the file being loaded.
Use an absolute directory pathname when starting elsewhere.
Explicit \code{inc} and \code{inc-lib} calls do not search this list;
\code{(inc-lib "lib/future")} instead names the file directly.
The path is session configuration; after \code{:reset}, set it again.

Symbols are converted to lowercase names. String arguments retain their
case during search; prefer symbols or explicitly lowercase strings.
\code{provide} normalizes both symbols and strings to lowercase.

For your own library directory, use a directory pathname, save the
previous search path, and restore it after the experiment:

\begin{codeblock}
\begin{lstlisting}
[saved-library-directories := *library-directories*]
[*library-directories* :=
  (cons #P"/absolute/path/to/my-libraries/" *library-directories*)]
; (require 'my-module) ; loads my-module.abclr2 from the search path
[*library-directories* := saved-library-directories]
\end{lstlisting}
\end{codeblock}

The absolute example directory is a placeholder to replace, not a
required installation path.
Putting it first makes its source win over later directories. A typical
module includes \code{(provide 'my-module)} to register its name in
\code{*modules*}. \code{require} does not register a module automatically:
without \code{provide}, a second require can execute the file again.

Once a name is registered, \code{require} skips loading entirely, even
with \code{:use-source t} or \code{:compile t}. To intentionally reload
a known file while editing, use its explicit base pathname with
\code{inc-lib}. Be aware that this reexecutes definitions and effects;
it is not an atomic replacement of every reference to earlier objects.
If \code{provide} ran before a later failure in a module, its recorded
name alone is not evidence that all initialization succeeded.

A missing module also returns \code{nil} without loading in this
implementation. That return value is not a success certificate or a
reliable distinction between ``already registered'' and ``not found.''
Inspect the transcript, expected definitions, and search path. Direct
\code{inc} and \code{inc-lib} instead signal an error for missing source.

Included libraries are \code{future}, \code{multicast}, \code{express},
\code{heavy-meta}, \code{copy}, \code{watch}, and \code{report}.
Their walkthroughs explain setup and use. Historical alternatives with
shared feature or generator names require separate fresh sessions, not
an assumption that \code{require} can distinguish filenames by itself.

\paragraph{Evidence and check your understanding.}
Run \code{make -C tutorial loading-check}. Thirty-five probes check
option defaults, compile-versus-load effects, independently marked Lisp
and FASL cache selection, forced source, missing-source rejection,
module registration/reload/search, and local versus global environment,
group, and evaluator behavior. Tests modify only generated cache files
in \filepath{build/loading-tests/}; original sources remain unchanged.
The small fragment above is copied there before testing.

Why might changing a module and calling \code{require} show no change?
\emph{Answer:} The registered name bypasses the load, including its
source option. Why can \code{comp} create output without creating the
application object? \emph{Answer:} Preparation does not execute its
ordinary ABCL expressions. Why does a file see a local value with
\code{inc} but the global value with \code{inc-lib}?
\emph{Answer:} Loading selects an execution environment, not just a file.


\section{Futures and Express Messages}
\label{sec:message-extensions}

Future-type messages are a library feature. Load the library before
using them:

\begin{lstlisting}
(require 'future)
\end{lstlisting}

Futures allow a result to be requested and consumed later, separating
the request from the point at which the result is needed.
See \filepath{lib/future.abclr2} and
\filepath{examples/misc/future-ex.abclr2} for the specific syntax and examples.

\subsection{A Future Is a Reply Stream}
\label{sec:futures}

The supplied future library can collect more than one reply from a request.
The send syntax uses a dollar sign:
\code{[receiver <= message \$ future-object]}.
Each \code{!} in the receiver's handler appends another reply to the future.
This is more general than a promise that resolves to one value.

To run the following example from the repository root:

\begin{lstlisting}
(inc "tutorial/examples/04-futures" :use-source t)
[range-client <== [:collect 1 4]]
; => (1 2 3 4)
\end{lstlisting}

The file loads the supplied future library. Here is the server:

\wholelisting[firstline=4,lastline=11]{tutorial/examples/04-futures.abclr2}

The server sends each integer and then \code{:done}.
The terminal message is part of our protocol: futures themselves do not
know that a producer has finished.

\wholelisting[firstline=12]{tutorial/examples/04-futures.abclr2}

The initial send is asynchronous. Each \code{next-value} waits if no reply
is currently available. Thus waiting occurs at the point where the client
needs a particular element, rather than at the initial request.
The sentinel is removed from the stream but excluded from the result.
The range from 3 to 2 is empty and returns \code{nil}; the server still
sends \code{:done}, so the client does not wait forever.

\subsection{Peeking, Consuming, and Readiness}

\begin{longtable}{L{0.39\textwidth}L{0.51\textwidth}}
\textbf{Operation} & \textbf{Behavior in the supplied implementation}\\\hline
\endhead
\code{(make-future)} & Create and initialize an empty future.\\
\code{(ready? f)} & Return whether at least one reply is buffered.\\
\code{(next-value f)} & Wait for and remove the next reply.\\
\code{(next-value f :remove nil)} & Wait for the next reply without removing it.\\
\code{(all-values f :wait nil)} & Return the current buffered replies,
or \code{nil} if none exist; removal defaults to true.\\
\code{(reset-future f)} & Send a reset request; use only when old producers
and readers have finished.\\
\end{longtable}

The actual macros in \filepath{src/user.lsp} default \code{:remove}
to true for \code{next-value} and \code{all-values}; an old comment
in the library header gives a different default.
\code{all-values} is a snapshot of buffered values, not an operation
that waits for a terminal marker. Its \code{:wait t} default waits
when the buffer is empty, but does not establish producer completion.

For stream consumers, use \code{next-value} and an explicit terminal
marker. Use \code{all-values} with \code{:wait nil} to inspect the replies
currently buffered; a snapshot is not a producer-completion signal.

\begin{codeblock}
\begin{lstlisting}
(let ((f (make-future)))
  [range-server <= [:range 8 8] $ f]
  (list (next-value f :remove nil)
        (next-value f)
        (next-value f)
        (ready? f)))
; => (8 8 :DONE NIL)
\end{lstlisting}
\end{codeblock}

The implementation records one waiting consumer destination.
Treat each future as a stream with one active consumer; use a separate
future per request rather than sharing one among unrelated readers.

\paragraph{Exercise 4.}
Change the server to send squares, keeping the terminal marker.
Then change the payload to \code{[:value n]} so that a legitimate data
value cannot be confused with \code{:done}.
Compare the two message shapes with Section~\ref{sol:stream}.


\subsection{Express Messages}

Express messages can preempt a script invoked by a normal message.
They require the sender to use a suitable evaluator and the receiver
to use a suitable metaobject generator. The \code{express} library
provides these components. Once configured, the sending syntax is:

\begin{lstlisting}
[receiver <= message :express]
[receiver <== message :express]
\end{lstlisting}

The documented purpose of \code{atomic} is to prevent express-message
preemption while its subexpressions execute. The first experiment below
studies routing to an idle receiver. Express messages are an advanced facility:
first become comfortable with ordinary past-type and now-type messages.

\subsection{An Express Exchange}

An express send needs both ends configured: the sender uses an express
evaluator, and the receiver uses an express metaobject generator.
Begin with the original clock while it is idle. In a fresh session at
the repository root, load its library and definitions:

\begin{codeblock}
\begin{lstlisting}
(inc-lib "lib/express" :use-source t)
(inc "examples/misc/express-ex" :use-source t)
\end{lstlisting}
\end{codeblock}

The example defines the group \code{express} and its clock.
Place the client in that group too, so its evaluator understands the
extra \code{:express} send argument:

\begin{codeblock}
\begin{lstlisting}
[object finite-express-sender
 (group express)
 (script
  (=> :peek ![clock <== :peek :express])
  (=> :past-peek @ R [clock <= :peek @ R :express])
  (=> :atomic-result !(atomic 42)))]
\end{lstlisting}
\end{codeblock}

Try these short exchanges:

\begin{codeblock}
\begin{lstlisting}
[clock <== :peek]                         ; => 0
[finite-express-sender <== :peek]          ; => 0
[finite-express-sender <== :past-peek]     ; => 0
[finite-express-sender <== :atomic-result] ; => 42
\end{lstlisting}
\end{codeblock}

The first query is ordinary. The second goes through an express-capable
sender. The third forwards the caller's reply destination through a
past-type express send; \code{@ R}, not \code{<=} alone, supplies that
reply path. The last expression shows that \code{atomic} returns its
body's value.

This exercise studies routing to an idle receiver, not interruption of
a running computation. Section~\ref{sec:original-express} follows the
original clock's scripts, mode transitions, and evaluator front end.


\section{Reflection and Hybrid Group Architecture}
\label{sec:reflection}

\subsection{What Reflection Adds}

So far, our messages have been about application data: a balance, a job,
or a value in a stream. \emph{Reflection} lets a program inspect and
operate on representations of its own structure and execution.
It is not merely another domain object or another name for inheritance.
For example, an account's application protocol deposits money; a
reflective protocol can inspect the environment holding its balance,
the scripts available to it, or the evaluator executing those scripts.

Call ordinary application activity the \emph{base level}.
Activity that describes or controls it is at a \emph{meta-level}.
These are roles in a relationship, not dedicated CPU threads.
A \emph{metaobject} represents an individual object's execution machinery.
An \emph{environment} represents bindings and lexical lookup.
An \emph{evaluator} executes expression representations in an environment
and context, reporting the result to a continuation.
A \emph{scheduling policy} decides which pending work to run next;
evaluation describes how an expression is executed. The two can interact:
an evaluator object can queue requests and apply its own scheduling policy.
Do not confuse that object with the host runtime's scheduling machinery.
A \emph{metaobject generator} is a factory for metaobjects.
A \emph{group manager} represents a group and participates in creation
using that group's generator and evaluator.
The evaluator, metaobject generator, and group manager are the group's
\emph{kernel objects}.

Two distinctions will keep the examples readable.
First, changing the machinery that executes a script is different from
changing the application's message payload.
Second, replying to the application caller is different from reporting
that an expression has finished executing.
The original inheritance reading project makes both visible.

\emph{Reification} exposes an execution representation for reflective
computation. ABCL/R2 avoids eagerly constructing an
infinite stack of such representations: higher levels are created lazily.
Its \emph{partial compilation} compiles ordinary computations while
keeping selected operations available through reflective protocols.
The compiled representations passed to an evaluator are therefore not
simply source strings to be re-read.
A \emph{non-reifying} object's script is compiled for direct execution
rather than this reflective expression interface; it cannot use
\code{reify-args} or \code{reify-args-to-meta}. This does not prohibit
ordinary messages to other objects' reflective interfaces: the breakpoint
inspector later uses environment objects in exactly that way.
These implementation techniques are explained in Masuhara's thesis and
the distribution's compiler documentation~\cite{masuhara,guide}.

An object definition can therefore contain additional parameters
alongside \code{state} and \code{script}:

\begin{codeblock}
\begin{lstlisting}
; Schematic: these kernel objects must be defined first.
[object named-object
  (meta-gen my-meta-generator)
  (group my-group-manager)
  (priority 10)
  (script
    (=> :hello !:ok))]
\end{lstlisting}
\end{codeblock}

The interpretation of additional parameters depends on the chosen
runtime objects; \code{priority} has no universal scheduling meaning.
To explore a custom evaluator, study the files
\filepath{express-ex.abclr2}, \filepath{built-ins.abclr2}, and
\filepath{reify-args.abclr2} in \filepath{examples/misc/}.

The special forms \code{reify-args} and \code{reify-args-to-meta}
send compiled argument expressions, together with their evaluation
context, to an evaluator or to the current object's metaobject,
respectively. They are unavailable in scripts of non-reifying objects.

This makes debugging, scheduling, resource management, and alternative
evaluation strategies expressible as programs in the language.
Reflection adds an abstraction level, so encapsulate it in small,
well-understood libraries.

\subsection{Two Reflective Towers, Two Kinds of Control}

The Hybrid Group Architecture distinguishes reflection on an individual
object from reflection on the group in which objects execute.
Masuhara's thesis, Section 3.3, develops this distinction~\cite{masuhara}.

\begin{longtable}{L{0.26\textwidth}L{0.32\textwidth}L{0.32\textwidth}}
\textbf{Level} & \textbf{What it controls} & \textbf{Example}\\\hline
Individual metaobject & The object's state environment, scripts, and
normal message queue. & Inspect a counter's state or change its evaluator.\\
Group kernel objects & Shared evaluation and object creation policy. &
Give a group a custom evaluator for a new control form.\\
Higher meta-levels & The corresponding behavior of metaobjects and
meta-groups. & Reflect on the evaluator itself.\\
\end{longtable}

This conceptual hierarchy is not allocated as an infinite collection
of objects. Chapter 4 of the thesis explains lazy creation and lightweight
metaobjects: forwarding an ordinary message should not force expensive
reification. Section 4.7 discusses non-reifying objects; Chapter 7
summarizes the use of partial compilation~\cite{masuhara}.
Sequential host computations can be compiled while selected operations
still pass through the reflective protocols.

\subsection{Inspecting an Object's State}

Load the reflection example in a fresh session:

\begin{lstlisting}
(inc "tutorial/examples/05-reflection" :use-source t)
[observed-counter <== :inc]                  ; => 1
[[[meta observed-counter] <== :state]
 <== [:value-of 'n]]                         ; => 1
[[[meta observed-counter] <== :state]
 <== [:set 'n 40]]                           ; => 40
[observed-counter <== :inc]                  ; => 41
\end{lstlisting}

The steps are significant. \code{[meta observed-counter]} obtains the
counter's metaobject. The \code{:state} query returns an environment
object, not a plain property map. Sending \code{[:value-of 'n]} to that
environment retrieves a binding; \code{[:set 'n 40]} changes it.
The quoted symbol is the variable's name as data.
The next ordinary increment observes that change: the meta-level
representation is connected to the computation it describes.

This bypasses the object's application protocol. Use it to understand
reflection and debugging; ordinary clients should continue to use
the counter's defined messages. Default metaobjects also answer
\code{:scripts}, \code{:queue}, and \code{:evaluator}.
Their precise replies are runtime objects or internal queue entries.
The operation \code{[den metaobject]} returns the associated lower-level
object, called its \emph{denotation}. It is not a copy or a variable's
value. In the original inheritance project we will check both identity
directions and distinguish this relationship from group membership.

\subsection{Defining a Lazy Control Form}

An ordinary function evaluates its arguments before receiving them.
A control form must instead decide which expression to evaluate.
The next example implements a choice using the group's evaluator:

\wholelisting[linerange={9-18,25-26}]{tutorial/examples/05-reflection.abclr2}

The new evaluator handles the \code{:choose} protocol and forwards all
other messages to a primary evaluator. The internal helper
\code{r2sys::create-primary-eval} matches the helper used by the
repository's examples; it is specific to this implementation.
Inside a group definition, \code{group} names the group being constructed.

The expression variables \code{predicate}, \code{then}, and \code{else}
contain compiled expression representations. They are not ordinary
argument results. The \code{[:do expression Env Id Gid Eval]} protocol
asks the evaluator to execute one representation in its original context.
Read the context fields by role, not as application arguments:

\noindent\begin{minipage}{\linewidth}
\begin{longtable}{L{0.13\textwidth}L{0.77\textwidth}}
\textbf{Field} & \textbf{Meaning in this evaluation request}\\\hline
\code{Env} & Environment used for variable lookup.\\
\code{Id} & Metaobject identity of the object whose expression is being
executed; evaluating \code{Me} produces \code{[den Id]}.\\
\code{Gid} & Group manager for that execution context; evaluating
\code{group} produces this value.\\
\code{Eval} & Evaluator used for the expression's evaluation operations.\\
\code{C} & Continuation receiving this expression's result, not
necessarily the application's original caller.\\
\end{longtable}
\end{minipage}

When a metaobject's own handler executes, \code{Me} is that metaobject,
not the application object it manages. Likewise, \code{Gid} is not
automatically the group of \code{Id}: a metaobject belongs to a group
at its own level. Keep object, metaobject, and group relationships separate.

The first send evaluates the predicate. Its continuation selects
one branch representation and forwards the final destination \code{C}.
No unselected branch is evaluated.

\sourcelisting[linerange={28-36,41-41}]{tutorial/examples/05-reflection.abclr2}

\begin{codeblock}
\begin{lstlisting}
[choice-test <== [:choose 3]]   ; => 2
[choice-test <== :selected]     ; => :ODD
[choice-test <== [:choose 8]]   ; => 4
[choice-test <== :selected]     ; => :EVEN
\end{lstlisting}
\end{codeblock}

The current evaluator receives compiled argument expressions through
\code{reify-args}. The corresponding form for the current object's
metaobject is \code{reify-args-to-meta}. Both preserve the evaluation context,
and neither is available in a non-reifying object's script.
This example uses a reflective control form to add behavior; it does
not change the meaning of the ordinary \code{if} operator.

The tag supplied to \code{reify-args} is evaluated first. The remaining
argument expressions are passed as compiled representations in a message
of the form \code{[Tag Expressions Env Id Gid Eval] @ C}; they have
not already been reduced to argument values~\cite{guide}.

\paragraph{Exercise 5.}
Extend the choice group with a \code{:unless} protocol that evaluates
one body only when its predicate is false. Forward the final result
to the original continuation in both cases.
Check both branches against Section~\ref{sol:unless}.

\subsection{Extension Parameters and Host Events}

The user's guide identifies additional parameters on message transmission,
object creation, and assignment as extension points~\cite{guide}.
They carry information to the execution machinery; they are not extra
elements of the application's message payload. For example, in the
schematic send \code{[target <= payload @ reply :RVT stamp]}, the
evaluator representation is
\code{[:send-past Target Payload Reply :RVT Stamp]}.
Capitalized names here denote evaluated values.
It is the selected evaluator's responsibility to interpret this extension.
An ordinary receiver still receives \code{Payload}.

Likewise, \code{(priority n)} in an object definition becomes a
\code{:priority} argument in its creation representation, rather than
automatically creating a state variable or a universal scheduling rule.
The priority-control examples supply machinery that gives this parameter
meaning. An extended assignment \code{[x := value tag argument]}
passes the evaluated extras to the environment's \code{:set} request.
Custom environment objects can therefore participate in update policy.
Do not assume the default environment acts on an arbitrary tag.

Two library facilities illustrate this principle without extending the
core syntax. Script-pattern \code{where} constraints require
\code{heavy-meta}; constraints inside an ordinary \code{match} form
remain unsupported. Multicast requires \code{multicast} and a suitable
evaluator created by \code{multicast-eval-gen}; merely loading the
library does not replace every object's evaluator.
The original \filepath{lib/multicast.abclr2} constructs a template for
nested target lists and collects replies into their corresponding
positions. This is a library-defined protocol, not the unsupported
brace-delimited synchronous-send form.
Section~\ref{sec:multicast} works through its original example and reply paths.
These libraries need their own setup and reply/completion analysis,
just as the express library does.

The guide also describes a boundary for host events~\cite{guide}.
At Lisp level, \code{r2sys:send-signal} delivers its arguments as a
\code{[:signal ...]} message to the object bound to
\code{*signal-handler*}. The default handler ignores events; a
\code{nil} handler disables delivery.
Alternatively, the host can set \code{r2sys:*signal-request*};
the runtime checks and clears it before forwarding the request.
This is a single pending-request slot, not a general event queue.
The ABCL side handles an event as a message. It does not follow that
every operating-system signal or timer is automatically connected to
this interface: that requires host-specific integration.
The interface also does not promise express preemption.

\subsection{A Safe Experiment with Synthetic Host Events}
\label{sec:host-events}

The event interface turns a host notification into an ABCL message.
You can explore that boundary without a timer, an operating-system signal,
or another thread. The following supplemental exercise uses the interface
described by the user's guide and implemented in \filepath{src/signal.lsp};
there is no corresponding runnable event example in \filepath{examples/}.
It changes no original source and installs no operating-system callbacks.

\paragraph{Receive and inspect an event.}
The file \filepath{tutorial/examples/14-host-events.abclr2} defines an
ordinary recording object. Loading it does not install it as a handler:

\wholelisting[firstline=2,lastline=8]{tutorial/examples/14-host-events.abclr2}

The dotted pattern \code{[:signal . parameters]} binds the argument tail
as a list. Each saved record contains that list, the reply destination,
and the sender. A host notification arrives with both reply and sender
\code{nil}: it is not an application call waiting for \code{!}.
The separate \code{:events} query does reply, allowing us to inspect
what the handler has processed.

Start a fresh interpreter at the repository root and save its current
handler before temporarily replacing it:

\begin{codeblock}
\begin{lstlisting}
(inc "tutorial/examples/14-host-events" :use-source t)
[saved-event-handler := *signal-handler*]
[*signal-handler* := event-inbox]
(r2sys:send-signal :tick 7)
[event-inbox <== :events]           ; => (((:TICK 7) NIL NIL))
\end{lstlisting}
\end{codeblock}

The \code{r2sys:} prefix names the exported host interface in the
ABCL/R2 system package. \code{send-signal} accepts an ordinary Lisp
\code{\&rest} argument list and constructs \code{(:signal :tick 7)}.
It submits the notification; returning from that function is not an
acknowledgment from your handler. This finite interactive exercise lets
the runtime dispatch the event before inspecting the inbox. Do not infer
a general ordering guarantee between independent producers, or use a
notification as proof that unrelated work finished.

Passing one list is different from passing its elements as arguments:

\begin{codeblock}
\begin{lstlisting}
[event-inbox <== :clear]
(r2sys:send-signal '(:tick 7))
[event-inbox <== :events]           ; => ((((:TICK 7)) NIL NIL))
\end{lstlisting}
\end{codeblock}

The parameter list now has one element, itself a list. Similarly,
\code{(r2sys:send-signal)} sends an event with an empty parameter list,
while \code{(r2sys:send-signal nil)} sends one with parameter list
\code{(nil)}. Neither call means ``disable events.''

\paragraph{Explore the single pending-request slot.}
The other interface is a \emph{host Lisp variable},
\code{r2sys:*signal-request*}. It is distinct from a variable stored in
an ABCL environment. An ABCL assignment such as
\code{[r2sys:*signal-request* := value]} would not be the host-side
write this interface requires. The teaching file supplies small host
functions to make that boundary explicit:

\wholelisting[firstline=11,lastline=17]{tutorial/examples/14-host-events.abclr2}

These function bodies run in host Lisp. Their \code{setf} writes the
actual host variable; \code{demo-event-slot} reads it. They are teaching
helpers, not additional built-in ABCL/R2 operations.
In the same fresh session, first check that the slot is empty:

\begin{codeblock}
\begin{lstlisting}
(demo-event-slot)                  ; => NIL
[event-inbox <== :clear]
(demo-post-event '(:tick 7))
[event-inbox <== :events]           ; => ((((:TICK 7)) NIL NIL))
(demo-event-slot)                  ; => NIL: consumed by the runtime
\end{lstlisting}
\end{codeblock}

A non-nil slot value becomes \emph{one} argument following \code{:signal}.
A list stored there is not spliced into multiple arguments.
The original polling function makes this explicit:

\wholelisting[firstline=57,lastline=61]{src/signal.lsp}

It saves the value, clears the slot, and forwards \code{(:signal value)}.
The scheduler calls this function at the top of its dispatch loop in
\filepath{src/process.lsp}. The slot is not a background service that
keeps an otherwise idle interpreter polling, and its check does not
interrupt a host Lisp computation halfway through execution.

Two writes before that polling point do not create two pending events:

\begin{codeblock}
\begin{lstlisting}
[event-inbox <== :clear]
(demo-replace-event :first :second)
[event-inbox <== :events]           ; => (((:SECOND) NIL NIL))
\end{lstlisting}
\end{codeblock}

Both writes happen inside one host function, with no ABCL scheduler
step between them. Only the last value remains to be polled. This
demonstrates overwrite, not queueing or reliable event aggregation.
Writing \code{nil} means no pending request; use
\code{(r2sys:send-signal nil)} if you actually need an event whose
argument is \code{nil}. These helpers assume an isolated fresh session;
do not overwrite a pending request owned by another event producer.

\paragraph{Disable delivery, then restore the saved handler.}
The handler binding controls both interfaces:

\begin{codeblock}
\begin{lstlisting}
[event-inbox <== :clear]
[*signal-handler* := nil]
(r2sys:send-signal :dropped)
(demo-post-event :also-dropped)
(demo-event-slot)                  ; => NIL after polling
[*signal-handler* := event-inbox]
[event-inbox <== :events]           ; => NIL: no replay
[*signal-handler* := saved-event-handler]
(eq *signal-handler* saved-event-handler) ; => T
\end{lstlisting}
\end{codeblock}

Disabling delivery does not buffer notifications for later. A polled
slot value is cleared even when no handler is installed. Conversely,
changing the handler binding is not a cancellation of messages already
sent to a previous handler. Always restore the object you saved, rather
than assuming that the default handler is still the desired one.
The normal default handler prints an ignored-signal diagnostic; ``ignored''
does not necessarily mean silent.

No step in this exercise enables keyboard interrupts or timer delivery.
The historical keyboard-interrupt wrapper in \filepath{src/signal.lsp}
contains host-specific branches; its fallback simply executes its body.
Do not infer an SBCL operating-system hook from the presence of this API.
Nor is this an Express-message experiment: the original Interrupt and
Express examples have their own reflective execution protocols.

\paragraph{Evidence and exercise.}
Run \code{make -C tutorial host-events-check}. Twenty-five probes in two
fresh processes check handler preservation, direct argument shapes,
nil sender/reply fields, slot consumption and overwrite, disabled delivery,
absence of replay, and restoration. Inputs, copied teaching source and
transcripts remain in \filepath{build/host-event-tests/}. Tests install
no operating-system callbacks and modify no runtime or original sources.

Why does a list-valued slot notification differ from
\code{(r2sys:send-signal :tick 7)}? \emph{Answer:} The slot carries one
value; the function accepts several arguments. Why is a zero-valued
slot request still delivered? \emph{Answer:} Only \code{nil} is false
in Lisp; zero is non-nil. Why should the signal handler not rely on
\code{!} to acknowledge the host? \emph{Answer:} The notification has
no application reply destination. Define a separate acknowledgment
protocol if your application needs one.



\section{Troubleshooting and Good Practice}

\begin{enumerate}
  \item Test handlers that return results with now-type messages first.
        Check that their patterns match and their replies are correct.
  \item Then introduce past-type messages and specify when replies
        and completion acknowledgments are expected.
  \item Keep shared state in a responsible object rather than spreading
        mutable lexical variables across many objects.
  \item Use clear message tags, such as \code{:deposit},
        \code{:withdraw}, and \code{:shutdown}.
  \item After an error, this distribution may display a Lisp error message
        and attempt to continue execution. Use \code{:reset} when the
        session's state is unclear.
  \item Build programs incrementally: objects, patterns, state,
        asynchronous communication, and finally reflection.
\end{enumerate}

\section{Working from the Original Examples}
\label{sec:originals}

Read each project from the application's request to its result before
following the machinery underneath. The tree walker uses the futures
you already know. The next projects connect the reflection introduction
to lazy control forms, inheritance, and interpreters; only then do we
combine those mechanisms in libraries, debugging, and scheduling.
Pcontrol and Time Warp are the longest projects and can be studied in
separate sessions. Appendix~\ref{sec:source-map} is a return map to their
files, not an alternative order in which to learn the concepts.

These projects follow the distributed programs in \filepath{examples/}.
The calculator, account service, and job service are teaching extensions;
use them to prepare for reading the original applications.
Line references use code-only views, omitting source-maintenance comments.

\paragraph{Read each result in its own protocol.}
The reflective examples contain several overlapping conversations.
Use this key when following a destination or an apparent completion:

\noindent\begin{minipage}{\linewidth}
\begin{longtable}{L{0.27\textwidth}L{0.63\textwidth}}
\textbf{Conversation} & \textbf{What its response establishes}\\\hline
Application request & A reply answers the application's caller.
Its destination is often named \code{R} or \code{Reply}.\\
Expression evaluation & A result sent to an evaluation continuation,
often \code{C}, finishes that expression, not necessarily the work
it has requested elsewhere.\\
Object dispatch & A \code{:end} message in the ordinary dispatch cycle
lets the metaobject advance its queue; it is not a reply to the client.\\
Object creation & A generator's reply supplies a new metaobject to
the creation machinery; it is not an application-method result.\\
\end{longtable}
\end{minipage}

The names are conventions in the examples, not fixed roles for every
variable with that spelling. A now-type send can use its evaluation
continuation as the application's reply destination, deliberately joining
two of these paths. Follow the actual sends before deciding that two
destinations are different, or that one result means all work is complete.

\subsection{Tree Leaves: The Original Future Protocol}
\label{sec:original-future}

\filepath{examples/misc/future-ex.abclr2} supplies a useful
starting point: a producer walks a nested tree and sends its leaves
through a future. Here is its original producer, without adaptation:

\wholelisting[firstline=17,lastline=30]{examples/misc/future-ex.abclr2}

The first two patterns reshape nested list structure; the third emits a
leaf and advances. Notice that \code{!A} occurs more than once for one
request. The final \code{!:finish} terminates the application protocol,
not the future mechanism itself.
The original consumer makes this explicit:

\wholelisting[firstline=36,lastline=46]{examples/misc/future-ex.abclr2}

In a fresh session using the rebuilt port at the repository root,
load the unchanged original source:

\begin{codeblock}
\begin{lstlisting}
(inc "examples/misc/future-ex" :use-source t)
[print-leaves <= '((a b) (c (d . e) f) g)]
; Prints A, B, C, D, E, F, G on separate lines.
\end{lstlisting}
\end{codeblock}

Loading interactively writes generated translations next to the sources.
For an isolated run that writes only under \filepath{build/}, use the
example checks in Appendix~\ref{sec:book-build}.
This book's range producer deliberately removes the tree-walking
detail while retaining this reply-stream and sentinel protocol.
Once you understand it, return to the original and trace each pattern
against \code{((a b) (c (d . e) f) g)} by hand.

\subsubsection{Read the Original in Three Passes}

First identify the participants. \code{print-leaves} coordinates one request,
\code{tree2leaves} produces leaf replies, and \code{printer} performs output.
A future object buffers the replies between producer and consumer.
The file's first expression loads that library; the three object
definitions then establish the application's protocol.

Second read the consumer. Its \code{temporary} bindings belong to one
handler invocation: \code{leaf} holds the current value and \code{leaves}
is a fresh future. The dollar send starts the walker without waiting for
the entire traversal. Each \code{next-value} removes one reply or waits
until one arrives. The consumer forwards ordinary leaf values to the
printer and stops when it sees the keyword \code{:finish}.

Third read the producer's loop. Its \code{tree} variable is a working
representation of the remaining traversal, not an unchanged input tree.
The first two \code{match} clauses rewrite the remaining work.
The third sends one reply and removes the emitted head.
Only one matching clause executes on an iteration.

\begin{longtable}{L{0.27\textwidth}L{0.63\textwidth}}
\textbf{Pattern} & \textbf{How to read it}\\\hline
\endhead
\code{[[A] . B]} & The head is a one-element proper list.
Remove that extra layer, keeping the outer tail \code{B}.\\
\code{[[A . B] . C]} & The head is a nonempty pair.
Make its head and tail the next two pending items, followed by \code{C}.\\
\code{[A . B]} & The outer head is the next leaf.
Send \code{A} and continue with the outer tail \code{B}.\\
\end{longtable}

These patterns overlap. A one-element list is also a pair, so the first
clause must precede the second. A nested pair is also an outer head, so
the structural clauses precede the emission clause.
Uppercase names in the source are binding variables, not literal
symbols that the tree must contain.

\subsubsection{Trace a Small Tree Completely}

Use the data \code{'((a b) (c . d))}; it exercises both a proper nested
list and a dotted pair. The following is a logical trace of the walker,
not a claim about how its messages interleave with printing.
Unquoted uppercase letters in the table denote the already quoted
input symbols, not new variable lookups.

\begin{longtable}{L{0.30\textwidth}L{0.13\textwidth}L{0.32\textwidth}L{0.12\textwidth}}
\textbf{Before iteration} & \textbf{Clause} & \textbf{After iteration} &
\textbf{Reply}\\\hline
\endhead
\code{((A B) (C . D))} & 2 & \code{(A (B) (C . D))} & none\\
\code{(A (B) (C . D))} & 3 & \code{((B) (C . D))} & \code{A}\\
\code{((B) (C . D))} & 1 & \code{(B (C . D))} & none\\
\code{(B (C . D))} & 3 & \code{((C . D))} & \code{B}\\
\code{((C . D))} & 2 & \code{(C D)} & none\\
\code{(C D)} & 3 & \code{(D)} & \code{C}\\
\code{(D)} & 3 & \code{nil} & \code{D}\\
\end{longtable}

On the first iteration, pattern variables are \code{A = 'a},
\code{B = '(b)}, and \code{C = '((c . d))}.
The construction \code{[A B . C]} puts the value of \code{A} first,
the value of \code{B} second, and uses \code{C} as the final tail.
It therefore preserves the remaining work while exposing the next leaf.
When the working list becomes \code{nil}, the loop exits and the handler
sends one more reply: \code{:finish}.

The original sample's longer input adds further rewrites but follows
the same rules. It yields \code{A B C D E F G}.
No recursive call to \code{tree2leaves} is necessary: the loop and
rewrites implement the traversal.

\subsubsection{Where Does Completion Travel?}

Distinguish three message paths:
\begin{enumerate}[nosep]
  \item The top-level request goes to \code{print-leaves}. Its past-type
        send does not wait for a return value.
  \item Leaf values and the terminal keyword go from \code{tree2leaves}
        to the future's reply destination, then to the waiting consumer.
  \item The consumer sends each leaf to \code{printer} without requesting
        a reply. Printing is a side effect, not the future stream.
\end{enumerate}
The terminal marker tells the consumer that no more leaf values are part
of this request. It does not acknowledge completion of every printer
handler, and the consumer does not reply to its original sender.
Changing the outer invocation to a now-type send would therefore require
adding an application completion protocol; merely changing the arrow
would not create a reply.

\paragraph{Understanding check.}
Why is a leaf equal to the keyword \code{:finish} problematic?
The consumer would treat that data value as the terminal marker.
The original example assumes its data does not contain that marker.
Tagged values such as \code{[:value x]} solve the ambiguity, as in the
later stream exercise. Why not use \code{all-values} as a completion test?
Its buffer snapshot does not say whether the producer has finished.

\paragraph{A boundary of the original algorithm.}
The loop expects an outer proper list of pending items.
An empty nested list can be emitted as a \code{nil} leaf by the final
pattern; it is not automatically skipped as an empty subtree.
Read the actual patterns rather than assuming the behavior of a
general-purpose tree-flattening function.


\subsection{Lazy Evaluation: The Original IF with EF}
\label{sec:original-if}

The first example in \filepath{examples/misc/reify-args.abclr2} contrasts
an ordinary function with a reflective control form.
The function \code{ef} selects between already evaluated arguments:

\sourcelisting[firstline=21,lastline=22]{examples/misc/reify-args.abclr2}

The custom evaluator instead receives representations of expressions.
It evaluates the predicate first and then uses \code{ef} to select which
branch representation to execute:

\wholelisting[firstline=24,lastline=39]{examples/misc/reify-args.abclr2}

The original sample session calls \code{[:ef 3]} and \code{[:my-if 3]}.
Both return 2, but the ordinary function prints both \emph{Even} and \emph{Odd};
the reflective form prints only \emph{Odd}. Equal return values alone would
miss the evaluation-order distinction.
This book's \code{choice-group} is a teaching extension of this
original protocol, using a state marker to make branch selection testable.

Study the two experiments separately. The first occupies code-only lines
1--51; Section~\ref{sec:interrupt} follows the interrupt experiment.
The example checks prepare separate executable excerpts for these projects.
After running them, load the first experiment and compare the calls:

\begin{codeblock}
\begin{lstlisting}
(inc "build/original-tests/reify-first" :use-source t)
[my-if-test <== [:ef 3]]     ; prints Even and Odd; returns 2
[my-if-test <== [:my-if 3]]  ; prints Odd; returns 2
\end{lstlisting}
\end{codeblock}

\subsection{Interrupting a Reflective Control Form}
\label{sec:interrupt}

The second experiment in \filepath{examples/misc/reify-args.abclr2}
defines \code{with-interrupt}. Unlike the preceding \code{my-if}, it
customizes an individual object's metaobject rather than the group's
evaluator. It illustrates \emph{abandoning} an evaluation and selecting
an alternative result. It is not an operating-system interrupt handler,
an exception mechanism with stack unwinding, or state rollback.

\subsubsection{Read the Application Before the Machinery}

The original worker counts, then returns one of two strings:

\wholelisting[firstline=194,lastline=202]{examples/misc/reify-args.abclr2}

The three arguments are the body, the success expression, and the
failure expression. They are passed as compiled representations, not
evaluated argument values. The body prints \code{Count 0.} through
\code{Count n-1.} if allowed to finish. Its ordinary \code{dotimes}
result is not the application's answer: the selected branch supplies
the value of \code{with-interrupt}, and the outer \code{!} sends that
value to the waiting caller.

The original controller starts the worker and later interrupts it:

\wholelisting[firstline=204,lastline=212]{examples/misc/reify-args.abclr2}

The worker's request is asynchronous with \code{@ R}, so the worker
answers the original caller directly. The controller does not wait for
the worker before running its dummy loop. After \code{m} empty iterations,
it addresses \code{[meta test-with-interrupt]} directly. Sending an
ordinary interrupt request to the application would queue it behind
the running script; the direct meta-level request can change evaluation
policy while that script is active. It has no application reply of its
own, and the controller does not execute a second \code{!}.

The commented historical session demonstrates these two outcomes:

\begin{codeblock}
\begin{lstlisting}
[test-with-interrupt-2 <== [:run 10 3]]
; Count 0., Count 1., Count 2.
; Being interrupted... / Context is caught.
; => "Time is up."

[test-with-interrupt-2 <== [:run 10 30]]
; Count 0. through Count 9.
; Interruption is ignored.
; => "Yeah! I did it!"
\end{lstlisting}
\end{codeblock}

The checked excerpt reproduces those counts and branch results on the
current image. The parameter \code{m} is a workload measured in empty
loop iterations, not seconds or a timeout deadline. Do not generalize
the exact point of interruption to another evaluator or scheduler.
Even diagnostic output order needs care: in the current transcript the
late interruption's ``ignored'' message appears after the check has
already received the success result. The worker's reply does not join
every outstanding action in the controller.

\subsubsection{How the Control Form Reaches the Metaobject}

The group's \code{meta-gen} selects \code{meta-with-interrupt-gen}.
The generator itself belongs to \code{[group-of group]}, the surrounding
group's meta-group. For each application object it constructs a
non-reifying evaluator wrapper and a non-reifying metaobject. These
machinery objects use direct message handlers, while the application
script can use reflective forms.
The macro exposes the protocol:

\wholelisting[firstline=190,lastline=192]{examples/misc/reify-args.abclr2}

Its expansion sends
\code{[:with-interrupt [Body Success Fail] Env Id Gid Eval] @ C}
to the current application's metaobject. \code{Env} is the captured
evaluation environment; \code{Id} identifies that application's
metaobject; \code{Gid} is its group context; \code{Eval} is its current
evaluator wrapper. \code{C} receives the result of the control form.
In this worker, that continuation leads to the outer \code{!}; it is
not simply the controller's application reply destination \code{R}.

Read the metaobject's control handlers in this order:

\wholelisting[firstline=105,lastline=121]{examples/misc/reify-args.abclr2}

On entry to \code{:with-interrupt}, it stores the failure representation
and the result/context tuple \code{[C Env Id Gid Eval]}. It requests
evaluation of \code{Body}; the attached continuation receives
\code{Result}, ignores that body result, and requests evaluation of
\code{Success} with destination \code{C}. The caller waits for one
selected branch, not for all three arguments.

An active metaobject accepts \code{[:interrupt]} and tells its evaluator
wrapper where to report an intercepted request. A dormant metaobject
only prints that interruption is ignored. The mode comes from the
ordinary message-dispatch lifecycle, not from a timer or a stack of
active \code{with-interrupt} regions.

\subsubsection{Capture the Next Request, Then Abandon It}

Here is the wrapper inside the original generator:

\wholelisting[firstline=78,lastline=93]{examples/misc/reify-args.abclr2}

The final delimiters also close the enclosing temporary-variable binding.
Initially \code{interrupted} is \code{nil}. A normal request \code{any}
is forwarded to the underlying group evaluator, retaining its reply
destination \code{R}. Here the capitalized \code{Eval} in the wrapper's
fallback comes from the generator's creation request: it names the
underlying evaluator. It is different from the incoming control
handler's \code{Eval}, which names the wrapper through which the
application is currently evaluated.

The \code{[:interrupt R]} handler saves its notification address and sets
a flag. It does not immediately stop host Lisp code or return a context.
The next ordinary evaluation request consumes the flag. Instead of
forwarding that request, the wrapper sends \code{[:context [any . R]]}
to the saved address. The dotted pair retains the request and its
completion destination. Interruption occurs at an evaluation-message
boundary; a host operation already executing is not preempted inside
its machine instructions.

The metaobject's \code{[:context \_]} handler intentionally ignores that
intercepted request. It evaluates \code{fail-form} using the separately
saved \code{context-for-recover}, with the same \code{C} as the control
form's original result destination. This bypasses the body continuation
that would have selected success. The interrupted computation is not
resumed later. Contrast this with the newer Express library, which
stores a suspended evaluation request and resumes it after the express
script has finished (Section~\ref{sec:original-express}).

Follow an early interruption through five distinct messages:

\begin{enumerate}
  \item The controller sends the worker \code{[:run n] @ R}.
  \item The worker's reflective form asks its metaobject to evaluate the
  body and remember a failure alternative.
  \item The controller sends that metaobject \code{[:interrupt]}; the
  metaobject sends its wrapper \code{[:interrupt Me]}.
  \item The wrapper intercepts a later evaluation request and reports
  \code{:context}; the metaobject substitutes failure evaluation at
  the control form's saved destination \code{C}.
  \item Failure returns through the worker's \code{!} to \code{R}.
  The script's separate execution-completion continuation eventually
  sends \code{:end} to its metaobject, advancing the ordinary queue.
\end{enumerate}

The remaining generator handlers reproduce the familiar queue protocol:
\code{:message} enqueues, \code{:begin} matches and starts one script,
and \code{:end} starts another request or returns to dormant mode.
State initialization uses \code{:do-evlis}, \code{pairlis}, and
\code{env-gen}, as in the inheritance walkthrough. Extra creation
parameters are matched by \code{Extra} but not used here. The
\code{:evaluator}/\code{:set-evaluator} handlers expose the evaluator
choice; they do not themselves define interruption.

\subsubsection{Effects, Scope, and Safe Reproduction}

Printed counts remain printed. A supplemental worker using the original
generator increments persistent state in its body and returns either
\code{(:SUCCESS count)} or \code{(:FAILED count)}. Its early interrupt
leaves a positive count below the requested 50 increments. A subsequent
ordinary state query sees that same count, and a new successful run adds
two. Thus branch substitution abandons further body work but does not
undo earlier effects or automatically reinitialize the worker.

The implementation has one \code{fail-form} and one recovery tuple per
metaobject, not a stack of nested regions. Use it around one body interval
rather than as a general nested cancellation API. There is no timeout clock, cleanup
protocol, or guarantee of arbitrary resource release.

For a separate executable version of this experiment, run
\code{make -C tutorial interrupt-check}. Then load the prepared definitions
in a fresh ABCL/R2 session:

\begin{codeblock}
\begin{lstlisting}
(inc "build/interrupt-tests/interrupt-excerpt" :use-source t)
\end{lstlisting}
\end{codeblock}

The excerpt contains the interrupt generator, its group, the worker and
the controller. Thirteen checks exercise the two controller calls,
empty and normal bodies, reuse, and the supplemental state experiment.
Sources, translations, inputs and logs stay in
\filepath{build/interrupt-tests/}.

\paragraph{Check your understanding.}
Why not use \code{reify-args} here? The handler is on the application's
metaobject, not the group evaluator. Why does the failure path not send
the intercepted request back to the wrapper? It substitutes another
expression rather than resuming that request. Does \code{m = 3} mean
three seconds? No. Does interruption restore the count? No; prior
effects remain. Why do both success and failure use \code{@ C}? Both
must complete the same control expression and resume its caller's
remaining evaluation, including the application's reply.


\subsection{Inheritance Without a Superclass Declaration}
\label{sec:original-inherit}

The control-form projects changed which expressions run. The original
\filepath{examples/misc/inherit.abclr2} now uses the same reflective
access to choose behavior from another object.
Familiar OOP gives us the question: can an emu reuse a bird's
behavior while overriding its flight behavior?
ABCL/R2's answer here is an explicit \emph{script-acquisition protocol},
not a class declaration or a built-in superclass search.

\subsubsection{Start with the Application, Not the Generator}

Load the original file in a fresh session:
\begin{lstlisting}
(inc "examples/misc/inherit" :use-source t)
\end{lstlisting}
The file defines a macro, a generator, and a group. Loading it does not
create bird or emu; enter their definitions from its closing comment:
\begin{codeblock}
\begin{lstlisting}
[object bird
 (script
  (=> :has-feather? !t)
  (=> :can-fly? !t))
 (group inherit-group)]

[object emu
 (script
  (=> :can-fly? !nil)
  (=> Any @ Reply from Sender
      (inherit Any Reply Sender bird)))
 (group inherit-group)]
\end{lstlisting}
\end{codeblock}

Both objects belong to \code{inherit-group}, which arranges for the
custom metaobject generator to be used.
The bird answers two questions locally. The emu answers the flight
question locally and forwards every other request through \code{inherit}.
The source's sample results are:
\begin{lstlisting}
[bird <== :can-fly?]           ; => T
[emu <== :can-fly?]            ; => NIL
[emu <== :has-feather?]        ; => T
\end{lstlisting}
\code{nil} in the second result is a reply containing false, not a
missing reply. The local emu handler sends it explicitly.

Check the identities separately from those application results:
\begin{lstlisting}
(eq [den [meta emu]] emu)             ; => T
(eq [meta emu] emu)                   ; => NIL
\end{lstlisting}
Emu's metaobject is another object, and its denotation leads back to emu.
Neither operation is a superclass lookup. These identity checks and
the three application replies pass on the current image.
Here \code{[group-of emu]} returns \code{inherit-group}. The generator
exposes its creation argument as \code{Group} for identity creation,
retaining \code{Group-manager} for execution. MCP, Meta-Meta, Fault-Tolerant,
and the Time-Warp application generator use the same convention.
This is application ownership, not the group of its metaobject or inherited method source.

For non-reifying metaobjects, the generator assigns its known creation
group to the denoted identity before initialization. The host expression
\code{(setf (r2sys::object-data-gid [den Metaobject]) GMgr)} writes that
identity's group slot; it does not inspect or expose the NRM state.

For the feather request, \code{Any} becomes \code{:has-feather?}.
\code{Reply} holds the caller's reply destination and \code{Sender}
holds the request's sender. These are context fields, not extra elements
inside the keyword message.
The fallback asks for bird's matching script with that context.
It does not first send the ordinary question \code{:has-feather?}
to bird and then copy bird's answer.

\subsubsection{The Reflective Boundary}

The macro from the Lisp reading guide expands the fallback call into:
\begin{lstlisting}
(reify-args-to-meta :inherit Any Reply Sender bird)
\end{lstlisting}
\code{reify-args-to-meta} sends compiled representations of these four
argument expressions, with their evaluation context, to emu's metaobject.
That is the object controlling emu's script execution.
The keyword \code{:inherit} selects a handler added by the custom generator.
The representations are not yet the four application values.

Read the core original handler now:
\wholelisting[linerange={33-35,38-48}]{examples/misc/inherit.abclr2}

The handler proceeds in three stages:
\begin{enumerate}
  \item \code{:do-evlis} asks the evaluator to evaluate the argument
        representations in their original environment. Its continuation
        receives the list \code{[message reply sender inherit-source]}.
  \item It asks the source's metaobject to find a script matching that
        message, reply destination, and sender. A successful reply contains
        pattern bindings and an executable script representation.
  \item It extends the current environment with those bindings and asks
        the evaluator to execute the acquired script in the current
        execution context, forwarding the meta-level continuation \code{C}.
\end{enumerate}
\code{Env} supplies the current bindings; \code{Id} identifies the
execution context used for \code{Me}; \code{Gid} supplies its group context;
\code{Eval} names the evaluator.
The acquired feather script sends \code{t} to the application's
\code{Reply}, answering the feather question. Separately, \code{C}
tells the reflective computation that the acquired expression has finished.
If no script is found, the failure branch warns and completes through
\code{C} with \code{nil}. It does not thereby send an application failure
reply to \code{Reply}.

\subsubsection{Reconstruct the Metaobject's Ordinary Dispatch}

Return to the generator's state and its normal-message handlers:
\wholelisting[firstline=23,lastline=30]{examples/misc/inherit.abclr2}

The application's initializers are described by \code{StateSpec}.
\code{LexEnv} names its outer lexical environment and \code{Scripts}
its script list. The group's execution components are \code{Eval}
and \code{GMgr}. The dotted \code{Extra} collects additional parameters.
The generated metaobject has its \emph{own} state: application queue,
application environment, scripts, evaluator, group manager, and mode.
The variable \code{state} is initially \code{nil}; it is filled by the
initialization code, not assumed to be the application's data already.

Read lines 49--71 of the original file with this state machine in mind:
\begin{enumerate}[nosep]
  \item \code{[:message Message Reply Sender]} queues one request.
        A dormant metaobject marks itself active and schedules \code{:begin}.
  \item \code{:begin} dequeues the request and finds its first matching
        script. The evaluator executes it in an environment containing
        the pattern bindings over the application's state.
  \item The execution continuation sends \code{:end} to the metaobject.
        That handler starts another queued request, or becomes dormant.
\end{enumerate}
\code{Me} in these metaobject handlers denotes the metaobject itself.
\code{[den Me]} obtains the application object it describes.
The custom \code{:inherit} handler is added to this machinery; it does
not replace the ordinary application inbox by a superclass hierarchy.

\subsubsection{Initialization and Reflective Queries}

Lines 73--92 expose lookup and update operations such as \code{:state},
\code{:scripts}, \code{:find-script}, and \code{:set-evaluator}.
In particular, the three-field \code{:find-script} query preserves sender
and reply information; the one-field query uses \code{nil} for those fields.
That difference matters for patterns that depend on the sender.

Finally, lines 95--104 evaluate state initializers with \code{:do-evlis}.
\code{stateSpec-varNames} supplies names and \code{pairlis} associates
them with the resulting values. \code{env-gen} creates an environment
with those bindings and the lexical parent \code{LexEnv}.
The generator sets that environment on the new metaobject and replies
with the metaobject to its creation caller.
Read initialization last, once you know what its state must support.

\paragraph{Understanding check, with answers.}
Why does emu's \code{:can-fly?} not consult bird?
Its specific local pattern is first and matches, so the catch-all never runs.
Would moving the catch-all first preserve that override? No.
Does the example copy bird's state or permanently attach bird's scripts
to emu? No: it finds a script for each inherited request and executes it
using the current environment and context. It is not a guarantee of the
field or superclass semantics of another OOP language.
Why is \code{from Sender} retained? Acquired script matching may depend
on who sent the original request, not just its payload.

\paragraph{Reading milestone.}
You should now be able to locate the application reply, the expression
completion, the queue transition, and the creation reply independently
in this file. The original baseline checks the three bird/emu queries
without changing any original expression.
That success does not certify arbitrary scripts, runtime script
mutation, or every use of the generator.


\subsection{Reading the Built-in Protocol Sketches}
\label{sec:builtin-reading}

Inheritance used an environment, scripts, and a queue. Before reading
another metaobject generator, inspect those components in
\filepath{examples/misc/built-ins.abclr2}. It describes environments,
script objects, and queues in familiar object notation. It is a
\emph{protocol sketch}, not a replacement runtime to load:
angle-bracket placeholders, spelling mistakes such as \code{Nmae}
and \code{Vlaue}, an extra bracket, and the \code{:bound} fallback
prevent that interpretation. Consult \filepath{src/primary-env.lsp}
and \filepath{src/primary-queue.lsp} for the supplied implementations.

An environment is an object representing variable bindings. Its local
frame is an association list: \code{((n . 2) (limit . 10))}
binds \code{n} to 2. A second environment, \code{lexEnv},
represents the enclosing lexical scope.

\wholelisting[firstline=9,lastline=20]{examples/misc/built-ins.abclr2}

Trace a lookup of \code{limit}: \code{assoc} finds its pair, and
\code{cdr} supplies 10 to the reply destination. If no local pair exists,
the request goes outward, preserving \code{@ R}.
A local binding whose value is \code{nil} still counts:
the test examines the pair, not its value. Assignment by \code{:set}
updates the nearest existing binding; it does not necessarily create a
new local binding. The actual implementation also accepts
\code{[:set-local Name Value]}, updating or creating a binding in
this frame. The evaluator uses this when naming a new object.
\code{:alist} and \code{:lexEnv} inspect the frame and its parent.
The global environment ends the chain and reports an unbound lookup.
The sketch's outer \code{:bound} request is a typo, not a second
protocol: the implementation forwards \code{:boundp}.

A script object separates two jobs usually hidden by a handler
declaration. \code{[:match Message Reply Sender]} returns either
\code{nil} for failure or \code{[Bindings ExecutableExpression]} for
success. The bindings are an association list; the executable expression
may be \code{[:script ScriptObject]}. The evaluator then asks that object
to execute with \code{[:exec Env Id Gid Eval] @ C}. Here \code{Env}
contains the bindings, \code{Id} is the current metaobject identity,
\code{Gid} the group manager, \code{Eval} the evaluator, and \code{C}
the execution continuation. Pattern and body inspection are separate
requests, not commands to run the script.

The queue sketch stores arrivals by \code{cons} at the front and removes
the last element. Arrivals A, then B produce \code{(B A)}, but dequeue
returns A first: this is \emph{first in, first out} (FIFO), not a stack. Its \code{waiters}
list similarly retains reply destinations for dequeue requests made
while empty. Such a request need not stop the queue from accepting a
later enqueue. The built-in queue uses lower-level machinery instead;
do not infer that its storage representation is this pair of lists.
It additionally provides \code{:listify}, used by the metaobjects below.

\paragraph{Check your understanding.}
An inner frame is \code{((n . nil))}; its parent binds \code{n} to 7.
What does looking up \code{n} return? \emph{Answer:} \code{nil},
because the local binding shadows the parent's binding.
Would storing A with \code{cons} make it the next queue item?
\emph{Answer:} only if no older item exists; removal is from the
opposite end.

\subsection{A Script Object with a Sender-sensitive Match}
\label{sec:user-script-reading}

The sketch separated script matching from execution; now follow an
executable script object in \filepath{examples/misc/user-script.abclr2}.
Loading it creates four objects; the installation commands at the end
are inside a block comment and must be issued separately.
\code{Monkey} initially replies \emph{Scratch!} to any message.
\code{Monkey-Showman} and \code{Poor-Audience} each handle \code{:try}
by sending \code{:play-trick} to Monkey and forwarding their caller's
reply destination. Monkey sees different \emph{senders}
even though the requests have the same content.

The fourth object, \code{monkey-script}, implements the script-object
contract. It is not Monkey's application handler until installed
in Monkey's script set. For \code{:play-trick}, its match method tests
\code{(eq Sender Monkey-Showman)}. On success, it returns two elements:

\begin{codeblock}
\begin{lstlisting}
[ [ ['R . Reply] ] [:script Me] ]
; First element: association list binding R to Reply.
; Second element: expression requesting this script's execution.
\end{lstlisting}
\end{codeblock}

These are nested lists, not three bindings. \code{R} is quoted
because it is a binding name; \code{Reply} is evaluated because its value
is the destination being bound. Here \code{Me} is the script object,
not Monkey or the Showman.
For a different sender the match method replies \code{nil}.
The next catch-all \code{:match} method also replies \code{nil};
without it, script search could wait rather than move on to
Monkey's original catch-all handler.

\wholelisting[firstline=51,lastline=66]{examples/misc/user-script.abclr2}

The execution method asks the environment for \code{R}. Its continuation
sends \emph{Hansei} to that application reply destination.
Separately it sends \code{nil} to \code{C}, completing the evaluator's
request. These destinations have different meanings.
There is an ordering boundary: the lookup uses \code{<=}, not
\code{<==}. The source does not wait for its continuation to deliver
\emph{Hansei} before sending completion to \code{C}.
Completion is not a receipt from the application caller.

Try these commands in a fresh session, preferably with a scratch copy:

\begin{codeblock}
\begin{lstlisting}
(inc "examples/misc/user-script" :use-source t)
[Monkey-Showman <== :try] ; "Scratch!"
[Poor-Audience <== :try]  ; "Scratch!"
[[meta Monkey] <= [:add-script monkey-script]]
[[meta Monkey] <== :scripts] ; inspect installation
[Monkey-Showman <== :try]
[Poor-Audience <== :try]
\end{lstlisting}
\end{codeblock}

The final replies are \emph{Hansei} and \emph{Scratch!}.
Installation puts the new script first. It handles the Showman's request;
its failed match for Audience lets search reach Monkey's original handler.
Removing a script with \code{[:remove-script message]} removes only the
first script matching that message with \code{nil} reply and sender fields,
not every matching script or a script selected by identity. The custom
sender test here matters: a removal request with a \code{nil} sender does
not select \code{monkey-script}. Script-list changes affect later dispatch,
not a handler already being evaluated.
\code{:add-script} does not send a reply: changing its arrow to
\code{<==} waits for a reply which the protocol never promises.

\paragraph{Check your understanding.}
How can the script distinguish Showman and Audience without adding a
field to \code{:play-trick}? \emph{Answer:} matching receives the sender
as well as content. Why must a failed match reply?
\emph{Answer:} search needs an explicit failure before trying another
candidate.

\subsection{MCP: Following a Meta-circular Interpreter}
\label{sec:mcp-reading}

We have followed individual environment and script protocols. Now combine
them: \filepath{examples/misc/MCP.abclr2} expresses a metaobject generator,
an evaluator, and a group manager in ABCL/R2 itself.
\emph{Meta-circular} means that the interpreter is described using
the language it interprets. It does not mean this file removes all
dependencies on host Lisp and existing kernel objects.
Read it in three passes: the group manager at the bottom, the generated
metaobject, and finally the evaluator's expression cases.

The ordinary object \code{Meta-circular} implements the group-manager
protocol. Its state selects \code{MC-meta-gen} and \code{MC-eval}.
Its \code{:new} handler delegates to the generator, adding the
evaluator and \code{Me} as group manager.
\code{:replace-meta-gen} and \code{:replace-evaluator} change these
configuration fields, not necessarily every previously created object.
\code{:new-group} delegates to the existing group machinery and sends
\code{:initialize} to the resulting manager.

The generator receives state specifications, lexical environment,
scripts, evaluator, and group manager. Extra arguments are warned about
and ignored. It creates a metaobject holding queue, state environment,
script set, evaluator, manager, and a mode initially \code{:dormant}.
Initializer expressions are evaluated by \code{:do-evlis};
\code{pairlis} pairs variable names with the returned values.
An environment extending \code{LexEnv} becomes the state, and the
generator replies with the metaobject.

For application traffic, follow this cycle:

\wholelisting[firstline=33,lastline=43]{examples/misc/MCP.abclr2}

\begin{enumerate}
\item \code{[:message Message Reply Sender]} enqueues the three fields.
If dormant, the metaobject becomes active and sends itself \code{:begin}.
\item \code{:begin} dequeues one request, finds its script, and extends
the state environment with match bindings.
\item The evaluator executes the body. Its completion continuation sends
\code{:end} to the metaobject.
\item \code{:end} starts another queued request or returns to dormant mode.
\end{enumerate}

Application replies use \code{Reply}; advancement of the cycle uses the
execution continuation. Reflective handlers expose state, scripts, queue, and
evaluator. The one-field \code{[:find-script Message]} query supplies
\code{nil} for reply and sender; ordinary dispatch uses the actual fields.

In \code{MC-eval}, \code{Exp} is a tagged representation, not necessarily
original source. \code{:compiled} invokes a compiled function with the
context. \code{:variable} looks up ordinary names in \code{Env};
\code{Me} and \code{Group} yield \code{[den Id]} and \code{Gid}.
\code{:self-evaluate} invokes a supplied function.
\code{:script} invokes script-object execution.
An unknown representation warns; this is not an unrestricted Common Lisp
evaluator.

The send cases expose the difference between reply and completion:

\begin{longtable}{L{0.25\textwidth}L{0.65\textwidth}}
\textbf{Expression tag} & \textbf{Completion and delivery}\\\hline
\endhead
\code{:send-past} & Sends \code{nil} to \code{C}, then forwards the
request to the target's metaobject with its separate \code{Reply} field.\\
\code{:send-future} & Sends \code{nil} to \code{C}, obtains the future's
reply destination, and forwards the request with that destination.\\
\code{:send-now} & Uses \code{C} as the application's reply destination;
the target's reply supplies the expression's result.\\
\end{longtable}

These describe the non-null-target paths. A now-send to \code{nil}
does not reply in this definition.
\code{:object-def} chooses the specified group or the current one,
creates an object, and, if named, installs its denotation through
\code{:set-local}. An explicit generator receives the group's evaluator;
otherwise creation goes through the group manager.
\code{:group-def} creates and optionally binds a group manager.

\code{:do-progn} evaluates expressions sequentially, returning the last
result, or \code{nil} for an empty body.
\code{:do-evlis} evaluates sequentially and reconstructs a list of
results. Both use continuations to wait for each expression.

Here is a small teaching probe of the interpreter, not another original example:

\begin{lstlisting}
(inc "examples/misc/MCP" :use-source t)
[object mc-counter
  (group Meta-circular)
  (state [n := 0])
  (script
    (=> :increment [n := (+ n 1)] !n)
    (=> :value !n))]
[mc-counter <== :increment] ; 1
[mc-counter <== :increment] ; 2
[mc-counter <== :value]     ; 2
[[meta mc-counter] <== :evaluator] ; MC-eval
\end{lstlisting}

All four observations pass on the current image. They establish this
creation/state/reply path, not coverage of every evaluator case.

\paragraph{Check your understanding.}
Why can a past-send complete before the receiver answers?
\emph{Answer:} its evaluator sends \code{nil} to \code{C} independently
of the application reply. A now-send instead lets the receiver answer
directly to \code{C}.

\subsection{Meta-meta: Reflecting on the Message Interpreter}
\label{sec:meta-meta-reading}

\filepath{examples/misc/meta-meta.abclr2} takes the queue/dispatch
cycle one level higher. Its application object, \code{twice},
replies with \code{(+ x x)}. The interesting change is which objects
interpret its requests.

Read the four numbered source steps in order.
\code{the-meta-meta-gen} is placed in
\code{(group-of (group-of group))}, and \code{the-meta-gen} in
\code{(group-of group)}. Each \code{group-of} selects the group of
its argument, climbing the group relationship in this setup;
\code{group} is the current-group pseudo-variable.
The second generator creates a metaobject explicitly configured with
\code{(meta-gen the-meta-meta-gen)}. That object's own execution
is therefore interpreted by a custom meta-metaobject.
The third step creates \code{a-group} using \code{the-meta-gen};
the fourth puts \code{twice} in that group.

Keep three identities distinct:

\begin{lstlisting}
twice                  ; application object
[meta twice]           ; interprets application messages
[meta [meta twice]]    ; interprets interpreter messages
\end{lstlisting}

For \code{[twice <== 3]}, the first interpreter receives a
\code{:message} envelope with application content 3.
The higher interpreter sees messages involved in executing that
interpreter, including its begin/end cycle.
The generators retain the familiar state initialization, matching,
environment extension, and completion protocol.
Compare their repeated bodies with \code{MC-meta-gen}: logging is in
the generated objects' \code{:message} handlers, not merely in the
factories' creation handlers.

The actual labels are \emph{Meta-gen receievs message} and
\emph{Meta-meta receievs message}; spelling is the original's.
Initialization also produces traffic. Not every log line is another
application request, and numeric object identifiers are not stable.

\begin{codeblock}
\begin{lstlisting}
(inc "examples/misc/meta-meta" :use-source t)
[twice <== 3] ; 6, with meta-level diagnostic output
\end{lstlisting}
\end{codeblock}

The commented historical transcript refers to \code{foo}, although
the file defines \code{twice}, and one logging label differs.
Treat it as an old illustration, not the literal session to reproduce.
The isolated check confirms 6, both actual logging markers, and
distinct identities at all three levels. It does not require a fixed
log order or number of lines.

\paragraph{Check your understanding.}
Does a meta-metaobject make \code{twice} a subclass of its interpreter?
\emph{Answer:} no. These are separate objects with an execution-control
relationship, not an inheritance hierarchy.

\paragraph{Reproducing the checks.}
Run \code{make -C tutorial reflection-check}. Eighteen checks use
distributed copies of \code{user-script}, \code{MCP}, and \code{meta-meta};
Their metaobjects retain distinct application and execution-group bindings.
Inputs, translations, and logs stay in \filepath{build/reflection-tests/}.
An additional 38 checks exercise addition, first-match removal, duplicate
scripts, reinstallation, retained state, and queued replies on both
primary dispatch paths. They guard script-list/cache consistency, not
arbitrary concurrent modification of an executing handler.


\subsection{Guarded Scripts: The Original Heavy-Meta Counter}
\label{sec:guards}

We now return from interpreter construction to a useful change in
application dispatch. A structural pattern can bind \code{n}, but it does not establish that
\code{n} is a positive integer. A script's \code{where} clause supplies
an additional acceptance test, called a \emph{guard}. This is part of
script selection, not an \code{if} inside an already selected body.
If the guard is false, the metaobject tries the next script for the
same request. The user's guide describes this facility as a library
extension~\cite{guide}; it does not add guards to the ordinary
\code{match} form.

The original example is inside the closing block comment of
\filepath{lib/heavy-meta.abclr2}, not in a separate example file:

\wholelisting[firstline=133,lastline=144]{lib/heavy-meta.abclr2}

Loading the library defines the generator \code{heavy-meta}, not this
commented counter. In a fresh session, load the library before entering
the counter definition above. Then replay its original commands:

\begin{codeblock}
\begin{lstlisting}
(inc-lib "lib/heavy-meta" :use-source t)
; Enter the original counter definition after loading the library.
[counter <= [:add 123]]
[counter <= [:add -100]]
[counter <== [:value]] ; => 123
\end{lstlisting}
\end{codeblock}

The two add requests are past-type sends. Neither add script replies;
do not replace them by now-type sends expecting an acknowledgement.
The value request is now-type and explicitly replies with \code{c}.
The reset request likewise changes state without replying.
The automated \code{guards-check} extracts original lines 133--144
byte-for-byte and replays lines 146--148 separately, leaving the
selected code unchanged.

\paragraph{Trace both structural success and guard failure.}
For \code{[:add 123]}, the first pattern binds \code{n = 123}.
\code{integerp} is true and \code{plusp} is true, so the body adds 123
to \code{c}. For \code{[:add -100]}, that same pattern still matches
and binds \code{n = -100}. Only the guard fails. The metaobject then
tries \code{[:add \_]}, whose wildcard accepts the argument without
binding it. That script has no body and intentionally ignores the
request, but its execution still finishes the dispatch cycle.

The \code{and} evaluates left to right and stops at the first false
result. Consequently an argument such as \code{"7"} fails \code{integerp}
without asking \code{plusp} to compare a string. Zero is an integer
but not positive, so it also reaches the fallback. Lisp's \code{nil}
is false; other values, including zero, are true when used directly
as guard results. A guard need not return exactly \code{t}.

For \code{[:value]}, neither add pattern matches; their guards are
not evaluated at all. The value script has no guard and replies
immediately once selected. Ordering still matters: placing the
unguarded add fallback first would accept every add request before
the guarded script could be considered.

\paragraph{Ensure that the object actually uses heavy-meta.}
The compiler records the guard separately from the body. In the
ordinary counter definition, it selects \code{heavy-meta} automatically
because neither a metaobject generator nor a group is specified.
The current \filepath{src/translate.lsp} makes that condition explicit.
The guide's shorter description of automatic selection should not be
read as overriding any custom group or metaobject configuration.

If you specify a group, explicitly select a compatible generator:

\wholelisting[firstline=26,lastline=32]{tutorial/examples/09-guards.abclr2}

Loading the library alone does not replace the generator used by every
group. A diagnostic object with an explicit ordinary group but no
\code{meta-gen} accepts its first structural match even when its stored
guard is \code{nil}; the ordinary metaobject does not perform this
guard protocol. Likewise, the non-reifying object translator warns
that a \code{where} field is ignored. Do not use a non-reifying object
as a guarded receiver. If another reflective generator is needed,
it must implement compatible guard-aware selection; naming both
features does not compose their algorithms automatically.

\paragraph{Read the library's continuation chain.}
The heavy metaobject retains the familiar ordinary queue, application
state environment, scripts, evaluator, group manager, and active/dormant
mode. Its \code{:begin} dequeues one request as \code{MRS}, a triplet
of message, reply destination, and sender. Selection is then a series
of asynchronous requests and result continuations:

\begin{enumerate}
  \item Ask the candidate script \code{[:match . MRS]}. Structural
  failure immediately moves on to the remaining scripts.
  \item On success, the result contains \code{[Bindings ScriptBody]}.
  Ask \code{env-gen} for a frame containing those bindings, with the
  object's persistent state as its outer environment.
  \item Ask the script for \code{:where-exp}. A \code{nil} reply means
  there is no guard. Otherwise this reply is a compiled expression
  representation, not its boolean result.
  \item Evaluate that representation with \code{[:do ...]} in the new
  environment. Its result continuation either starts \code{ScriptBody}
  in that same environment or tries the remaining scripts.
  \item Attach \code{[cont \_ [Me <= :end]]} to body execution.
  Completion advances the ordinary queue or restores dormant mode.
\end{enumerate}

The central decision is visible in the original implementation:

\wholelisting[firstline=74,lastline=89]{lib/heavy-meta.abclr2}

Its two \code{if} tests answer different questions: whether a guard
representation exists, and whether evaluating it accepted the request.
An explicitly false guard such as \code{where nil} still has a non-nil
representation; it must be evaluated and rejected, not treated as absent.
The completion continuation is also separate from the script's
application reply. The empty fallback can finish execution without
sending anything to the original caller.

The library first presents a synchronous sketch under
\code{\#+COMMENT}, then the active continuation-passing version.
Continuation-passing style, abbreviated CPS in its comments, makes
each next step a handler of the preceding result. Read the active
chain as the implementation; the sketch is explanatory, not an
additional code path executed before it.

\paragraph{Which bindings can the guard see?}
The guard runs after pattern matching and environment creation. It
can use pattern variables, persistent state, outer lexical bindings,
and the bound reply/sender fields of its script header. The supplemental
example checks the last two explicitly:

\wholelisting[firstline=15,lastline=24]{tutorial/examples/09-guards.abclr2}

In the first driver request, \code{Me} denotes the driver, which is
also the sender received as \code{S}. The guard succeeds and the body
confirms that \code{R} is a non-nil reply destination. The other request
passes \code{nil} as the expected sender, so the guard fails and the
fallback replies \code{nil}.

Body-local \code{temporary} bindings are different: their initialization
belongs to body execution, after guard selection. Do not refer to those
bindings as though they were available before the body starts.
The parser's header order is pattern, optional \code{@ R}, optional
\code{from S}, then optional \code{where expression}.

\paragraph{A failed guard does not undo its effects.}
The supplemental \code{guarded-choice} increments \code{checks} inside
its guard before testing the argument against state binding \code{limit}.
Requests for 5, 11, and the symbol \code{wrong} return
\code{:rejected}, \code{:accepted}, and \code{:rejected}, while
\code{checks} becomes 3. Even failed candidates have executed their
guards. The example demonstrates the mechanism, not a recommendation
to hide side effects in acceptance tests.

\wholelisting[firstline=4,lastline=13]{tutorial/examples/09-guards.abclr2}

Prefer simple predicates. A guard that waits on an ordinary request
to its own receiver can deadlock: that request cannot start while
the current guard is still selecting a script. Other objects can also
change shared state; guard evaluation is not a multi-object transaction.

Reflective script lookup is not the same as guard-aware dispatch.
The heavy metaobject's \code{[:find-script Message]} query calls
\code{search-script}, which checks structural matches. In the original
counter it returns the first add script even for \code{[:add -1]}.
Its \code{:where-exp} remains available for inspection, but the lookup
query does not certify that evaluating that guard would accept the request.

\paragraph{Evidence and exercise.}
Run \code{make -C tutorial guards-check}. Twenty-four checks cover the
original counter, script/guard inspection, state and sender/reply bindings,
failed-guard effects, explicit-group setup and non-reifying restrictions.
Translations, extracted sources and transcripts stay in
\filepath{build/guard-tests/}. The runnable supplemental definitions are in
\filepath{tutorial/examples/09-guards.abclr2}.

Why does \code{[:add -100]} not raise an error or block later value queries?
Its structural match succeeds, its guard fails, and the empty fallback
completes dispatch. Why not put the same test inside the first body's
\code{if}? That would select the first script regardless of the test;
it would not ask dispatch to try the next script. Why must a guard's
side effects be considered separately from acceptance? Evaluation has
already happened when its false result selects another candidate.

\paragraph{Which Heavy-Meta library should you load?}
Use \filepath{lib/heavy-meta.abclr2} for guarded programs. The distribution
also contains \filepath{lib/heavy-meta-old.abclr2}, a historical variant
whose guard protocol does not match the current translator. It is not
needed by the original programs in \filepath{examples/} and is not an
interchangeable alternative. Both files define \code{heavy-meta} and
provide the feature \code{heavy-meta}; do not load both into the same
session. The old file is historical reference material, not another
step in this book's learning path.


\subsection{Multicast: One Request, a Structured Result}
\label{sec:multicast}

An ordinary send has one receiver. The original \filepath{lib/multicast.abclr2}
adds a library-defined operation: send the same message to every object in
a target tree, and optionally reconstruct a tree of replies. This is not
network multicast, a transaction, or the unsupported brace-delimited send
syntax. Its targets are ordinary ABCL/R2 objects in Lisp cons structures.
Where guards changed the receiving object's script selection, multicast
changes how the sender's evaluator interprets a target. Keep the future
reply-stream model and the evaluator context from earlier sections in mind.

\subsubsection{Replay the Original Library Example}

The library ends with a commented example, rather than a separate file in
\filepath{examples/}. Its three objects are:

\wholelisting[firstline=173,lastline=183]{lib/multicast.abclr2}

In a fresh session at the repository root, load the library, wrap the
top-level evaluator, enter those definitions, and replay the original send:

\begin{codeblock}
\begin{lstlisting}
(inc-lib "lib/multicast" :use-source t)
[*eval* := [multicast-eval-gen <== [:new *eval*]]]
; Enter the original twice, square, and identity definitions above.
[[twice square identity] <== 12]
; => (24 144 12)
\end{lstlisting}
\end{codeblock}

The bracket expression \code{[twice square identity]} constructs a list
of three object values. Each receives 12: \code{twice} returns 24,
\code{square} returns 144, and \code{identity} returns 12. The result list
uses the targets' positions, not the order in which replies arrive.
The same top-level installation is documented in the user's guide~\cite{guide}.
Loading the library alone only defines its machinery; it does not enable
list-valued targets in every evaluator.

\subsubsection{Follow the Shape, Not Just the Values}

Try a target tree containing both objects and literal data:

\begin{codeblock}
\begin{lstlisting}
[[twice [square :label] identity nil] <== 12]
; => (24 (144 :LABEL) 12 NIL)
[[twice twice] <== 3]                  ; => (6 6)
\end{lstlisting}
\end{codeblock}

Only object leaves receive the message. The keyword \code{:label} and
\code{nil} are retained as data; they are not receivers. Repeating an
object creates two target occurrences and two sends, not a deduplicated
recipient set. Each occurrence has its own result position.

The host Lisp function \code{pickup-target-objects} recursively walks
\code{car} and \code{cdr}. It constructs a fresh cons template with
\code{nil} in every object position, and records each object with a path
into that template. For the nested example:

\begin{longtable}{L{0.21\textwidth}L{0.37\textwidth}L{0.32\textwidth}}
\textbf{Receiver} & \textbf{Path from the template root} & \textbf{Stored reply}\\\hline
\code{twice} & \code{(:car)} & 24\\
\code{square} & \code{(:cdr :car :car)} & 144\\
\code{identity} & \code{(:cdr :cdr :car)} & 12\\
\end{longtable}

The template is \code{(nil (nil :label) nil nil)}. A path is a sequence
of operations: \code{(:cdr :car :car)} means take the outer tail, its
first nested list, and that list's first cell. The function builds its
target records with \code{push}, so their order is reversed relative to
the traversal. That does not change the reconstructed result.
\code{store-returned-value} follows a path and updates the final
\code{car} or \code{cdr} with \code{setf}.

This is a cons-tree walk, not arbitrary container traversal. A dotted
pair of object targets also works: \code{(cons twice square)} with message
3 yields \code{(6 . 9)}. Arrays are not recursively decomposed.
Use finite, acyclic target trees. A receiver that sends \code{!nil}
has replied; a \code{nil} result is not the same as a missing reply.

\subsubsection{Separate Submission from the Aggregate Reply}

The evaluator wrapper intercepts two ordinary send representations:

\wholelisting[firstline=134,lastline=143]{lib/multicast.abclr2}

For a nonempty cons target, a now-type request sends the evaluator's
continuation \code{C} as the aggregate reply destination. A past-type
request instead supplies its explicit application destination \code{Reply},
and preserves \code{C} separately as the submission continuation.
Single-object targets and other evaluator requests are forwarded to
\code{original-eval}. Future-type sends and extended send representations
are not intercepted by these two patterns; do not assume that every
extension composes automatically with multicast.

Read the \code{multicaster} handler in this order:

\begin{enumerate}[nosep]
  \item Decompose the target tree into template and object/path records.
  \item Send \code{nil} to its own continuation \code{C}, if present.
        This is submission progress, not a collected application result.
  \item Without an aggregate destination, send the payload to each object
        and collect nothing.
  \item Otherwise create one collector, initialize it with the template,
        target count, and aggregate destination, then send to each object
        with a position-specific reply receiver.
\end{enumerate}

\code{partial-value-receiver-gen} creates a \code{cont} object for each
position. An application reply becomes \code{[:result value position]}
to the collector. Notice another forwarding effect: the target's sender
binding is \code{multicaster}, not the original application client.
Sender-sensitive scripts must take this into account.

\subsubsection{Put a Multicast Client in Its Own Group}

The supplemental \filepath{tutorial/examples/10-multicast.abclr2}
keeps this policy local to a client group instead of replacing the
interactive evaluator:

\wholelisting[firstline=4,lastline=17]{tutorial/examples/10-multicast.abclr2}

The targets may remain ordinary objects: it is the \emph{sender's}
evaluator that interprets a list-valued target. The internal
\code{r2sys::create-primary-eval} helper is the same implementation-specific
helper used earlier. In a fresh session, enter the original three object
definitions, then load the supplemental file:

\begin{codeblock}
\begin{lstlisting}
(inc "tutorial/examples/10-multicast" :use-source t)
[multicast-client <== [:now [twice square] 3]]
; => (6 9)
[multicast-client
 <== [:past [twice square] 3 multicast-inbox]]
; => :SUBMITTED
[multicast-inbox <== :status]           ; => (1 (6 9))
\end{lstlisting}
\end{codeblock}

The inbox stores an arrival count and the last aggregate. Inspect its
short definition in the supplemental file. \code{:submitted} is the
client handler's reply to the interactive caller; \code{(6 9)} travels
separately to the inbox. The submission reply does not acknowledge that
every recipient has finished. The checked session observes the aggregate
after the finite sends have run; it does not establish an arrival-time
guarantee for other schedules.

For an existing group, the original \code{multicast-installer} accepts
\code{[:install group-manager]} as a past-type request. The manager must
support \code{:replace-evaluator}. Install before creating clients that
will use multicast: replacing a group's evaluator does not automatically
rewrite every existing object's captured evaluator. The guide's
\code{*eval*} assignment is the explicit top-level alternative.

\subsubsection{Know What the Collector Can and Cannot Promise}

The essential collector handlers are:

\wholelisting[firstline=72,lastline=82]{lib/multicast.abclr2}

Initialization stores the template, remaining-result count, and reply
destination. Here \code{C} names the numeric count, not the evaluator
continuation used earlier; \code{R} is the aggregate destination.
Each \code{:result} stores a value and decrements the count.
Only the transition to zero sends the aggregate. This assumes
\emph{exactly one reply per target occurrence}.

A nonempty target list with no object leaves initializes a count of zero,
but initialization sends no aggregate. A silent recipient leaves the
count above zero. Neither case has a library timeout or failure result;
a now-type caller cannot obtain the missing aggregate. An empty list is
\code{nil}, not a cons, so it does not enter the multicast branch at all.
Use a nonempty set of object recipients for collected requests, and make
each recipient's one-reply contract explicit.

The collector does not remember which positions have already replied.
A controlled probe initializes two positions, supplies two results for
the first position, and observes a premature \code{(9 nil)} aggregate.
A later result for the second position mutates that same delivered list
to \code{(9 10)} without a second aggregate message. These are observed
protocol limits, not a supported duplicate-delivery or retry mechanism.
Do not use a multi-reply stream producer as a one-reply multicast recipient.
For fire-and-forget sends, replies are not collected and there is no
aggregate completion acknowledgment.

\paragraph{Evidence and exercise.}
Run \code{make -C tutorial multicast-check}. Twenty-four checks replay
the original commented session and test nested/dotted targets, duplicate
occurrences, \code{nil} replies, sender identity, group setup, explicit
reply routing, fire-and-forget effects, and the collector boundaries above.
The runner extracts original lines 171--185 and copies the library
byte-for-byte; inputs and transcripts stay in \filepath{build/multicast-tests/}.
No original source or runtime is patched.

Why does \code{:label} survive in the result? It occupies a literal
template position rather than an object position. Why are two instances
of \code{twice} not merged? Positions, not distinct identities, define
the requests. As an exercise, add a one-reply negation object to a nested
target tree, predict the template and paths, then compare the collected
result. Keep the original library unchanged.


\subsection{Object Copies Are Not Independent Computations}
\label{sec:copy}

The original \filepath{lib/copy.abclr2} creates another object from an
existing object's scripts, state information, and pending messages.
For an OOP reader, the important question is not just ``is the identity
different?'' but ``which references and pending effects remain shared?''
This library also underlies \code{report} and can move an object into
the \code{watching} group defined by \code{watch}.

\paragraph{What is shared.}
Copying preserves local state values but is neither deep nor atomic: new local bindings retain the
same referenced data, while outer lexical environments, scripts, and queued
reply destinations remain shared. The examples below separate these effects;
do not interpret them as independent replay of an active computation.

\subsubsection{Read the Original Fibonacci Experiment}

The library's final block comment defines a recursive Fibonacci object.
It creates a collector for each nontrivial request and sends two requests
to itself: both target \code{Me}, not newly created worker objects.
The later Pcontrol project (Section~\ref{sec:pcontrol-fib}) uses newly
created children instead; you do not need that project to follow this one.

\wholelisting[firstline=164,lastline=179]{lib/copy.abclr2}

The collector captures \code{R}, the reply destination of the original
request. Its local \code{half-result} stores the first answer; the next
answer is added to it and sent to \code{R}. The base cases are 0 and 1.
Without copying an active computation, \code{[fib <== 5]} returns 5.
The library then proposes this experiment:

\wholelisting[firstline=181,lastline=187]{lib/copy.abclr2}

To replay it, load \code{lib/copy} with \code{inc-lib}, enter the
original \code{fib} definition, then enter this expression. The automated
\code{copy-check} target extracts both source ranges without editing them.
The \code{dotimes} loop is a scheduling opportunity, not a timed delay.
The copy can inherit requests already waiting for \code{fib}; those
requests still carry references to the original reply continuations and
collectors. Starting the copy can therefore deliver additional answers
into the original computation. The original comment warns that the result
can seem ``funny.''

Do not expect one fresh, independent Fibonacci answer from this session.
In the checked run the same \code{Fib(5)=} label appeared several times
with values including 4, 5, and 6. Those values are observations, not a
specified sequence to reproduce on every schedule. A collector designed
for two answers cannot distinguish extra answers caused by duplicated
work. A successful ordinary \code{[fib <== 5]} after the experiment does
not certify independent continuation copying.

\subsubsection{Make Sharing Visible with a Small Object}

The supplemental \filepath{tutorial/examples/11-copy.abclr2} puts a
counter binding in an \emph{outer lexical environment}, not in an object's
local \code{state} frame:

\wholelisting[firstline=4,lastline=10]{tutorial/examples/11-copy.abclr2}

The outer assignment retains the object in the top-level binding
\code{copy-source}. The name \code{copy-subject} inside the \code{let}
is local to that environment; a printed object name is not a guarantee
that a top-level variable of that name exists.
In a fresh session, try:

\begin{codeblock}
\begin{lstlisting}
(inc "tutorial/examples/11-copy" :use-source t)
[clone := [copier <== [:copy copy-source]]]
(eq clone copy-source)                 ; => NIL
[copy-source <== :increment]           ; => 1
[clone <== :value]                     ; => 1
[clone <== :increment]                 ; => 2
[copy-source <== :value]               ; => 2
\end{lstlisting}
\end{codeblock}

The objects and their metaobjects have distinct identities, but the
scripts are reused and the outer environment is retained. Both objects
therefore update the same binding \code{shared}. This is not a deep copy
of all reachable Lisp data or of captured application objects.
It is also not the proposed \code{COPY} environment semantics discussed
as future work in the historical guide.

\subsubsection{Follow the Copier's Construction Protocol}

Start with the \code{:group} handler, which the default copy request
also uses:

\wholelisting[firstline=44,lastline=54]{lib/copy.abclr2}

It first queries the original metaobject for its queue, state environment,
and scripts. It asks that environment for its local association list and
outer \code{:lexEnv}. Then \code{alist->stateSpec} builds an object
creation specification, and the selected group manager receives
\code{[:new stateSpec lexEnv scripts]}. The reply is a \emph{metaobject}.
The continuation gives the new object a display name, arranges its pending
messages through \code{message-setter}, and forwards the final object
to the copier's caller.

\noindent\begin{minipage}{\linewidth}
\begin{longtable}{L{0.30\textwidth}L{0.60\textwidth}}
\textbf{Part} & \textbf{Copy protocol}\\\hline
Object/metaobject identity & Creates new identities; renaming does not redirect
existing application references.\\
Scripts & Passes the existing script objects to the new construction.\\
Local state & Reconstructs an initializer specification from current values;
creates separate bindings but retains referenced data.\\
Outer lexical environment & Reuses the original environment reference.\\
Pending messages & Enqueues the original message/reply/sender records on the
copy; leaves the original queue in place.\\
Active execution & Does not clone an executing evaluator context or
recursively replace captured continuations and collectors.\\
\end{longtable}
\end{minipage}

These are several reflective requests, not an atomic snapshot protocol.
If an application continues to change while they run, do not assume the
collected queue and state describe one indivisible instant.
Custom metaobjects must support the queried reflection protocol.

\paragraph{How local-state values are preserved.}
For a binding \code{(n . 7)}, the helper captures 7 in a zero-argument
function and uses that function as the operand of \code{:self-evaluate}.
The primary evaluator calls it with \code{funcall} to obtain the value.
This also preserves a symbol or function as data rather than accidentally
calling it. The binding is new, not the referenced value: assigning the
copy's numeric slot leaves the source's slot unchanged; mutating a shared
list is visible through both. Rebinding that list slot does not rebind its
aliases. Initializer order and the outer environment are unchanged.

\subsubsection{Choose the Group and Evaluator Deliberately}

The library accepts these ordered message patterns. In the table,
\code{obj}, \code{g}, \code{mg}, and \code{ev} stand for evaluated objects:

\noindent\begin{minipage}{\linewidth}
\begin{longtable}{L{0.48\textwidth}L{0.42\textwidth}}
\textbf{Request to \code{copier}} & \textbf{Construction policy}\\\hline
\code{[:copy obj]} & Original group, using that group's current generator
and evaluator.\\
\code{[:copy obj :group g]} & Selected group's generator and evaluator.\\
\code{[:copy obj :meta-gen mg]} & Selected generator, original object's
current evaluator, original group.\\
\code{[:copy obj :meta-gen mg :eval ev]} & Selected generator and evaluator,
original group.\\
\end{longtable}
\end{minipage}

In particular, the simplest copy does \emph{not} necessarily retain an
object-specific evaluator override: it goes through the group's current
creation policy. The explicit \code{:meta-gen} form retrieves the source
object's evaluator instead. The tests distinguish these cases by comparing
evaluator identities. These are message patterns, not unrestricted Common
Lisp keyword arguments; retain the indicated order and combinations.

The macro \code{(copy name)} assigns the copy back to the binding
\code{name}. It also accepts \code{:group} or \code{:meta-gen}, but not both,
and has no \code{:eval} keyword of its own. For example, if
\code{[old-reference := copy-source]} precedes \code{(copy copy-source)},
\code{old-reference} still denotes the original object. Rebinding one name
does not update all references held by other objects.

\subsubsection{Copy Pending Work without Silently Starting It}

The original queue installer is short enough to read completely:

\wholelisting[firstline=150,lastline=157]{lib/copy.abclr2}

It sends \code{:enqueue} for each saved record and replies with
\code{[den newMeta]}, the new application object. The \code{:begin}
send is deliberately commented out. Copying pending work therefore does
not start that work; it also does not remove the original requests.
When a copy already has a nonempty queue, send \code{:begin} to its
metaobject to start it. Do not send \code{:begin} blindly to an empty queue.
An empty-queue copy can receive ordinary application requests normally.
An ordinary query is not a passive queue probe: it can start execution.
Inspect the metaobject's \code{:queue} instead when checking pending work.

The bounded queue exercise stages one pending record explicitly, so it
does not depend on interrupting a running computation:

\begin{codeblock}
\begin{lstlisting}
[object copy-inbox
 (state [values := nil])
 (script
  (=> :values !(reverse values))
  (=> value [values := (cons value values)]))]
[[meta copy-source]
 <= [:enqueue [[:echo 42] copy-inbox copy-source]]]
[queued-clone := [copier <== [:copy copy-source]]]
[copy-inbox <== :values]                ; => NIL
[[meta queued-clone] <= :begin]
[copy-inbox <== :values]                ; => (42)
[[meta copy-source] <= :begin]
[copy-inbox <== :values]                ; => (42 42)
\end{lstlisting}
\end{codeblock}

Here \code{[:enqueue ...]} takes the internal triple
\code{[message reply sender]}, not an ordinary application payload.
Both copies retain \code{copy-inbox} as the reply destination. The
original queue still held its request after the clone ran; starting it
produces the second 42. Enqueueing this way is a controlled reflective
exercise, not the recommended way to submit normal application work.
Do not interpret copying as exactly-once processing, migration that
deletes the source, or independent replay of an active computation.

\paragraph{Evidence and exercise.}
Run \code{make -C tutorial copy-check}. Fifty-two successful probes
cover the original Fibonacci session, identities, shared lexical bindings,
script reuse, queue copying/startup, alias rebinding, and group/evaluator
selection. Twenty-three state checks cover values, reference identity,
independent bindings, shared mutation, and both construction routes. Source extracts,
copies, inputs, and transcripts stay in \filepath{build/copy-tests/}.

Why do both copies increment \code{shared}? They retain the same outer
environment. Why does the inbox receive two replies? Both queues contain
a request addressed to it; only the execution start is separate. As an
exercise, keep the original and a named copy, then change which binding
you query. Predict identity and lexical effects separately. Do not infer
deep-copy semantics from separate object identities.


\subsection{Watch: Observe Sends without Replacing the Application}
\label{sec:watch}

The original \filepath{lib/watch.abclr2} is a practical debugging aid.
It wraps an evaluator, reports message-send requests, then forwards them
to the original evaluator. It does not copy the application object when
you use \code{watch}. In particular, it can observe an ordinary stateful
object without invoking the state-copying path discussed in
Section~\ref{sec:copy}.

The key distinction is \emph{outgoing sends}, not incoming requests.
Watching a caller reveals the requests it sends and the replies it sends
onward. It does not show every arrival at that caller, every operation
inside the receiver, or every host Lisp computation. The separate
\code{report} library observes a different boundary: message arrival and
script execution through a replacement metaobject.

\subsubsection{Start with One Sender}

The supplemental \filepath{tutorial/examples/12-watch.abclr2} defines
an echo target and a driver with now-, past-, and future-type sends:

\wholelisting[firstline=5,lastline=14]{tutorial/examples/12-watch.abclr2}

The file loads \code{watch} and \code{future}; the latter is needed for
this driver's future request, not for the basic watch facility.
Run this in a fresh session:

\begin{codeblock}
\begin{lstlisting}
(inc "tutorial/examples/12-watch" :use-source t)
(watch watch-driver)
[watch-driver <== [:now 7]]             ; => 7, plus trace output
(nowatch watch-driver)
[watch-driver <== [:now 8]]             ; => 8, no watch trace
\end{lstlisting}
\end{codeblock}

These are the original library's public macros. Common Lisp reader case
folding makes \code{noWatch} and \code{nowatch} the same macro name.
The trace contains a send to \code{watch-target} with value 7 and a send
of the reply onward to the driver's caller. The top-level request
\code{[:now 7]} itself was sent by the interactive evaluator, not by the
watched driver, so it is not the driver's outgoing request being reported.

The default reporter prints lines shaped like
\code{[target <== payload from sender]}. Object serial numbers are not
stable expected output. For a past-type send with a reply destination it
also prints \code{@ reply}. Its historical future display uses
\code{<<=}; that is a trace label, not a replacement for this book's
\code{[target <= payload \$ future]} source syntax.
Trace lines are asynchronous reports, not acknowledgments that a
receiver has finished processing the message.

\subsubsection{Read the Evaluator Wrapper}

Here is the complete original generator:

\wholelisting[firstline=37,lastline=55]{lib/watch.abclr2}

The dotted patterns bind the rest of the send representation as
\code{Arg}, and the remaining evaluation context as \code{Context}.
For each of the three send types, the wrapper sends the reporter four
fields: the type tag, evaluated send arguments, evaluation context, and
continuation \code{C}. It then forwards the original evaluator request
with the same continuation. The application reply path is not redirected
through the reporter.

For the standard send protocol, \code{Context} is
\code{(Env Id Gid Eval)}. Its second element is the sender's meta-level
identity; \code{[den Id]} denotes the application sender. This is how
the default formatter obtains the displayed sender. The wrapper also
answers \code{:original-evaluator}, which is useful for checking or
explicitly undoing one layer of instrumentation.

The catch-all forwards other evaluator requests unchanged. This is not
an instruction-by-instruction profiler, and operations performed by host
Lisp or a non-reifying script need not cross these evaluator patterns.
Observation adds messages and work of its own; do not infer uninstrumented
scheduling or timing from a watched run.

\subsubsection{Collect Structured Events Instead of Printed Names}

The macros select \code{default-watch-reporter}. To supply another
reporter, use the original manager's explicit protocol. The teaching
file includes this small recording object:

\wholelisting[firstline=16,lastline=22]{tutorial/examples/12-watch.abclr2}

Each event retains the same four fields the wrapper sends. The helper
\code{watch-target-events} filters those records by target identity;
it avoids comparing unstable printed names. In a fresh session after
loading the teaching file:

\begin{codeblock}
\begin{lstlisting}
[watch-manager <= [:watch-object watch-driver watch-recorder]]
[watch-driver <== [:now 7]]
[watch-driver <== [:past 8]]
[watch-driver <== [:future 9]]
[events := (watch-target-events
             [watch-recorder <== :events] watch-target)]
(mapcar #'first events)
; => (:SEND-NOW :SEND-PAST :SEND-FUTURE)
(nowatch watch-driver)
\end{lstlisting}
\end{codeblock}

There are additional events for future creation, future consumption, and
the driver's replies. Filtering for \code{watch-target} isolates the
three intended sends. The driver returns 7, \code{:submitted}, and 9:
the past-type result still acknowledges submission, not target completion.

Manager commands are past-type requests. Their handlers do not promise
application replies: do not replace \code{<=} with \code{<==} merely
to wait for installation. The bounded sessions check observable evaluator
identities and finite traces, not a general synchronization protocol.
Keep a recording reporter outside the watched sender's evaluator policy,
and bound or clear its retained events in longer experiments.

For printed output with filtering, the original \code{reporter-gen}
accepts two host Lisp functions. The predicate decides whether to print;
the formatter returns a string. For example:

\begin{codeblock}
\begin{lstlisting}
[filtered := [reporter-gen <== [:new
 :pred #'(lambda (type args context cont) (eq type :send-now))
 :format #'(lambda (type args context cont)
             (format nil "WATCH-NOW ~S" (second args)))]]]
[watch-manager <= [:watch-object watch-driver filtered]]
[watch-driver <== [:now :probe]]
(nowatch watch-driver)
\end{lstlisting}
\end{codeblock}

The formatter sees evaluated arguments, not source text. These functions
are host lambdas: do not place ABCL/R2 sends in their bodies or assume they
can capture arbitrary surrounding ABCL bindings. Use ordinary Lisp
parameters or explicit Lisp closure construction as in the earlier
lambda-boundary discussion.

\subsubsection{Choose the Scope and Understand Restoration}

\begin{longtable}{L{0.33\textwidth}L{0.57\textwidth}}
\textbf{Public operation} & \textbf{Scope of the change}\\\hline
\code{(watch obj)} / \code{(nowatch obj)} & Replaces/restores this object's
evaluator through its metaobject. Requires compatible reflective operations.\\
\code{(watch-group g)} / \code{(nowatch-group g)} & Changes the group's
evaluator for subsequent object creation; does not rewrite existing objects.\\
\code{(watch-toplevel)} / \code{(nowatch-toplevel)} & Changes/restores
\code{*eval*}, the interactive evaluator; not every object's evaluator.\\
\end{longtable}

The manager saves the previous evaluator before installing the wrapper.
For an object it sends \code{[:set-evaluator wrapper]} to the metaobject;
for a group it sends \code{[:replace-evaluator wrapper]} to the group
manager. The top-level operation assigns \code{*eval*}. These are three
different places to install the same kind of wrapper.

In a group experiment, create one sender before \code{watch-group} and
another after it. The first retains its original evaluator; the second
captures the wrapper. After \code{nowatch-group}, a third sender gets the
restored group policy, but the second still holds its wrapper. Group
restoration is not a broadcast repair of all existing objects.

There is also an original, preconfigured \code{watching} group. An object
defined with \code{(group watching)} uses its tracing evaluator from
creation. Alternatively, after loading \code{copy}, the library suggests
\code{(copy name :group watching)}. That rebinds \code{name} to a
\emph{different object}; it is not equivalent to \code{(watch name)}.
All the queue-sharing and local-state-copy restrictions of
Section~\ref{sec:copy} still apply. The checked copy example has no local
state entries and does not establish general copying compatibility.

Wrappers obtained from a group's creation policy are not entries in the
manager's watched-object list. Thus \code{(nowatch obj)} alone cannot
remove such a preconfigured wrapper. For a known single wrapper, an
explicit object-level restoration is:

\begin{codeblock}
\begin{lstlisting}
[wrapper := [[meta obj] <== :evaluator]]
[underlying := [wrapper <== :original-evaluator]]
[[meta obj] <= [:set-evaluator underlying]]
\end{lstlisting}
\end{codeblock}

Use this only when you know that the current evaluator implements this
watch-wrapper protocol. It removes one layer from that object, not the
group policy. Do not replace an unrelated custom evaluator blindly.

\paragraph{Avoid repeated installation.}
The original manager does not make \code{watch} idempotent. Calling it
twice adds two wrappers. On disabling the object, it restores the most
recent saved evaluator but deletes all saved entries for that object.
A second \code{nowatch} therefore leaves the inner wrapper installed.
The tests observe this boundary and explicitly restore the evaluator
saved before the experiment. Pair one installation with one restoration;
do not rely on repeated object, group, or top-level requests behaving as
a balanced instrumentation stack.

\paragraph{Evidence and exercise.}
Run \code{make -C tutorial watch-check}. Forty-one probes check
now/past/future results and event payloads, sender context, default and
filtered reporters, restoration, group creation policy, top-level control,
the preconfigured group, stateful watching without copying, an
empty-local-state watched copy, and repeated
installation boundaries. Libraries are copied unchanged; transcripts and
generated files stay in \filepath{build/watch-tests/}.

Why watch the driver rather than just the echo target? The driver issues
the application requests of interest. Why does turning off the group not
silence every existing sender? Each already has an evaluator reference.
As an exercise, record only \code{:send-past} events and explain which
are application requests and which are replies forwarded by the driver.


\subsection{Report: Observe Arrivals and Script Execution}
\label{sec:report}

Where \code{watch} observes outgoing evaluator send requests,
\filepath{lib/report.abclr2} observes the receiving object's dispatch
cycle. It supplies a replacement metaobject generator,
\code{meta/report-gen}, that reports arrivals, dispatch attempts, and
completed script bodies. The public \code{report} macro copies the
application object using this generator and rebinds its name to the copy.
This is not an evaluator wrapper installed on the original object.

Read Section~\ref{sec:copy} first. The copy retains references to scripts
and the outer lexical environment, but has a new object identity and queue.
The original Copy library's state-initializer limitation also applies:
do not assume that \code{report} can instrument arbitrary stateful
objects on this image. Our first client deliberately has no local state.

\subsubsection{Start with the Public Macros}

The supplemental \filepath{tutorial/examples/13-report.abclr2} explicitly
loads Copy before Report, whose original source requires Copy. Its
application and saved alias are:

\wholelisting[firstline=5,lastline=11]{tutorial/examples/13-report.abclr2}

The \code{:silent} handler intentionally sends no application reply.
Start with a request that does reply:

\begin{codeblock}
\begin{lstlisting}
(inc "tutorial/examples/13-report" :use-source t)
(report report-subject)
(eq report-subject report-original)       ; => NIL
[report-subject <== [:echo 7]]            ; => 7, plus reports
(noreport report-subject)
(eq report-subject report-original)       ; => T
\end{lstlisting}
\end{codeblock}

The macro takes a symbol naming an object, not an arbitrary expression.
Its three keyword options, \code{:message}, \code{:begin}, and
\code{:end}, all default to \code{t}. For example,
\code{(report report-subject :message nil :end nil)} selects only begin
reports. Disabling all three options suppresses printed reports but
still creates a copy; it does not make \code{report} a no-op.

Here is the original macro's construction sequence:

\wholelisting[firstline=30,lastline=37]{lib/report.abclr2}

The copier creates a new object with \code{meta/report-gen}; the macro
retrieves its default reporter, wraps it in a filter, saves the pair
\code{(new-object . previous-object)} in \code{*reporting-objects*},
installs the filter, and finally rebinds the supplied name. Other aliases
still denote the previous object. Requests through \code{report-original}
therefore bypass this new reporting metaobject.

\subsubsection{Follow One Dispatch Cycle}

Compare the original receive and begin handlers:

\wholelisting[firstline=60,lastline=79]{lib/report.abclr2}

The reporting protocol is
\code{[:report metaObject type [Message Reply Sender]]}.
Its identity field is the application's \emph{metaobject};
\code{[den metaObject]} gives the application object. The final triple
retains the received payload, application reply destination, and sender.
These are not the evaluator's completion continuation.

For a matching request, follow three distinct boundaries:

\begin{enumerate}
  \item \code{:message}: the metaobject receives a \code{:message}
        envelope and reports it before enqueuing the triple. If dormant,
        it becomes active and sends itself \code{:begin}.
  \item \code{:begin}: it dequeues a triple and reports it \emph{before}
        looking for a matching script. Bindings from a successful match
        extend the object's state environment for body evaluation.
  \item \code{:end}: the evaluator calls the supplied \code{cont} after
        the body finishes. That continuation sends an end report and the
        metaobject's internal \code{:end} request. The latter either
        starts another queued request or makes the object dormant.
\end{enumerate}

Thus a begin report means a dispatch attempt, not proof that a script
matched. An end report means that this body returned through evaluator
completion, not that all downstream work finished. A script may reply
early and continue working; another may finish without replying at all.
In our \code{:silent} case, a past-type send produces all three events
although there is no application answer. Do not call it with \code{<==}
and expect the end report to satisfy the waiting caller.

Reports themselves are past-type messages. The event order for one
completed request is arrival, begin, end, but reports from different
objects may interleave. They add work and cannot establish the timing
or scheduling of an uninstrumented run.

\subsubsection{Record the Actual Protocol Fields}

The teaching file includes a non-reifying recorder. This keeps the
diagnostic handler simple and avoids recursively instrumenting it:

\wholelisting[firstline=13,lastline=20]{tutorial/examples/13-report.abclr2}

In a fresh session after loading the file, replace the copy's reporter:

\begin{codeblock}
\begin{lstlisting}
(report report-subject :message nil :begin nil :end nil)
[[meta report-subject] <= [:set-reporter report-recorder]]
[report-subject <== [:echo 7]]
[events := [report-recorder <== :events]]
(mapcar #'second events)                 ; => (:MESSAGE :BEGIN :END)
(mapcar #'(lambda (e) (first (third e))) events)
; => ((:ECHO 7) (:ECHO 7) (:ECHO 7))
\end{lstlisting}
\end{codeblock}

Each record is \code{(metaObject type triple)}. The lambda only transforms
its ordinary Lisp argument; it neither captures ABCL variables nor sends
messages. Pass context explicitly if a host helper needs it.

\code{:set-reporter} replaces the \emph{whole} reporter chain. The macro's
earlier keyword filter is no longer in that path, so the recorder receives
all three event types. To select only end events while retaining structured
records, install a new filter:

\begin{codeblock}
\begin{lstlisting}
[only-end := [standard-filter-gen <== [:new report-recorder
                         :message nil :begin nil :end t]]]
[[meta report-subject] <= [:set-reporter only-end]]
[report-recorder <== :reset]
[report-subject <== [:echo 8]]
(mapcar #'second [report-recorder <== :events])
; => (:END)
(noreport report-subject)
\end{lstlisting}
\end{codeblock}

The generator's message requires all three named fields in this order;
it is an object protocol, not a Common Lisp keyword-argument function.
The filter forwards enabled events unchanged and drops disabled ones.
It does not suppress application messages or change script execution.
Reporter installation uses \code{<=}: the setter promises no application
reply, so changing it to \code{<==} would introduce a wait with no answer.

\paragraph{Reading the default report.}
Prefer structured records when interpreting reply and sender identities.
The \code{default-reporter} receives \code{[M R S]} and passes those
fields to \code{default-reporting-function} in the same order.
Its \code{from} field prints the sender; \code{reply=} prints the
application reply destination, not the evaluation continuation.
Either field can be \code{nil}. Printed objects and messages are abbreviated by
\code{*print-length*} and \code{*print-level*}. Do not use the displayed
names or labels as a substitute for the recorded triple.

\subsubsection{Restoration and the Limits of Instrumentation}

\code{noreport} looks up the current object in \code{*reporting-objects*},
removes its association, and rebinds the supplied name to the saved object.
It does not destroy the reporting copy, remove other aliases, drain its
queue, or merge changes back into the previous object. A retained alias
to the copy still runs and reports after \code{noreport}.

Repeated \code{report} creates successive copies, not multiple reporting
layers around the same object. Each copy has its own reporting metaobject;
it does not forward dispatch through the previous copy's metaobject.
With an intact registry, successive \code{noreport} calls restore the
previous copy and then the original. This differs from Watch's repeated
installation limit (Section~\ref{sec:watch}). Do not confuse this name
history with a general facility for rolling back application effects.

For an active application, copying is not an atomic snapshot. Existing
references still target the old object. Pending queue entries preserve
their original reply and sender references; Copy installs them with
\code{:enqueue}, not \code{:message}, so they do not acquire new arrival
reports merely by being copied. They also do not start simply because
the copy was returned. Follow the explicit \code{:begin} rules in
Section~\ref{sec:copy}, and do not manually start an already active
metaobject. This book uses an idle, empty-state subject to keep those
separate issues out of the first experiment.

\paragraph{Check your understanding.}
Will \code{report-original} start reporting when the other name is rebound?
\emph{Answer:} No; it retains the old object identity. Does an end report
replace \code{!} for the application caller? \emph{Answer:} No; evaluator
completion and application replies have separate destinations. Does
\code{noreport} stop a reporting copy held under another alias?
\emph{Answer:} No; it restores one binding, not every reference.

\paragraph{What the checks establish.}
Fifty-nine bounded probes check identities, alias preservation, event
triples, all eight filter combinations, repeated copying/restoration,
completion without an application reply, and dispatch state.
The formatter checks cover arguments and output for all three event types
and \code{nil} reply/sender fields. Run \code{make -C tutorial report-check}; copies,
inputs and transcripts remain in \filepath{build/report-tests/}.
These are bounded observations, not certification of arbitrary stateful
copying or transparent instrumentation.


\subsection{Retry: Repairing a Message Acceptance Condition}
\label{sec:retry}

Watch and Report observed execution. This project changes dispatch when
a request cannot yet be accepted. The question mark in the opening comment of
\filepath{examples/misc/fault-tolerant.abclr2} is worth keeping in mind.
This is a small experiment in reflective recovery, not a distributed
fault-tolerance service. Its receiver refuses ordinary requests until it
has received an initialization message. The sender's metaobject repairs
that condition and retransmits the refused request.
No machine crash or lost network packet is involved.

\subsubsection{Start with the Application}

Read the two application objects before their generator:

\wholelisting[firstline=130,lastline=147]{examples/misc/fault-tolerant.abclr2}

The server owns \code{amount}, initially \code{nil}. Its initialization
handler resets that value to zero, and \code{[:count-up n]} adds an amount
and prints the resulting count. The client has no application state:
\code{[:start server]} simply sends \code{[:count-up 10]} asynchronously.
Both explicitly select \code{FT-meta-gen}, so both participate in the
recovery protocol. The client's empty initialization handler is still
important: its metaobject must accept the client before it can run
\code{:start}.

In a fresh session at the repository root, the file's own entry sequence is:

\begin{codeblock}
\begin{lstlisting}
(inc "examples/misc/fault-tolerant" :use-source t)
[client <= [:initialize]]
[client <= [:start server]]
; Prints Count is 10.
\end{lstlisting}
\end{codeblock}

The server was deliberately not initialized by these commands. A successful
count therefore depends on recovery. Do not change \code{amount}'s initializer
to zero: that would conceal the problem the original example demonstrates.
For scratch-copy checks that leave the original directory untouched, run
\code{make -C tutorial retry-check}.

\subsubsection{Trace the Rejection and Retransmission}

Follow one request in six steps. This is a causal explanation, not a promise
about the order of unrelated console output.

\begin{enumerate}
\item The client's metaobject accepts \code{[:initialize]}, sets its private
acceptance flag, and executes the empty client handler.
\item The accepted \code{[:start server]} handler sends
\code{[:count-up 10]} to the server. Its sender is the client.
\item The server's metaobject dequeues the request. Its acceptance flag is
false, so it does not execute the arithmetic on \code{nil}.
\item It sends \code{[:reject Receiver Message Reply Sender]} to the
client's metaobject. Here \code{Receiver} is the server and
\code{Message} is the original \code{[:count-up 10]}.
\item The client's metaobject sends \code{[:initialize]} to the server,
then forwards the original message envelope to the server's metaobject.
\item The server processes initialization before the retransmitted request.
Initialization sets \code{amount} to zero; the accepted count request
then sets it to 10 and prints the count.
\end{enumerate}

The original recovery handler is short:

\wholelisting[firstline=47,lastline=51]{examples/misc/fault-tolerant.abclr2}

Notice the two levels of addressing. The initialization goes to the
application object \code{Receiver}. Retransmission goes directly to
\code{[meta Receiver]}, using the interpreter's \code{:message} envelope.
This preserves the old \code{Reply} and \code{Sender} fields rather than
creating a new application request with the recovering metaobject as sender.
In this trace the recovered request was a past-type send, so it has no
waiting application reply to deliver.

\subsubsection{Read the Acceptance State Machine}

The generator constructs the familiar queue, environment, script set,
evaluator, group manager, and \code{:dormant}/\code{:active} dispatch mode.
Its extra \code{initialized} flag belongs to the \emph{metaobject};
it is not the server's application variable \code{amount}.

\wholelisting[firstline=53,lastline=60]{examples/misc/fault-tolerant.abclr2}

The \code{:message} handler enqueues a three-field application envelope.
If dormant, the metaobject activates itself and schedules \code{:begin}.
That handler removes one envelope and applies the acceptance condition
before looking for an application script:

\wholelisting[firstline=66,lastline=80]{examples/misc/fault-tolerant.abclr2}

\code{[:initialize . \_]} matches a list beginning with the keyword
\code{:initialize}; the wildcard ignores its remaining tail. The flag is
set \emph{before} the application handler is found or executed. Thus the
comment saying that initialization has executed is stronger than the code:
the flag establishes that an initialization-shaped message was encountered.
It does not validate a successful initialization result.

For rejection, \code{Me} is the server's metaobject and \code{[den Me]}
recovers the application server it represents. \code{[meta Sender]}
selects the sender's recovery machinery. The rejected request has already
been removed from the queue; \code{:end} advances dispatch rather than
putting it back locally. The sender is responsible for retransmission.

For acceptance, lines 83--90 use \code{find-script} to obtain bindings and
a script body. The evaluator executes that body in an environment extending
the application's state. Its completion continuation schedules \code{:end}.
This continuation is an interpreter completion notice, not a client/server
acknowledgment. Lines 94--97 either process the next queued request or return
to dormant mode. The generator's final \code{:do-evlis} and \code{pairlis}
sequence builds the initial application environment, as in the earlier
inheritance and MCP generators; that construction is separate from the
application's \code{[:initialize]} protocol.

\subsubsection{What the Experiment Does and Does Not Guarantee}

A second client start prints \emph{Count is 20.}: the server is already
accepted, so no recovery reset is needed. An explicit server initialization
resets the count to zero; the next client start prints \emph{Count is 10.}
Six isolated checks verify these state transitions, the printed counts,
and an empty final server queue using the distributed example.
Reflective state queries inspect \code{amount}; they are test probes, not
new server API methods.

The example has no retry budget, timeout, persistent log, duplicate detection,
or exactly-once delivery mechanism. A sender without the expected
\code{:reject} handler cannot supply this recovery. Multiple rejections can
cause multiple initialization messages, each resetting server state.
The checked two-object sequence demonstrates one repaired precondition,
not correctness for every interleaving or failure.

\paragraph{Check your understanding, with answers.}
Why initialize the client even though its handler does nothing?
Its metaobject also gates application requests; the empty handler satisfies
the chosen protocol without adding application work.
Why preserve the original sender on retry?
Script matching can depend on the sender, and subsequent rejection must
reach the original sender's recovery machinery.
Would automatically resetting the server after every request be equivalent?
No: it would lose the accumulated count. Recovery belongs to the failed
acceptance path, not the normal successful path.


\subsection{Breakpoints: Inspecting a Suspended Lexical Environment}
\label{sec:break}

Retry repaired a precondition and resent a request; here the application
deliberately waits while we inspect its current bindings.
\filepath{examples/misc/break.abclr2} implements an interactive debugger
in ABCL/R2 itself. It does not invoke the Common Lisp debugger or install
an operating-system interrupt handler. An application calls a macro,
waits for a reply, and lets another object inspect its lexical environment
before execution continues. This makes reflection useful for a concrete
task rather than merely printing an object's representation.

\subsubsection{Reproduce the Original Session}

Load the debugger before defining code that uses its macro. At the repository
root, enter the following. The definition of \code{foo} comes from the
original file's commented session, not from a replacement example.

\begin{codeblock}
\begin{lstlisting}
(inc "examples/misc/break" :use-source t)
[object foo
 (state [x := 0])
 (script
  (=> y
      (temporary [z := (+ x y)])
      (abclr2-break)
      !z))]
[foo <== 123]
\end{lstlisting}
\end{codeblock}

The handler binds \code{y} to 123 and computes the temporary \code{z}
as \code{x + y}, initially 123. The breakpoint now prints an object
identifier and an \code{Inspect(...)} prompt. Object numbers and generated
reply-variable names vary between sessions; they are not expected results.
At this prompt you are talking to the environment inspector, not the
ordinary \code{ABCL/R2>} reader.

Use the following sequence. For \code{:set}, the next two inputs answer
the \code{Name =} and \code{Value =} prompts.

\begin{codeblock}
\begin{lstlisting}
:l             ; lists Z in the temporary frame
z              ; prints Z=123
:u
:l             ; lists Y and a generated reply binding
y              ; prints Y=123
:u
:l             ; lists X in the application-state frame
:set
x
10
:d
:d
:set
z
999
:r
0
\end{lstlisting}
\end{codeblock}

The first \code{:u} reaches the handler's pattern bindings, and the second
reaches the object's state. The two \code{:d} commands retrace those steps.
The resumed handler sends 999, because you changed its current \code{z}.
The value 0 supplied after \code{:r} is the result of
\code{(abclr2-break)}, which this handler ignores before executing
\code{!z}.

Now send \code{[foo <== 123]} again. At the new inspector prompt, enter
\code{:r}, then \code{t}. The handler returns 133: \code{x} remained 10,
but a new invocation computed a new temporary \code{z = 10 + 123}.
This is the original demonstration's distinction between persistent state
and an invocation-local binding. The inspector changes existing bindings;
it does not automatically recompute an already calculated \code{z}
when you change \code{x}.

\subsubsection{The Inspector's Small Command Language}

\begin{longtable}{L{0.25\textwidth}L{0.65\textwidth}}
\textbf{Input} & \textbf{Effect}\\\hline
\endhead
\code{:l}, \code{:list} & List the names in the selected local frame.\\
\code{:a}, \code{:alist} & Print that frame's association list of names and values.\\
A variable name & Look up its value, searching outward if necessary;
print an unbound notice if lookup fails.\\
\code{:u}, \code{:up} & Select the outer lexical environment, if one exists,
and remember the current frame.\\
\code{:d}, \code{:down} & Return to the most recently remembered frame,
if there is one.\\
\code{:s}, \code{:set}, \code{:=} & Read a name and a value, then update
or create that binding in the selected local frame.\\
\code{:h}, \code{:help}, \code{?} & Print the built-in command help.\\
\code{:r} & Read a return value, reply to the breakpoint caller, and finish
the inspector handler.\\
\end{longtable}

These commands are read by Lisp's \code{read}, but are interpreted by
the inspector's \code{case}, not evaluated as arbitrary Lisp forms.
Enter \code{x}, not \code{'x}, when asked for a name. Entering
\code{(+ 1 2)} as a value supplies a list; it does not calculate 3.
In particular, \code{:r} is a command followed by a separate value input.
It is missing from the original printed help, although the loop's
termination test implements it.

Local listing and outward lookup are different operations. You can look up
\code{x} while positioned at the temporary frame even though \code{:l}
there lists only \code{z}. However, \code{:set} uses \code{:set-local},
not the outward-searching \code{:set} environment operation. Setting
\code{x} in the temporary frame would introduce a local binding there,
not update the object's existing state. Navigate to the state frame first
when you intend the persistent change shown above.

\subsubsection{How the Macro Obtains the Environment}

Read the original monitor:

\wholelisting[firstline=7,lastline=16]{examples/misc/break.abclr2}

The macro at lines 65--66 expands into a now-type request to this monitor.
It creates an anonymous, empty-script \emph{dummy object} at the call site
and includes that object, an optional format string, and a list of format
arguments in the message. The macro uses backquote, comma, and
comma-at as introduced in the Lisp reading guide: it builds code while
inserting the caller's forms.

Why create an object that handles no messages? Its state environment
retains the lexical environment in which it was created.
\code{[meta dummy-object]} obtains its metaobject;
the \code{:state} query obtains its state environment; and
\code{:lexenv} follows that environment's outer link. That link is the
application environment at the breakpoint. It is a live environment object,
not a copied map or a snapshot of a machine stack.

The monitor receives the breakpoint request's reply destination as
\code{R} and its application sender as \code{sender}. It prints that sender,
then sends \code{[environment nil]} to \code{environment-inspector},
forwarding \code{@ R}. The initial \code{nil} is an empty navigation
history. The monitor does not manufacture the final return value or wait
for a second result: the inspector replies directly to the waiting caller.
If a format string is present, \code{~?} applies it to the supplied
argument list to print a contextual message.

\subsubsection{Navigation, Mutation, and Resume}

The inspector explicitly selects \code{non-reifying-meta}. Its handler is
compiled for the non-reifying execution path rather than interpreted by a
custom reflective evaluator. It can still send messages to the environment
objects it was given. Non-reifying does not mean that the supplied
environment has become inaccessible.

\wholelisting[firstline=21,lastline=29]{examples/misc/break.abclr2}

The loop begins with the command \code{:skip}, whose case does nothing.
Its step expression prints the prompt and reads the next command.
The termination test recognizes \code{:r}; it reads one more datum and
sends it to \code{R}. Until this happens, the application's now-type
breakpoint request has not received its reply. The suspended application
handler resumes after that reply, not after a local-listing or edit command.

\wholelisting[firstline=31,lastline=49]{examples/misc/break.abclr2}

For upward navigation, \code{env-record} is a stack of previously selected
environments, represented as a Lisp list. The bracket expression
\code{[environment . env-record]} constructs a new head/tail pair.
Downward navigation retrieves its \code{car} and removes it with
\code{cdr}. It therefore retraces your own upward moves; it does not
search for every inner scope in the program. An upward move with no outer
environment or a downward move with empty history does nothing.

The setting branch uses \code{let*} so the name is read before the value,
then sends the local-binding update asynchronously. The original session
and isolated checks show those updates visible when the handler resumes.
Do not extend that observation to arbitrary debugger clients or scheduler
policies without checking their message ordering.

\paragraph{Evidence and scope.}
Run \code{make -C tutorial break-check}. Six checks replay the original
definition and editing session from unchanged source copies, verify
999 followed by 133, inspect persistent \code{x}, and confirm that
a supplementary handler receives \code{(+ 1 2)} as raw breakpoint-return
data. Transcript assertions also check local listings, variable values,
unbound-name handling, help, and an optional formatted message.
They do not certify every command alias or concurrent breakpoint arrival.
The single inspector reads shared interactive input; this example is not
a multi-user debugger, a rollback facility, or a whole-system pause.

\paragraph{Check your understanding, with answers.}
Why does the second call not return 999?
The first call's \code{z} was temporary; the second creates a fresh binding.
Why can its result still depend on the first session?
The state binding \code{x} was changed in the persistent state frame.
Does entering \code{:r} and then 0 force the application to return 0?
No. It returns 0 from the breakpoint expression; the handler still decides
what application reply to send afterwards.


\subsection{The Original Express Clocks}
\label{sec:original-express}

The breakpoint could wait for an inspector's reply. A running clock poses
a different problem: how can a new request affect a handler that has
not finished? An ordinary object finishes one script before starting the next queued
request. This creates a problem for a clock whose \code{:start} script
loops until another request sets \code{stop}: an ordinary stop request
waits behind the very script it must stop. Express messages provide an
alternative dispatch policy, not faster transport or operating-system
threads. Read \filepath{examples/misc/express-ex.abclr2} with
\filepath{lib/express.abclr2}; its older companion uses
\filepath{lib/express-old.abclr2}.

\paragraph{Configuration and the application first.}
The original group installs both halves of the protocol:

\wholelisting[firstline=9,lastline=12]{examples/misc/express-ex.abclr2}

The evaluator recognizes express sends; the metaobject generator equips
receivers to accept them. Both \code{clock} and \code{test} belong to this
group. A keyword in a send does not by itself change the primary evaluator
or the receiver's ordinary queue policy. The top-level only sends an
ordinary request to \code{test}; express sends occur inside that object.

The clock starts with \code{time = 0} and \code{stop = nil}.
Its \code{:start} script resets only \code{stop}, not \code{time}:

\wholelisting[firstline=19,lastline=31]{examples/misc/express-ex.abclr2}

In Common Lisp, \code{(do () (stop) ...)} has no loop variables and tests
\code{stop} before each iteration. \code{(dotimes (i 5))} runs five empty
iterations; it is neither a delay in seconds nor a clock measurement.
Each completed pair adds two. The intermediate value between increments
is odd, so excluding express observation of that intermediate state is
the point of \code{atomic}. Without it, an odd sample is permitted, not
guaranteed: a particular schedule may still produce only even samples.

\wholelisting[firstline=44,lastline=50]{examples/misc/express-ex.abclr2}

The comments calling these scripts ``express'' do not declare a special
kind of script. The \emph{send} chooses the mode; an ordinary \code{:peek}
uses the same matching script. \code{:peek} replies with the current count.
\code{:stop} changes a binding but does not reply. Stopping is completed
when the resumed loop notices that binding and replies \code{:finished}.

\par\noindent\begin{minipage}{\linewidth}
\paragraph{Following the controller's two phases.}
The controller supplies the orchestration rather than the clock itself:
\par\vspace*{\medskipamount}

\wholelisting[firstline=55,lastline=71]{examples/misc/express-ex.abclr2}
\end{minipage}\par

Trace the original \code{[test <== :start]} as follows:

\begin{enumerate}
  \item \code{R} is the caller's reply destination. The controller sends
  \code{clock} an ordinary asynchronous \code{:start}, attaching a
  continuation instead of waiting for the loop to finish.
  \item Twenty now-type express \code{:peek} requests each wait for one
  count and print it. They are sequential observations by this controller,
  not twenty simultaneous requests.
  \item A past-type express \code{:stop} requests the state change without
  waiting for a reply that the stop script never sends.
  \item The clock's \code{:finished} reply invokes the continuation.
  Its \code{\_} discards that value and sends \code{Me} a request for the
  second phase, explicitly retaining \code{R}.
  \item The second phase repeats the experiment without the paired
  increments' atomic wrapper. The count continues from the first phase.
  Its completion continuation sends \code{:end} to the controller.
  \item Only \code{:end}'s \code{!:finished} answers the original caller.
\end{enumerate}

Thus the immediate controller script finishing, an express peek replying,
and the entire two-phase test completing are three different events.
Replacing the asynchronous start by a now-type start would prevent this
controller from issuing its peeks and stop. Replacing stop by a now-type
stop would wait for a nonexistent application reply.

\paragraph{How the newer library interrupts evaluation.}
The generator creates a dispatcher, a metaobject, and an evaluator front
end for each receiver. The dispatcher finds the script and constructs
its binding environment over the object's persistent state. The front
end forwards ordinary evaluation to the group's evaluator, but can
retain the next \code{[:do . Context]} request instead of executing it.
This is cooperative interruption at an evaluator boundary: it cannot
stop a host Lisp operation halfway through its machine instructions.

\begin{longtable}{L{0.18\textwidth}L{0.72\textwidth}}
\textbf{Mode} & \textbf{Intended handling of an arriving express request}\\\hline
\endhead
\code{:dormant} & Dispatch it directly; there is no normal script to suspend.\\
\code{:active} & Ask the front end to \code{:grab}; save the normal
evaluation request and its completion destination, then dispatch the express script.\\
\code{:express} & Queue it; one express script does not preempt another.\\
\code{:lock} & Queue it until the atomic region releases the lock.\\
\end{longtable}

Normal arrivals during active, express, or locked execution enter the
normal queue. At an express script's completion, the intended order is
pending express requests, resumption of the saved normal context, and
then queued normal requests. The saved context includes the destination
\code{R} of the evaluator request, not necessarily the application
caller's reply destination. Resuming it must preserve evaluation's
completion path as well as expression and environment.

The \code{atomic} macro in \filepath{src/user.lsp} asks \code{[meta Me]}
for a now-type \code{:lock}, evaluates its body, sends a past-type
\code{:release}, and returns the body's last value. In the newer library,
locking saves the previous mode and sets \code{:lock}; releasing restores
that mode and can start an express request deferred during normal
execution. This is object-local protection against this express protocol,
not a transaction over all objects. The macro uses no unwind-protection:
do not infer automatic release after an error or a non-local exit.

\paragraph{Why the older peek locks too.}
The July version has no grabbing front end or suspended normal context.
It dispatches express scripts without putting them into the normal
message queue, while normal evaluation can continue. Its
\code{express-count} tracks outstanding express executions. A separate
lock object stores waiting reply destinations and grants them in order.
Both a writer and an observer must acquire that same receiver-local lock:

\wholelisting[firstline=41,lastline=46]{examples/misc/express-old-ex.abclr2}

An unlocked read in that version could run between the two protected
increments, because the reader does not participate in the lock protocol.
The newer example deliberately removes \code{atomic} from \code{:peek}:
under its intended suspension policy, the normal writer cannot advance
while the express reader executes, and a reader arriving during a locked
pair must wait. These are different algorithms, not interchangeable files.

Both libraries define the same public generator names and provide the
feature \code{express}. Compare them in separate fresh sessions, explicitly
loading the matching library first. The old example's \code{require}
otherwise names \code{express}, not \code{express-old}; its commented
sample also still names \code{express-ex}. Use the actual old filename:

\begin{codeblock}
\begin{lstlisting}
; New version, in a fresh session:
(inc-lib "lib/express" :use-source t)
(inc "examples/misc/express-ex" :use-source t)

; Old version, in a different fresh session:
(inc-lib "lib/express-old" :use-source t)
(inc "examples/misc/express-old-ex" :use-source t)
\end{lstlisting}
\end{codeblock}

\paragraph{Connect the reading to the first experiment.}
The idle-clock exchange introduced earlier isolates sender configuration,
reply forwarding, and the receiver's express dispatch.
Read \code{:start-non-atomic} next: it starts normal clock work, samples
the value, sends stop, and waits for the clock's completion reply.
Follow the roles of the two reply destinations separately before tracing
the atomic phase.

The queue-handling branches dequeue the selected envelope and bind
\code{[Message Reply Sender]} before dispatch. Compare this with the
saved execution context: an envelope describes an application request,
whereas the saved context describes suspended expression evaluation.
The \code{express-check} sessions provide short examples of both
library versions, queue order, and reply routing.

\paragraph{Check your understanding.}
Why is \code{:stop} sent asynchronously? Why is \code{time} not reset
between phases? Could the non-atomic phase print only even values? Why
does configuring only the clock fail to establish express sending?
Answers: stop has no reply; only \code{stop} is reset; scheduling need
not expose an odd intermediate state; and the sender's evaluator must
translate the extra send argument into the express meta-level protocol.


\subsection{Pcontrol: Scheduling a Recursive Object Computation}
\label{sec:pcontrol}

The programs in \filepath{examples/pcontrol/} move reflection from an
individual object's behavior to a group's execution policy. A recursive
algorithm can create many objects before earlier work has finished.
The example asks whether choosing small pending computations first can
reduce this growth. Here, \emph{parallelism control} means controlling the
progress of concurrent work, not creating operating-system threads.
Read the application first, then follow the machinery that schedules it.

\subsubsection{A Checked Session and the Loading Order}
\label{sec:pcontrol-load}

Start a fresh interpreter with its working directory set to a scratch
copy of \filepath{examples/pcontrol/}. Relative names in
\filepath{start.abclr2} refer to that directory. Loading creates translations
and compiling creates compiled files; the experiments also overwrite
their own log files. The automated reading check creates this scratch
copy for you:

\begin{lstlisting}
make -C tutorial pcontrol-check
\end{lstlisting}

For an interactive session in such a copy, use:

\begin{codeblock}
\begin{lstlisting}
(inc "start" :use-source t)
(inc "fib" :use-source t)
[test-pc <== [:start 4]]
; Prints Fibonacci(4) = (:ANS 3); replies T.
[test-np <== [:start 4]]
; Prints Fibonacci(4) = (:ANS 3); replies T.
\end{lstlisting}
\end{codeblock}

The printed \code{(:ANS 3)} is the Fibonacci answer message, whereas
\code{T} is the test object's final reply. These are different protocol
results. Use the argument shapes in \filepath{README.abclr2} and the
actual handlers: some commands in the commented block at the end of
\filepath{start.abclr2} have obsolete extra arguments.
Keep inputs small; the original README warns that these experiments
are expensive.

\filepath{start.abclr2} compiles both applications but loads only the
four machinery files. Their order supplies the dependencies:

\begin{enumerate}[nosep]
  \item \filepath{para-monitor.abclr2} defines the monitor.
  \item \filepath{pqueue-eval.abclr2} defines the queue and evaluator factories.
  \item \filepath{priori-meta.abclr2} defines application metaobjects.
  \item \filepath{pcontrol-group.abclr2} connects them into a group.
\end{enumerate}

The subsequent \code{inc "fib"} or \code{inc "qsort"} chooses the
application. Both application files define \code{test-pc} and
\code{test-np}; loading the second replaces those bindings. Use a fresh
session when changing applications.

\subsubsection{Fibonacci: Create Children, Then Collect Replies}
\label{sec:pcontrol-fib}

The convention is \code{F(1) = F(2) = 1}. The source does not provide a base
case for zero or negative inputs. A factory \code{fib-gen-gen} creates
a \code{fib-gen} in the \code{pcontrol} group. That generator creates
one \code{fib} object per recursive subproblem, with its problem size
as its priority. The child objects inherit the creation context's group.

\wholelisting[firstline=18,lastline=33]{examples/pcontrol/fib.abclr2}

The literal patterns \code{1} and \code{2} precede the general
variable pattern \code{n}. A base case sends \code{[:ans 1]} to its
request's reply destination, then reports logical completion to the
monitor. A larger problem saves its caller's destination in state
\code{reply}. Each nested \code{fib-gen <== [:new ...]} waits for
the newly created child object, not for its Fibonacci result.
The following \code{<=} starts that child asynchronously with
\code{@ Me}, so the child sends its answer back to the parent.

The parent does not block while collecting answers. Its first
\code{[:ans n]} stores the value in \code{sub-node-value}; its
second adds both values and forwards \code{[:ans sum]} to the
saved destination. Because all valid results are positive, \code{nil}
can safely mean no first answer yet. This is a one-computation-per-object
protocol: the state is not reset for a second request.

Trace the five application nodes for input 4:

\begin{enumerate}
  \item The root for 4 saves its caller and creates children for 3 and 2.
  \item The node for 3 saves the root's identity and creates children
        for 2 and 1.
  \item Each base node replies \code{[:ans 1]} to its parent.
  \item The node for 3 combines two answers into \code{[:ans 2]}.
  \item The root combines 2 and 1 into \code{[:ans 3]}.
\end{enumerate}

This is a dependency trace, not a guaranteed temporal ordering.
At each join either child may reply first. Addition makes the
aggregation independent of that order. The test waits on the root,
prints its answer, requests that monitoring stop, and replies \code{T}.
Its \code{<==} is therefore a completion boundary for the computation,
not just for dispatching the first child.

\subsubsection{Where the Scheduling Policy Actually Lives}
\label{sec:pcontrol-policy}

The group definition is short enough to read in full:

\wholelisting[firstline=9,lastline=16]{examples/pcontrol/pcontrol-group.abclr2}

The group uses \code{priori-meta-gen} for each application object's
metaobject. Its evaluator is a \code{multi-eval} created with four
primary evaluators. The group forwards an object-creation request,
including \code{:priority priority}, to its metaobject generator.
Four evaluators here are four runtime objects; this declaration does
not establish execution on four processors.

There are two queues to distinguish. Each application metaobject in
\filepath{priori-meta.abclr2} has its own ordinary FIFO message queue.
On \code{:begin}, it selects a script and submits a \code{[:do ...]}
request to the shared evaluator, adding the object's priority after
the usual context fields. A continuation sends \code{:end} back to
that metaobject when script evaluation completes. That signal advances
the object's queue; it is separate from the script's application reply.

The shared \code{multi-eval} has a metaobject made by \code{PQ-meta-gen}.
That metaobject queues evaluation requests as
\code{[Priority [Message Reply Sender]]}.
An explicitly supplied priority is associated with the requesting
object's \code{Id}; requests without an explicit priority use the
recorded association, or zero if none exists.
The priority is thus scheduling information in the evaluation context,
not an extra field in \code{[:ans value]}.

Both tests keep this machinery and the \code{pcontrol} group.
They replace the shared evaluator metaobject's queue:

\wholelisting[firstline=44,lastline=45]{examples/pcontrol/fib.abclr2}

The \code{pc} variant installs a queue sorted by \code{<} on each
entry's \code{first} element. The \code{np} variant instead uses
\code{[queue-gen <== :new]}, the ordinary FIFO queue.
In the nested expression, \code{[pcontrol <== :evaluator]} obtains
the evaluator object; \code{meta} obtains its metaobject; and
\code{[:set-queue ...]} changes that metaobject's pending-work policy.
This does not change each Fibonacci object's inbox into a priority queue.

\wholelisting[firstline=12,lastline=18]{examples/pcontrol/pqueue-eval.abclr2}

\code{cons} inserts an entry into the list and \code{sort} reorders
that list. \code{:key \#'first} tells sorting to compare the priority
field. Smaller problem sizes are selected before larger ones among
entries pending in this queue. Do not infer a stable order for equal
priorities: the source uses \code{sort}, not a stable tie-breaking
policy. Nor is this express interruption: it selects pending requests
rather than preempting an already running script.

Finally, read \code{multi-eval}'s two handlers near the end of
\filepath{pqueue-eval.abclr2}. One recognizes an object-definition
representation carrying \code{:priority} and forwards creation to
the chosen group manager, then arranges the resulting name binding
and reply. The catch-all handler forwards other evaluation requests
to the next primary evaluator in a cyclic list, preserving their
completion destination \code{C}. This explains how the example
changes creation and scheduling without rewriting the Fibonacci
application's reply protocol.

\subsubsection{Quicksort: A Message-driven State Machine}
\label{sec:pcontrol-qsort}

\filepath{qsort.abclr2} uses the same group but a different application.
Its opening comment explains why scanning is implemented through
self-messages rather than a \code{do} loop. Read those self-messages
as transitions between phases, not as recursive host function calls.

Each sorting object captures \code{array}, \code{left}, and
\code{right} from its generator's lexical scope. Its priority is
\code{(- right left)}. For a nontrivial range, \code{:start}
saves its caller's destination, chooses the middle element as
\code{pivot}, initializes indices \code{i} and \code{j}, and sends
\code{:scan-left} to itself.

\begin{enumerate}
  \item \code{:scan-left} advances \code{i} until its element is
        not smaller than the pivot, then requests \code{:scan-right}.
  \item \code{:scan-right} moves \code{j} left until its element
        is not larger than the pivot, then requests \code{:test-swap}.
  \item \code{:test-swap} exchanges elements if \code{i <= j},
        advances both indices, and either resumes scanning or creates
        two child sorting objects for the remaining ranges.
  \item Each child sends \code{:fin} to its parent. The parent's
        \code{wait} flag records the first completion; the second
        forwards \code{:fin} to the parent's saved caller.
\end{enumerate}

\code{aref} reads an array element; \code{setf} changes that element
in place. The descendants share this array but work on partitioned
ranges. The first completion means only that one side is sorted;
the second establishes completion of both. A range with
\code{left >= right}, including an empty range, immediately replies
\code{:fin}.

\wholelisting[firstline=54,lastline=62]{examples/pcontrol/qsort.abclr2}

For \code{(3 1 3 2)}, the root covers indices 0 through 3 and
chooses the element at index 1, namely 1, as its pivot.
The left scan stops at index 0. The right scan moves from 3 to 1.
The swap produces \code{(1 3 3 2)}, with \code{i = 1} and
\code{j = 0}. Scanning ends, so the children cover the already
trivial range 0--0 and the remaining range 1--3. Their completions
join at the root; the final array is \code{(1 2 3 3)}.

In a fresh session with the same setup, a deterministic test avoids
the random initialization object:

\begin{codeblock}
\begin{lstlisting}
(inc "start" :use-source t)
(inc "qsort" :use-source t)
[input-array := (make-array 4 :initial-contents '(3 1 3 2))]
(coerce [test-pc <== [:start input-array 4]] 'list)
; => (1 2 3 3)
(coerce [test-np <== [:start input-array 4]] 'list)
; => (1 2 3 3)
(coerce input-array 'list)
; => (3 1 3 2)
\end{lstlisting}
\end{codeblock}

\code{coerce ... 'list} makes an array's contents convenient to inspect.
Both test handlers first copy their input with
\code{make-array :initial-contents}; they return the sorted copy,
not the caller's modified array. The original random entry point
\code{[init <= N]} instead stores an array in the outer binding
\code{a}; follow its \code{sample-gen} self-message loop to see
how it fills each index.

\subsubsection{Reading the Monitor Without Overinterpreting It}
\label{sec:pcontrol-monitor}

\filepath{para-monitor.abclr2} records two columns: a logical
\code{clock} and a counter. The evaluator sends \code{:tick}
when handling an evaluation request; this clock is not elapsed
wall-clock time. Application metaobjects announce creation and
transitions between dormant and active dispatch. The application
explicitly sends \code{:object-dead} at logical completion.
Nothing in that message destroys a Lisp object or measures garbage
collection.

The active count tracks announced dispatch activity, not the number
of physical CPUs busy. A node waiting for child answers can be
dormant yet still represent unfinished application work.
The object count combines creation announcements with explicit
logical-death announcements. For example, generator objects are
also created through the machinery. It is not simply a count of
the recursive leaves, allocated bytes, or all reachable host objects.
The \code{:record} messages are asynchronous observations of these
counters.

For Fibonacci input 4, the actual source names the files
\code{fib-active-pc-4.log}, \code{fib-object-pc-4.log}, and
the corresponding \code{np} files; quicksort uses a \code{qsort-}
prefix. Repeating the same run supersedes those files.
The directory README's generic names omit these application prefixes.
Compare traces as experiment output, not as proof that one policy
always reduces memory or is faster.

\paragraph{Check your understanding.}
Does \code{test-np} remove priority-bearing metaobjects?
\emph{Answer:} No. It replaces the shared evaluator's work queue with
a FIFO queue while retaining the group and application metaobjects.
Why can a parent accept either child's answer first?
\emph{Answer:} Its state records the first answer and combines it
with the second; Fibonacci uses addition, while quicksort counts
two completion signals for disjoint subranges.

\paragraph{Reproducing the reading checks.}
\code{make -C tutorial pcontrol-check} runs twelve checks in fresh
sessions using copies of the distributed sources. They check priority-queue
ordering for distinct priorities, both Fibonacci policies at 4 and 6,
both quicksort policies on a duplicate-containing input, and preservation
of the caller's array. Translations, monitor files, inputs, and
transcripts stay in \filepath{build/pcontrol-tests/}.
These results establish the small examples' values and completion,
not general performance, queue fairness, or faithful historical
parallel execution.


\subsection{Time Warp: Virtual Time and Recovering from Late Events}
\label{sec:timewarp}

Pcontrol changed the order in which pending computations run. Here we
also give application events simulated times and recover when their
arrival order disagrees with that history.
The \filepath{examples/time-warp/} directory implements a reflective
simulation experiment. A \emph{discrete event simulation} advances by
processing events such as a car arriving or a washer becoming available.
Each event has a simulated time. Processing an event at time 10 does
not require waiting ten real seconds.

Time Warp is an \emph{optimistic} approach: an object processes available
events without first proving that no earlier event will arrive later.
If one does arrive, the object must restore an earlier state and redo
work. This is \emph{rollback}, not an ordinary error retry: the computation
may have been correct for the events known then, but its history was
incomplete. Reflection lets this example put time and recovery policy
in the execution machinery rather than in every application handler.

\subsubsection{Local Time, Receive Time, and Three Queues}

\code{LVT}, the local virtual time, belongs to a simulation object's
metaobject. The custom evaluator makes that value available as the
pseudo-variable \code{LVT} in its scripts. A timed send has the form:

\begin{lstlisting}
[target <= message @ reply :vrt receive-time]
\end{lstlisting}

The evaluator sends the target metaobject a six-field envelope:

\begin{lstlisting}
[:message Message Reply Sender VST VRT]
\end{lstlisting}

\code{VST} is the sender's virtual send time; \code{VRT} is the
intended virtual receive time. Neither is a wall-clock timestamp.
The application still matches \code{Message}, not this entire envelope.
The explicit reply destination remains a separate field.

Distinguish three queues. A simulation metaobject has an ordinary
message queue and an undoable, time-ordered event queue.
Separately, the evaluator's own metaobject has a queue of evaluation
requests. Lowest-LVT-first versus FIFO changes this third queue,
not the timestamp semantics of the event queue.
An ordinary send without \code{:vrt} goes through the ordinary
message path; the source does not automatically timestamp every send.
For example, the washer's self-messages advance its work loop without
advancing virtual time.

\subsubsection{Run the Original Car-Wash Experiment}
\label{sec:tw-loading}

As with pcontrol, start in a scratch copy of the example directory:
\filepath{tw-load.abclr2} uses relative file names and creates
translations and compiled files. The automated check prepares the copy:

\begin{lstlisting}
make -C tutorial timewarp-check
\end{lstlisting}

In a fresh interpreter whose working directory is that copy:

\begin{codeblock}
\begin{lstlisting}
(inc "tw-load" :use-source t)
[starter <== [:start [3 4 6] 3]]
; Runs Lowest LVT first, then FIFO.
; Prints each run's RESULT and counters.
; Replies (:PRIORI (:ROLLBACK ...) (:EXPS ...)
;          :FIFO   (:ROLLBACK ...) (:EXPS ...)).
\end{lstlisting}
\end{codeblock}

\filepath{tw-load.abclr2} first compiles ten files, then incorporates
those same ten files in the same order. Compilation alone does not run
their object definitions (Section~\ref{sec:loading}). Read the loading
sequence as four stages, rather than opening an arbitrary helper first:

\begin{longtable}{L{0.44\textwidth}L{0.46\textwidth}}
\textbf{Files, in loader order} & \textbf{What the stage supplies}\\\hline
\endhead
\filepath{tw-undo-env}, \filepath{tw-undoable-pqueue},
\filepath{tw-output-history} & State frames, replayable input, and
speculative-output history for recovery.\\
\filepath{tw-pqueue}, \filepath{tw-monitor} & Evaluator work queues and
the rollback counter, distinct from the replayable input queue.\\
\filepath{tw-meta}, \filepath{tw-meta-eval}, \filepath{tw-eval} &
Application dispatch/recovery, scheduling of evaluator requests, and
expression evaluation with time/output context.\\
\filepath{tw-group}, \filepath{tw-wash1} & Connect the machinery, then
define the timed car-wash application and its controllers.\\
\end{longtable}

The names in this table omit \code{.abclr2}, just as the loader does.
It does not load the earlier \code{tw-wash0}, the independent
\code{tw-test}, or the unfinished \code{tw-scheduler} fragment.

The service times are 3, 4, and 6 simulated units; the last argument
requests three cars. \code{starter} compares scheduling policies,
whereas \code{starter-1} performs one simulation. Its factories
create a new attendant and new washers for each run.
The returned report contains counters, not the car assignment list.
That list is printed as \code{RESULT=(:FIN (...))}.
Do not assert particular expression counts as a language rule.

\filepath{tw-wash0.abclr2} is an earlier, ordinary-message version:
two queues pair waiting cars with waiting washers, and a two-signal
join completes the starter's request. Its own comment says it has
no rollback or Time Warp; its \code{service-time} argument is not
used to model a timed completion. The loader instead selects
\filepath{tw-wash1.abclr2}, which adds timed events and explicit
simulation completion.

\subsubsection{Follow Three Cars Through the Application}
\label{sec:tw-application}

\code{washers-gen} numbers the washers from zero. It initially
announces washer 0 at virtual time 0, washer 1 at time 1, and washer
2 at time 2. \code{starter-1} injects car arrivals at times 1, 2,
and 3, followed by an arrival sentinel \code{:fin} at time 4.
These initial injections go directly to the attendant's metaobject
with \code{:VRT}; they are not ordinary application messages to
the attendant.

The attendant stores pairs \code{[time car]} and \code{[time washer]}.
When both lists are nonempty, it removes their first entries and
starts washing at the maximum of those two times:

\wholelisting[firstline=79,lastline=87]{examples/time-warp/tw-wash1.abclr2}

This maximum prevents service before either resource is available.
\code{job-record} stores washer, car, and \emph{start} time, not finish
time. Each washer executes a short self-message loop to supply runtime
work, then announces its next availability at
\code{(+ service-time LVT)}:

\wholelisting[firstline=16,lastline=20]{examples/time-warp/tw-wash1.abclr2}

For the checked small run, the assignments are:

\begin{center}
\begin{tabular}{rrrr}
\textbf{Washer} & \textbf{Car} & \textbf{Start} & \textbf{Available again}\\\hline
0 & 1 & 1 & 4\\
1 & 2 & 2 & 6\\
2 & 3 & 3 & 9\\
\end{tabular}
\end{center}

This is a simulated-time trace, not the order in which console lines
must appear. \code{job-record} uses \code{cons}, so the printed
assignment list is in reverse recording order.
After the sentinel, the attendant sets \code{finished} and waits
until all washers are available before sending its result to
\code{starter-1}. The latter joins that result with its own
end-of-injection signal using \code{wait-flag}. The final reply
therefore does not mean merely that all arrival messages were sent.

\subsubsection{What Happens When an Earlier Event Arrives Late?}
\label{sec:tw-recovery}

Open \filepath{tw-meta.abclr2}. The ordinary path resembles the
metaobjects read earlier: enqueue, find a script, evaluate it,
and continue on \code{:end}. The timed path uses \code{:begin-tw}.
For a positive event with \code{VRT >= LVT}, it saves the current
state frame, sets \code{LVT} to \code{VRT}, and evaluates the
selected script with output-history and time added to the context.

Suppose an object has processed time 10 and then receives an event
for time 5. Continuing at 10 would violate the intended event
history. The other branch requests coordinated recovery:

\wholelisting[linerange={96-101,103-105}]{examples/time-warp/tw-meta.abclr2}

\begin{enumerate}
  \item The undoable environment restores a saved frame from before
        the offending receive time and returns its saved time.
  \item The input queue rewinds previously dequeued events at or after
        the cutoff so they can be considered again.
  \item The output history supplies sends made at or after that
        cutoff. \code{anti-messenger} sends matching cancellation
        envelopes to their destinations.
  \item Dispatch resumes, allowing the revised history to be processed.
\end{enumerate}

An \emph{anti-message} is a cancellation record, not an application
request to undo a particular business operation. The queue can remove
a matching positive event that has not yet executed. Cancellation of
already processed work can require the receiving object to roll back,
which can propagate recovery further.
The original queue compares message contents, reply, sender, and times;
it does not introduce a separate globally unique event identifier.
Read these as the experiment's concrete protocols, not a complete
general-purpose distributed simulation framework.

\subsubsection{Read the Three Histories Separately}
\label{sec:tw-histories}

In \filepath{tw-undo-env.abclr2}, \code{:push} saves a frame and
copies its association-list pairs before later binding updates.
\code{backFrame} selects a saved time strictly less than the rollback
cutoff. For a component trace, start with \code{x = 0}, save at 0,
set \code{x = 5}, save at 10, and set \code{x = 8}.
Rollback to 5 restores the frame at 0 and \code{x = 0}.
It does not interpolate a value between snapshots.

This is a shallow copy of binding pairs, not a deep copy of arbitrary
Lisp objects stored in them. Updating a binding through \code{:set}
is different from mutating an array or cons cell reachable through
that binding. Updates delegated to an outer lexical environment also
do not become part of the local frame's saved history.
Do not assume arbitrary host side effects are reversible.

\filepath{tw-undoable-pqueue.abclr2} keeps pending events in
\code{queue} and dequeued events in \code{old}. After processing
events at times 2 and 4, rollback to 3 makes the event at 4 available
again while retaining the earlier history. It also contains
positive/anti-message cancellation and sorting helpers.
Some alternative enqueue code is guarded by \code{\#+COMMENT};
read the active handlers rather than combining all printed alternatives.

\filepath{tw-output-history.abclr2} records the target, message,
reply destination, send time, and receive time. Its cutoff compares the
record's \emph{send time}, the fourth field, not its receive time.
For sends at times 1 and 5, a cutoff of 4 extracts the second send
even if its receive time is 20. Extraction removes those records
from the retained history. Restoring bindings alone would leave
such speculative output effects uncorrected.

\subsubsection{A Complete Two-Object Rollback Cascade}
\label{sec:tw-cascade}

The original three-car run does not force recovery. To see the original
machinery actually restore state, replay inputs, and cancel an output,
use the supplemental application in
\filepath{tutorial/examples/08-timewarp-rollback.abclr2}.
It introduces no new simulation algorithm: the unchanged original
loader, metaobjects, evaluator, environments, and history queues do the
recovery. The new application merely makes their effects easy to inspect.

The sink stores the latest received value and a reverse-order trace:

\wholelisting[firstline=6,lastline=12]{tutorial/examples/08-timewarp-rollback.abclr2}

The source adds each input to its total and sends that total to the sink
one simulated unit later:

\wholelisting[firstline=14,lastline=20]{tutorial/examples/08-timewarp-rollback.abclr2}

Neither application script contains an undo operation. Assignment goes
through the original environment protocol; the timed send is recorded
by the original evaluator. The trace uses \code{cons} to construct a new
list and rebinds \code{trace}, rather than destructively modifying a list
shared with earlier snapshots. That distinction matters for shallow
binding-history restoration.

\paragraph{Arrange a late event deliberately.}
Run \code{make -C tutorial timewarp-check}; the runner copies the
original directory and this supplemental file into
\filepath{build/timewarp-tests/}. For an interactive reproduction,
start a fresh interpreter in that prepared directory and load
\code{(inc "08-timewarp-rollback" :use-source t)}.
The example loads \code{tw-load} and resets the rollback monitor before
creating the two objects. It initially uses Lowest-LVT-first scheduling.

Inject a time-10 event, let the finite computation settle, and only
then inject a time-5 event:

\begin{codeblock}
\begin{lstlisting}
[[meta rollback-source] <= [:message [:add 10] nil nil 0 10]]
; Source total becomes 10 at LVT 10.
; Sink receives (:VALUE 10) at LVT 11.

[[meta rollback-source] <= [:message [:add 1] nil nil 0 5]]
; Source rolls back; its earlier output is cancelled at the sink.
; Final source total = 11, sink value = 11.
\end{lstlisting}
\end{codeblock}

These are deliberate meta-level event injections, not ordinary
application requests. The two \code{nil} fields are the initial reply
destination and sender; the final numbers are \code{VST = 0} and
\code{VRT = 10} or 5. No application reply is requested. The source's
subsequent timed sends name the source object as sender automatically.
The current interactive top-level drains runnable work after each
entered form. The checks exploit that boundary to ensure the first
output has \emph{already executed}, rather than merely waiting in the
sink's queue. This staged arrival order is a test arrangement, not a
promise that arbitrary concurrent clients see the same schedule.

\paragraph{Follow the histories, not just the final sum.}
After the first injection, the source has saved its initial frame at
time 0, set \code{LVT = 10}, and changed \code{total} from 0 to 10.
Its output history contains
\code{[rollback-sink [:value 10] nil 10 11]}.
The sink has likewise saved its initial frame at 0 before processing
time 11. Its value is 10 and its trace is \code{((11 10))}.

Now the late positive event has \code{VRT = 5 < LVT = 10}.
The source's \code{:begin-tw} handler takes the recovery branch.
It restores the saved frame at 0, so both its local time and total
return to zero. Its input queue has already dequeued the offending
time-5 event and the previous time-10 event. Rewinding to 5 makes both
available for replay, sorted by receive time. It must not discard the
late event after using it to trigger recovery.

The source extracts outputs whose send time is at least 5. The old
\code{[:value 10]} send qualifies because its send time was 10;
its receive time 11 is preserved in the cancellation envelope.
Here is the original anti-messenger's active handler:

\wholelisting[firstline=135,lastline=142]{examples/time-warp/tw-meta.abclr2}

The first pattern ends recursion on an empty history with \code{:fin}.
The second uses \code{bracket*} to separate one output record from the
remaining history. It sends directly to the target's metaobject,
changing the envelope's tag to \code{:anti-message}, but preserving
message, reply, sender identity, and both timestamps. Then it processes
the remaining records with the same \code{R}.

The source waits for the anti-messenger's completion before sending
itself \code{:end}. That completion means the cancellation requests
were issued, not that every receiving simulation object has completed
its own rollback. The target sends are asynchronous.
In this checked staged run, the sink receives:

\begin{codeblock}
\begin{lstlisting}
[:anti-message [:value 10] nil rollback-source 10 11]
\end{lstlisting}
\end{codeblock}

Because the positive event has already executed, it is in the sink's
\code{old} input history, not its pending queue. The anti-message
therefore cannot simply remove a pending positive event. The sink
dequeues the anti-message and rolls back to a frame strictly before
11: its initial frame at 0. This restores \code{value = 0} and
\code{trace = nil}. Rewinding the input history brings the old positive
and its anti-message together; \code{remove-anti} removes that pair,
so the invalid \code{[:value 10]} event is not replayed.
The sink has no timed outputs of its own to cancel in this application.

The source then processes its revised inputs at 5 and 10. At 5 it
adds 1 to zero and sends value 1 for receive time 6. At 10 it adds
10 to that corrected state and sends value 11 for receive time 11.
These are new outputs of replayed work, not the old output being
retained under another name. The sink's revised history is therefore:

\begin{center}
\begin{tabular}{rll}
\textbf{Receive time} & \textbf{Application event} & \textbf{Resulting sink value}\\\hline
6 & \code{[:value 1]} & 1\\
11 & \code{[:value 11]} & 11\\
\end{tabular}
\end{center}

The final reverse trace is \code{((11 11) (6 1))}; the speculative
\code{(11 10)} entry has disappeared. The source's final local time is
10, the sink's is 11, and the shared monitor records two rollbacks:
one caused by a late positive input, one caused by cancellation of
already processed output. Console cancellation and rollback lines
remain visible: restoring simulation bindings does not erase printing.

\paragraph{Inspect without injecting another application event.}
The original \code{TW-meta} does not implement an application-state
query. Its own ordinary metaobject can expose the simulation machinery's
bindings. For example:

\begin{codeblock}
\begin{lstlisting}
[source-meta-state := [[meta [meta rollback-source]] <== :state]]
[source-state := [source-meta-state <== [:value-of 'state]]]
[source-state <== [:value-of 'total]] ; => 11
[source-meta-state <== [:value-of 'LVT]] ; => 10
\end{lstlisting}
\end{codeblock}

The inner \code{meta} selects the simulation object's interpreter;
the outer selects that interpreter's ordinary metaobject.
Its \code{:state} reply is an environment whose binding named
\code{state} points to the undoable \emph{application} environment.
\code{LVT}, \code{pqueue}, and \code{output-history} instead belong
to the simulation metaobject's own frame. Do not confuse these two
state environments or treat the returned objects as copied snapshots.

\paragraph{Evidence and limits.}
Twelve assertions run for this cascade with each evaluator scheduler,
in separate fresh processes: Lowest-LVT-first and FIFO. They check
the speculative values and output record, corrected values and trace,
both final local times, the replacement output history, the replayed
source input history, the sink's retained corrected inputs, empty
pending event queues, and the two-rollback counter. Transcript
assertions require the cancellation and both restoration reports.
The finite result is the same under these two checked arrangements;
the test is not a performance comparison.

This closes the gap between isolated history tests and a concrete
metaobject/evaluator recovery cascade. It does not certify every
cancellation arrival order, simultaneous or duplicate event,
zero-time rollback boundary, or timestamped now-type send. In
particular, output-history matching is based on the original event
fields rather than unique event identifiers. The same original
\code{remove-anti} helper compares envelope tails, not the tags;
do not infer a general duplicate-delivery policy from this one pair.
The experimental NR/NS evaluator variants remain separate cases.

\paragraph{Check your understanding.}
Why is restoring only the source's total insufficient? Its speculative
output has already changed the sink. Why is the final value 11 rather
than 1? The time-10 input is replayed after the new time-5 input.
Why does the sink not end with both time-11 values in its trace?
The old input is cancelled during rewind and its earlier state effect
is rolled back. Does waiting for the anti-messenger join the entire
cascade? No; it joins cancellation issuance, not all recipient work.


\subsubsection{The Evaluator and a Source Reading Route}
\label{sec:tw-evaluator}

\filepath{tw-group.abclr2} connects \code{TW-meta-gen} and
\code{TW-eval-gen}. Scheduler replacement transfers the old
evaluator queue's entries to the new queue before installing it.
The queues in \filepath{tw-pqueue.abclr2} support \code{:enq},
\code{:deq}, \code{:listify}, and \code{[:set ...]} for that purpose.
The priority variant sorts by an entry's first field; the normal
variant appends entries for FIFO order.

\filepath{tw-meta-eval.abclr2} wraps evaluation requests as
\code{[LVT Message Reply Sender]} when the message has the
extended eight-field context; other requests receive key zero.
Its counter counts incoming evaluation requests, not cars,
virtual time units, or elapsed execution time.
\filepath{tw-monitor.abclr2} separately counts rollback announcements.

Read the first evaluator definition in \filepath{tw-eval.abclr2}
before its variants. It recognizes compiled expressions, variable
lookup, sends, object/group creation, sequential execution, and
argument-list evaluation. The added \code{Outputs} and \code{LVT}
fields carry the simulation context. A cache indexed by
\code{Id} supplies those fields to subsequent requests lacking
them. Timed past-type sends both construct a timestamped envelope
and record the output; ordinary sends use their separate path.

The same file also defines non-reifying and no-scheduling variants,
selected by the experimental \code{NR} and \code{NS} controllers
in \filepath{tw-wash1.abclr2}. They are not aliases for the two
policies compared by \code{starter}; Section~\ref{sec:tw-variants}
explains their separate setup and evaluator architecture. Timestamped now-type sends need separate
validation; the car-wash application uses timed past-type sends.

\filepath{tw-test.abclr2} supplies independent countdown objects
\code{foo}, \code{bar}, and \code{baz}. Their timed self-sends
increase virtual receive time by 1, 2, and 3 respectively.
Their \code{start} object's handler dispatches the three requests;
it has no final application reply. Do not call it with \code{<==}
and assume it is a completion join.
\filepath{tw-scheduler.abclr2} is an unfinished object-definition
fragment and is not loaded by \code{tw-load}. The operative scheduler
objects come from \code{tw-pqueue}, not that fragment.

\subsubsection{NR and NS: Replacing the Evaluator, Not Just Its Queue}
\label{sec:tw-variants}

The original \code{NR} and \code{NS} controllers in
\filepath{tw-wash1.abclr2} explore different implementations of the
Time Warp evaluator. They are not additional names for the priority
and FIFO policies already compared by \code{starter}.
The latter two replace the \emph{queue} of one schedulable evaluator;
NR and NS install entirely new evaluator objects.

\begin{longtable}{L{0.15\textwidth}L{0.36\textwidth}L{0.39\textwidth}}
\textbf{Variant} & \textbf{Evaluator's metaobject} & \textbf{Meaning of the change}\\\hline
\endhead
Standard & Created by \code{TW-meta-eval-gen}. & Its evaluation-request
queue can be selected as Lowest-LVT-first or FIFO.\\
NR & Explicit \code{non-reifying-meta}. & Evaluation requests use direct,
non-reifying handlers instead of the custom scheduling metaobject.\\
NS & Default ordinary metaobject. & Removes the custom Time Warp
evaluation-request scheduler, but keeps ordinary object dispatch.\\
\end{longtable}

All three evaluator definitions retain a local virtual-time cache and
the extended \code{Outputs}/\code{LVT} context. All are used with
\code{TW-meta-gen} for simulation objects, so selecting NR does not
remove those objects' rollback histories. ``No scheduling'' in NS
does not mean that the runtime stops scheduling processes or ignores
simulation event times. It refers to the omitted customized evaluator
scheduling layer; ordinary queues and the undoable event queue remain.

The group selectors show the replacement directly:

\wholelisting[firstline=22,lastline=31]{examples/time-warp/tw-group.abclr2}

Each creates a fresh evaluator and assigns the group's \code{evaluator}
binding. Unlike \code{:change-scheduler}, these handlers do not transfer
the old evaluator's pending request queue. Nor do they rewrite the
\code{evaluator} binding held by an existing simulation metaobject.
The checks create an object before switching and another afterwards:
the former retains the old evaluator, the latter captures the new one.
This is why the original NR/NS controller switches before
\code{starter-1} creates its attendant and washers. A group's configuration
is a creation policy, not necessarily a live broadcast to all its members.
Use separate fresh sessions for these experiments; restoring the group's
standard evaluator does not retrofit older simulation objects.

\paragraph{Read the controller's result protocol carefully.}
Here is NR's original controller:

\wholelisting[firstline=127,lastline=142]{examples/time-warp/tw-wash1.abclr2}

It selects the new evaluator, resets the rollback monitor, and waits on
the familiar \code{starter-1} simulation. If that request completes,
it prints the result, reads the rollback counter, and replies
\code{[:fifo [:rollback count]]}. NS follows the same structure but sends
\code{:no-schedule-eval}. Both reuse the historical \code{:fifo} result
tag, even though neither installs \code{normal-queue-gen} here. The tag
therefore cannot identify the selected evaluator architecture.
The expression-counter reset/query and \code{:exps} result fields are
commented out. The custom counter belongs to \code{TW-meta-eval}; an
ordinary or non-reifying evaluator metaobject does not acquire its
\code{:count} handler merely because the group has that query method.

The original entry points, after loading \code{tw-load} in a scratch
directory, are:

\begin{codeblock}
\begin{lstlisting}
; Use a fresh session for each variant.
[NR <== [:start [3 4 6] 3]]
; In another fresh session:
[NS <== [:start [3 4 6] 3]]
\end{lstlisting}
\end{codeblock}

Both complete this small case. Use the returned controller result to
identify completion; printed car assignments alone are not a join.

\paragraph{Follow the evaluator's representations.}
The translator represents a simple literal such as \code{:fin} as
\code{[:self-evaluate function]}, where the zero-argument function
returns the literal. Each evaluator invokes that function with
\code{funcall}; a function-valued payload is returned as data.
The standard evaluator shows the protocol directly:

\wholelisting[firstline=43,lastline=44]{examples/time-warp/tw-eval.abclr2}

The \code{[:script scriptObject]} case delegates to
\code{call-script-body}, which accepts internal scripts, object-level
\code{:exec}, or compiled functions. Trace the representation before
deciding which case handles a request.

\paragraph{What the variant checks establish.}
Run \code{make -C tutorial timewarp-variants-check}. Fifteen probes per
variant check a fresh evaluator identity, old/new objects' captured
evaluators, the literal's returned representation and payload,
a synthetic five-argument compiled body that replies 42, and
\code{LVT} lookup at 7. The compiled-body probe exercises
\code{:compiled}, whereas the literal probe exercises
\code{:self-evaluate}; one working case does not certify the other.
Further checks cover numbers, nil, symbols, lists, function-valued payloads,
single invocation of an initializer, and an application script returning 42.

Two car-wash sessions add four checks each: controller result shape,
historical \code{:fifo} tag, and a nonnegative rollback count. Transcripts
also require the \code{:fin} result and all three expected car assignments.
The 53 checks cover these evaluator and application protocols.
Inputs and transcripts stay in \filepath{build/timewarp-variants/}.
Use small runs to understand the control flow before exploring
additional cancellation orders or larger simulations.

\paragraph{Check your understanding.}
Does NR eliminate application-state rollback? No; simulation objects
still use the Time Warp metaobject generator. Does NS mean no queues?
No; it omits a particular scheduling layer. Does a returned \code{:fifo}
tag prove FIFO evaluator scheduling? No; the old controllers reuse that
tag. Can a new group evaluator fix an already created object's evaluator?
Not through these selectors: existing metaobjects retain their captured evaluator.


\paragraph{Check your understanding.}
If a washer starts at time 2 and has service time 4, does its self-loop
advance its virtual time ten times?
\emph{Answer:} No. The ordinary self-messages do not supply timestamps;
the next timed availability is at 6.
Why restore both input and output histories?
\emph{Answer:} Inputs may need replay after state restoration, while
outputs from discarded work may already have affected other objects.

\paragraph{What the checks establish.}
Forty checks run in isolated scratch sessions: twelve exercise
binding restoration, input replay, pending cancellation, and output
history extraction; four check the car-wash report and nonnegative
rollback counters; twenty-four check the two-object rollback cascade
under both evaluator scheduling policies. The transcript additionally checks the three
assignment records under both scheduling policies.
Files stay in \filepath{build/timewarp-tests/}.
The small car-wash case can complete without any rollback, so its
success does not demonstrate a complete rollback cascade.
Section~\ref{sec:tw-cascade} supplies a dedicated recovery integration
check. Other cancellation arrival orders remain outside these recovery
checks; the separate variant probes in Section~\ref{sec:tw-variants}
check small NR/NS runs, not arbitrary rollback schedules.


\subsection{Navigating the Original Sources}

Appendix~\ref{sec:source-map} maps every original example file and
library to its reading section and role. Use it to return to a helper
after following the corresponding application trace; it is not a new
instruction to load every file separately. The directory READMEs are
historical instructions, built-ins is a schematic protocol sketch, and
\code{tw-scheduler} is unfinished. Their presence in the source tree
does not make them independent executable demonstrations.


\section{Final Project}
\label{sec:project}

Implement a small concurrent job service with three roles:
\code{dispatcher}, \code{worker}, and \code{monitor}.

\begin{enumerate}
  \item The dispatcher accepts \code{[:run jobs]}, where each job is
        \code{[id value]}, creates a batch monitor, and assigns each job
        to a worker with \code{[:square id value]}.
  \item Each worker computes a square or rejects an invalid value,
        sending exactly one \code{[:done id status result]} to the monitor.
  \item The monitor counts successful and failed jobs and sends
        \code{[:report successful failed results]} to the original caller
        once the batch is complete. Handle an empty batch explicitly.
  \item Use past-type messages between actors and an explicit
        reply destination for a completion acknowledgment.
  \item Optional advanced extension: add a priority or express-message variant using an
        appropriate library and compare it with ordinary message handling.
\end{enumerate}

This project brings together object definitions, pattern matching,
state, replies, past-type and now-type messages, and concurrency.
The worked solution below covers this ordinary-message protocol;
the optional extension returns to reflective libraries and is not
required to complete it.

\subsection{A Worked Job-Service Protocol}

The complete solution is \filepath{tutorial/examples/03-jobs.abclr2}.
It uses two workers to show that completion cannot be inferred from
dispatch order. Each input is a two-element list \code{[id value]},
where \code{id} is an integer. Each worker sends exactly one
\code{[:done id status result]} reply, including when the value is invalid.
This rule lets the monitor count completions without confusing a rejected
job with a missing response.

\noindent\begin{minipage}{\linewidth}
\begin{tabular}{L{0.24\textwidth}L{0.40\textwidth}L{0.26\textwidth}}
\textbf{Receiver} & \textbf{Message} & \textbf{Reply destination}\\\hline
Dispatcher & \code{[:run jobs]} & The original caller.\\
Monitor factory & \code{[:new count destination]} & The dispatcher;
returns a fresh monitor object.\\
Worker & \code{[:square id value]} & The batch's monitor.\\
Monitor & \code{[:done id status result]} & After the last completion,
sends a report to the original caller.\\
\end{tabular}
\end{minipage}

The workers deliberately use the same behavior; the second definition
creates another independent object:

\wholelisting[firstline=2,lastline=14]{tutorial/examples/03-jobs.abclr2}

The monitor factory captures the batch's final destination.
Each created monitor has its own counters and result list:

\wholelisting[firstline=16,lastline=33]{tutorial/examples/03-jobs.abclr2}

The dispatcher creates the monitor before submitting any work.
It routes even job identifiers to one worker and odd identifiers
to the other, keeping \code{@ monitor} on every request:

\sourcelisting[firstline=35]{tutorial/examples/03-jobs.abclr2}

The caller waits on the dispatcher, but the dispatcher handler ends after
dispatching. The monitor sends the eventual reply to the caller's \code{R}.
The empty case replies immediately because there will be no worker
completions to trigger the monitor.

\begin{codeblock}
\begin{lstlisting}
(inc "tutorial/examples/03-jobs" :use-source t)
[dispatcher <== [:run [[1 2] [2 3] [3 :bad]]]]
; => (:REPORT 2 1 ((... completion entries ...)))
[dispatcher <== [:run nil]]
; => (:REPORT 0 0 NIL)
\end{lstlisting}
\end{codeblock}

The completion entries are \code{(1 :OK 4)}, \code{(2 :OK 9)}, and
\code{(3 :FAILED :NOT-A-NUMBER)}. Their order is not part of the interface.
The automated check sorts entries by identifier before comparing them.
The counts, identifiers, statuses, and computed values are what the
protocol promises.

\subsection{Invariants and Limits}

For every monitor, \code{remaining + successful + failed} equals the
number of submitted jobs. A reply is sent only when \code{remaining}
reaches zero. Because each worker responds exactly once in this example,
the final report means the whole batch has completed.

The demonstration assumes well-formed pairs with integer identifiers,
no duplicate completion messages, and no worker failures outside the
declared invalid-value branch. Production extensions would need
duplicate detection, timeout or cancellation policy, and validation
of malformed submissions. A timeout is a policy decision, not proof
that a remote computation stopped.

\section{Worked Exercise Solutions}

Run \filepath{07-solutions.abclr2} after attempting the exercises:

\begin{lstlisting}
(inc "tutorial/examples/07-solutions" :use-source t)
\end{lstlisting}

\subsection{Exercise 1: Sum of Squares}
\label{sol:squares}

\sourcelisting[firstline=3,lastline=9]{tutorial/examples/07-solutions.abclr2}

\begin{lstlisting}
[sum-of-squares <== [:sum [2 3 4]]]   ; => 29
[sum-of-squares <== [:sum nil]]       ; => 0
\end{lstlisting}

The temporary accumulator belongs to one invocation.
An empty \code{dolist} leaves it at its initial value, so the result
is zero without a special case.

\paragraph{Back to the original.}
In Section~\ref{sec:original-future}, \code{tree2leaves} also keeps
handler-local working data, but rewrites the remaining tree instead of
adding to an accumulator. Why is a fixed \code{cdr} step not enough there?
\emph{Answer:} a matched clause may expose a nested head and tail as new
work. Follow the rewritten tree, not just the original list's next element.

\subsection{Exercise 2: An Account Factory and a Transfer}
\label{sol:transfer}

The factory returns an anonymous object for each \code{:new} request.
Its initializer is copied into a separate state binding.
Initial balances must be nonnegative integers; this is a precondition,
not a check performed by the factory. Deposit and withdrawal amounts
are checked by the account handlers.

\wholelisting[firstline=3,lastline=27]{tutorial/examples/02-accounts.abclr2}

\begin{lstlisting}
(inc "tutorial/examples/02-accounts" :use-source t)
[alice <== [:transfer bob 30]]     ; => :TRANSFERRED
[alice <== :balance]               ; => 70
[bob <== :balance]                 ; => 50
[alice <== [:transfer alice 10]]   ; => :REJECTED
\end{lstlisting}

The self-transfer guard avoids a synchronous wait on the current object.
The demonstration uses trusted recipients and a sequential caller.
The recipient is credited before the sender is debited, and the two
changes are not atomic together. Opposing simultaneous transfers can
form a wait cycle. A production transfer protocol would need reservations
and a defined commit/cancel procedure; changing one send to past-type
alone does not solve that problem.

\paragraph{Back to the original.}
Why can the inherited feather query in Section~\ref{sec:original-inherit}
forward a destination, while this transfer waits for the deposit result?
\emph{Answer:} the inherited script can answer the original caller itself;
the transfer handler still needs the recipient's result before deciding
whether to debit its own state and report success.

\subsection{Exercise 3: Count Only After Acknowledgment}
\label{sol:count}

The solution creates a fresh consumer for each run. Every item returns
its post-increment count to a future, and the producer consumes exactly
as many acknowledgments as it submitted items:

\wholelisting[firstline=11,lastline=25]{tutorial/examples/07-solutions.abclr2}

\begin{lstlisting}
[counted-producer <== [:run [2 3 4]]]   ; => 3
[counted-producer <== [:run nil]]       ; => 0
\end{lstlisting}

The final count is queried only after all increments have been acknowledged.
The local consumer prevents another batch from affecting the result.
Each \code{next-value} may wait, but no worker waits back on the producer,
so this protocol does not introduce a circular dependency.

\paragraph{Back to the original.}
Why does this client count replies, while \code{print-leaves} waits for
\code{:finish}? \emph{Answer:} this protocol submits a known number of
items and requires exactly one reply per item. The original tree walker
sends a stream of leaves whose length the consumer does not know in advance.
Pcontrol's two-child join (Section~\ref{sec:pcontrol-fib}) is another
known-count protocol; it must collect two answers, not merely the first.

\subsection{Exercise 4: A Tagged Stream}
\label{sol:stream}

\wholelisting[firstline=27]{tutorial/examples/07-solutions.abclr2}

\begin{lstlisting}
[square-stream-client <== [:collect 1 4]]
; => (1 4 9 16)
\end{lstlisting}

Values arrive inside \code{[:value x]}, while \code{:done} is a distinct
terminal message. The client unwraps only the value messages.
This separation remains useful even if the data domain later includes
keywords or lists.

\paragraph{Back to the original.}
Would renaming \code{:finish} to another keyword in the original tree
walker solve its marker collision? \emph{Answer:} no; that keyword could
also occur as a leaf. The value wrapper distinguishes the \emph{role}
of a message instead of reserving one possible data value.

\subsection{Exercise 5: Unless at the Meta-Level}
\label{sol:unless}

Add this handler before the choice evaluator's fallback:

\sourcelisting[firstline=19,lastline=24]{tutorial/examples/05-reflection.abclr2}

When the predicate is true, reply to \code{C} with \code{nil}.
Otherwise evaluate the compiled body in its captured environment.
The test object's \code{:unless} handler invokes this protocol:

\begin{codeblock}
\begin{lstlisting}
[choice-test <== [:unless 3]]   ; => 3
[choice-test <== :selected]     ; => :UNLESS
[choice-test <== [:unless 8]]   ; => NIL
[choice-test <== :selected]     ; => :UNLESS (unchanged)
\end{lstlisting}
\end{codeblock}

The unchanged state verifies that the body was skipped for an even input.
Both control-flow paths still send a result to the caller's continuation.

\paragraph{Back to the original.}
Compare \code{my-if} (Section~\ref{sec:original-if}) and
\code{with-interrupt} (Section~\ref{sec:interrupt}). Where does each
reified request go? \emph{Answer:} \code{my-if}, like our choice form,
uses the evaluator; \code{with-interrupt} uses the application's metaobject.
For either route, trace the saved continuation separately from the
application handler's eventual \code{!} reply.


\section{Further Reading in the Source Tree}

This book is the primary user documentation for the current SBCL
port. The historical documents below provide supplementary detail and
background; their platform-specific installation instructions describe
the original release, not the current build.

\begin{itemize}
  \item \filepath{historical/doc/users-guide-e}: reference for ABCL/R2-specific
        forms and implementation restrictions.
  \item \filepath{doc/Masuhara_1992.pdf}: Hidehiko Masuhara.\newline
        \emph{Study on Implementation of an Object-Oriented Concurrent
        Reflective Language}, University of Tokyo, February 1992.
        Covers implementation, reflective towers, and lazy creation.
  \item \filepath{examples/misc/}: small programs demonstrating futures,
        express messages, built-in objects, and reflection.
  \item \filepath{lib/}: library implementations to inspect and load.
\end{itemize}

\appendix
\section{Building the Book and Checking Its Examples}
\label{sec:book-build}

From the repository root, build the PDF and check the examples:

\begin{lstlisting}[language=bash]
make -C tutorial
make -C tutorial check
\end{lstlisting}

The PDF build runs \code{lualatex} three times to stabilize the contents,
page numbers, and citations.
It keeps build intermediates in \filepath{build/tutorial-build/} and
copies the finished document to \filepath{doc/tutorial.pdf}.
The build and checks create their scratch directories automatically;
no pre-existing contents of \filepath{build/} are required.
It requires SBCL for source-excerpt preparation and a LuaLaTeX installation
with the packages named in the preamble,
Libertinus Serif and Libertinus Sans for text, and Noto Sans Mono for code.
The native runner \filepath{tutorial/tests/run-tests.lisp} needs SBCL and
the current ABCL/R2 executable, built with \code{make} at the root;
no Python or external Lisp packages are needed.
The checks run example sessions in fresh processes and keep source copies,
translations and transcripts under \filepath{build/}.
Session inputs are in \filepath{tutorial/tests/sessions/}; expected
results are specified in \filepath{tutorial/tests/cases.lisp}.
Run \code{make -C tutorial reading-check} for the Lisp-reading examples,
\code{reflection-check} for reflective protocols, and
\code{group-check} for ownership and identity checks.

Interactive loading creates a neighboring \code{.lsp} file; use
\code{:use-source t} while editing to avoid cached translations.
Printed object identifiers are session-dependent, not stable expected outputs.

\section{Original Source Reading Map}
\label{sec:source-map}

This map covers all 34 original example files and nine library files.
Each entry points to the explanation of its role in a complete example,
not necessarily to a separate runnable program. Start with the linked
application or loading section before reading its supporting generators.
Consult the walkthrough's stated checks and limits: a source-map entry
is not a certification that every execution path works on this image.

\newcommand{\sourcefile}[3]{\filepath{#1} & Section~\ref{#2}: #3\\}

\paragraph{Individual examples and reflective mechanisms.}
\begin{longtable}{L{0.41\textwidth}L{0.49\textwidth}}
\textbf{Original file} & \textbf{Reading route and role}\\\hline
\endhead
\sourcefile{examples/misc/future-ex.abclr2}{sec:original-future}{Tree traversal, repeated replies, and the terminal marker.}
\sourcefile{examples/misc/inherit.abclr2}{sec:original-inherit}{Bird/emu application, inherited scripts, and reflective context.}
\sourcefile{examples/misc/built-ins.abclr2}{sec:builtin-reading}{Schematic environment, script, and queue protocols; not a runtime to load.}
\sourcefile{examples/misc/user-script.abclr2}{sec:user-script-reading}{Sender-sensitive script matching, dynamic installation, and separate application/evaluation replies.}
\sourcefile{examples/misc/MCP.abclr2}{sec:mcp-reading}{Meta-circular group manager, application metaobjects, and evaluator cases.}
\sourcefile{examples/misc/meta-meta.abclr2}{sec:meta-meta-reading}{Three levels of identities and the twice computation.}
\sourcefile{examples/misc/reify-args.abclr2}{sec:original-if}{Eager/lazy branch selection; Section~\ref{sec:interrupt} follows interruption and the separate control-flow experiment.}
\sourcefile{examples/misc/fault-tolerant.abclr2}{sec:retry}{Acceptance state, initialization, rejection, and retry with preserved reply/sender.}
\sourcefile{examples/misc/break.abclr2}{sec:break}{Breakpoint continuation, lexical environment navigation, and resume.}
\sourcefile{examples/misc/express-ex.abclr2}{sec:original-express}{Newer clock protocol, suspension context, atomic regions, and express-message routing.}
\sourcefile{examples/misc/express-old-ex.abclr2}{sec:original-express}{Legacy clock and reader/writer lock participation; separate fresh session.}
\end{longtable}

\paragraph{Priority control: read the application, then the execution policy.}
\begin{longtable}{L{0.41\textwidth}L{0.49\textwidth}}
\textbf{Original file} & \textbf{Reading route and role}\\\hline
\endhead
\sourcefile{examples/pcontrol/README.abclr2}{sec:pcontrol-load}{Historical entry instructions, argument shapes, and small inputs.}
\sourcefile{examples/pcontrol/start.abclr2}{sec:pcontrol-load}{Compile applications, then load the four machinery files in dependency order.}
\sourcefile{examples/pcontrol/fib.abclr2}{sec:pcontrol-fib}{Child creation, saved reply destination, and two-answer join.}
\sourcefile{examples/pcontrol/qsort.abclr2}{sec:pcontrol-qsort}{Partition self-messages, shared array, and completion join.}
\sourcefile{examples/pcontrol/pcontrol-group.abclr2}{sec:pcontrol-policy}{Connect priority metaobjects and four evaluator objects into one group.}
\sourcefile{examples/pcontrol/priori-meta.abclr2}{sec:pcontrol-policy}{Attach object priority to evaluation and announce dispatch activity.}
\sourcefile{examples/pcontrol/pqueue-eval.abclr2}{sec:pcontrol-policy}{Priority/FIFO work queues, priority associations, and cyclic evaluator routing.}
\sourcefile{examples/pcontrol/para-monitor.abclr2}{sec:pcontrol-monitor}{Logical clock, active/object counters, and output streams.}
\end{longtable}

\paragraph{Time Warp: separate application, histories, and evaluator scheduling.}
\begin{longtable}{L{0.41\textwidth}L{0.49\textwidth}}
\textbf{Original file} & \textbf{Reading route and role}\\\hline
\endhead
\sourcefile{examples/time-warp/README.abclr2}{sec:tw-loading}{Historical run sequence and virtual-time application interface.}
\sourcefile{examples/time-warp/tw-load.abclr2}{sec:tw-loading}{Ten-file compile/load sequence; distinguishes loaded machinery from alternatives.}
\sourcefile{examples/time-warp/tw-wash1.abclr2}{sec:tw-application}{Timed car/washer pairing, availability, finish sentinel, and controller join.}
\sourcefile{examples/time-warp/tw-wash0.abclr2}{sec:tw-loading}{Earlier ordinary-message application, without Time Warp or timed service completion.}
\sourcefile{examples/time-warp/tw-meta.abclr2}{sec:tw-recovery}{Ordinary/timed dispatch, stragglers, state restoration, and anti-messages.}
\sourcefile{examples/time-warp/tw-undo-env.abclr2}{sec:tw-histories}{Saved local binding frames and strict-before-cutoff restoration.}
\sourcefile{examples/time-warp/tw-undoable-pqueue.abclr2}{sec:tw-histories}{Pending/dequeued inputs, replay, cancellation, and active versus conditional alternatives.}
\sourcefile{examples/time-warp/tw-output-history.abclr2}{sec:tw-histories}{Record sent events and extract speculative output by send-time cutoff.}
\sourcefile{examples/time-warp/tw-group.abclr2}{sec:tw-evaluator}{Connect components, transfer scheduler queues, and select evaluator variants.}
\sourcefile{examples/time-warp/tw-eval.abclr2}{sec:tw-evaluator}{Expression interpretation and time/output context; Section~\ref{sec:tw-variants} compares NR/NS.}
\sourcefile{examples/time-warp/tw-meta-eval.abclr2}{sec:tw-evaluator}{Queue evaluator requests by LVT and count evaluation arrivals.}
\sourcefile{examples/time-warp/tw-pqueue.abclr2}{sec:tw-evaluator}{Priority and FIFO work queues; not the undoable application input queue.}
\sourcefile{examples/time-warp/tw-monitor.abclr2}{sec:tw-evaluator}{Rollback announcements and their counter.}
\sourcefile{examples/time-warp/tw-test.abclr2}{sec:tw-evaluator}{Independent countdown objects; the starter is not a completion join.}
\sourcefile{examples/time-warp/tw-scheduler.abclr2}{sec:tw-evaluator}{Unfinished fragment, not loaded by the application loader.}
\end{longtable}

\paragraph{Libraries: loading is not the same as installing a policy.}
\begin{longtable}{L{0.41\textwidth}L{0.49\textwidth}}
\textbf{Original file} & \textbf{Reading route and role}\\\hline
\endhead
\sourcefile{lib/future.abclr2}{sec:original-future}{Buffered reply destination, consumption, and explicit application completion; see also Section~\ref{sec:futures}.}
\sourcefile{lib/multicast.abclr2}{sec:multicast}{Evaluator installation, target templates, forwarding, and reply collection.}
\sourcefile{lib/heavy-meta.abclr2}{sec:guards}{Guard-aware script selection and candidate fallback.}
\sourcefile{lib/heavy-meta-old.abclr2}{sec:guards}{Historical reference only: use the regular variant with the current translator.}
\sourcefile{lib/express.abclr2}{sec:original-express}{Newer sender/receiver machinery and execution modes.}
\sourcefile{lib/express-old.abclr2}{sec:original-express}{Different legacy lock/count protocol; not an alias for the newer library.}
\sourcefile{lib/copy.abclr2}{sec:copy}{New identities versus shared lexical/script/reply references and state-copy restrictions.}
\sourcefile{lib/watch.abclr2}{sec:watch}{Outgoing-send evaluator wrapper, reporters, and restoration boundaries.}
\sourcefile{lib/report.abclr2}{sec:report}{Copy-based arrival/dispatch/completion reports and filtering.}
\end{longtable}


\section{Syntax at a Glance}
\label{sec:syntax}
This is a reading aid, not a complete grammar. Names such as \code{target},
\code{pattern}, and \code{body} are placeholders; ellipses stand for omitted
parts. The linked sections explain evaluation, matching, prerequisites,
and the limits of the supplied implementation. Lisp notation and control
forms are introduced in Section~\ref{sec:lisp-reading}.

\paragraph{Objects, state, and matching.}
\begin{longtable}{L{0.39\textwidth}L{0.39\textwidth}L{0.12\textwidth}}
\textbf{Construct} & \textbf{Meaning} & \textbf{Section}\\\hline
\endhead
\code{[object name (script ...)]} & Create an object and bind its name;
omit the name for an anonymous object. & \ref{sec:objects}\\
\code{(state [n := 0])} & Object-definition entry for persistent state. & \ref{sec:state}\\
\code{(temporary [x := 0])} & Handler-local binding, not persistent state. & \ref{sec:objects}\\
\code{[x := expression]} & Change a binding. & \ref{sec:state}\\
\code{(=> pattern body)} & A script handling a matching message. & \ref{sec:objects}\\
\code{(=> pattern @ R from S body)} & Additionally bind reply destination
and sender; either header field can be omitted. & \ref{sec:objects}\\
\code{[head . tail]} & Pattern binding the head and entire tail of a pair;
dot notation is not an extra list element. & \ref{sec:lisp-reading}\\
\code{(match value (is pattern body) ...)} & Match a value inside a script;
\code{otherwise} supplies a fallback clause. & \ref{sec:objects}\\
\code{(=> pattern where test body)} & Guarded script selection, requiring
compatible heavy-meta machinery; not an ordinary \code{match} constraint. & \ref{sec:guards}\\
\end{longtable}

\paragraph{Sending, replying, and continuing.}
\begin{longtable}{L{0.39\textwidth}L{0.39\textwidth}L{0.12\textwidth}}
\textbf{Construct} & \textbf{Meaning} & \textbf{Section}\\\hline
\endhead
\code{[target <= message]} & Past-type send; do not wait for a reply. & \ref{sec:messages}\\
\code{[target <== message]} & Now-type send; wait for a reply. & \ref{sec:messages}\\
\code{[target <= message @ R]} & Past-type send with a specified reply destination. & \ref{sec:messages}\\
\code{[target <= message \$ f]} & Future-type send; route replies through
the loaded future library. & \ref{sec:futures}\\
\code{!expression} & Send a reply for the current message; not a return
from the rest of the handler. & \ref{sec:objects}\\
\code{[cont pattern body]} & Create a short, non-reifying continuation object. & \ref{sec:messages}\\
\code{[target <= message :express]} & Express send; requires both sender
and receiver machinery, not just this extra parameter. & \ref{sec:message-extensions}\\
\code{(atomic body ...)} & Receiver-local protection in the express
protocol; not a multi-object transaction. & \ref{sec:message-extensions}\\
\end{longtable}

\noindent\begin{minipage}{\linewidth}
\paragraph{Groups, identities, and reflection.}
\begin{longtable}{L{0.39\textwidth}L{0.39\textwidth}L{0.12\textwidth}}
\textbf{Construct} & \textbf{Meaning} & \textbf{Section}\\\hline
\endhead
\code{[group name ...]} & Create a group manager with optional custom
kernel objects. & \ref{sec:reflection}\\
\code{(group manager)} & Object-definition parameter selecting its
creation group. & \ref{sec:reflection}\\
\code{(meta-gen generator)} & Object/group definition parameter selecting
a metaobject generator. & \ref{sec:reflection}\\
\code{(evaluator evaluator-object)} & Group-definition parameter selecting
its evaluator. & \ref{sec:reflection}\\
\code{[meta object]} & Obtain its metaobject. & \ref{sec:reflection}\\
\code{[den metaobject]} & Obtain its lower-level denotation, not a copy. & \ref{sec:reflection}\\
\code{[group-of object]} & Query its group manager; custom generators must
provide the corresponding identity metadata. & \ref{sec:original-inherit}\\
\code{Me} / \code{group} & Current object denoted by \code{Id} / group
manager in \code{Gid}, not necessarily base-level application objects. & \ref{sec:reflection}\\
\code{(reify-args tag forms ...)} & Evaluate the tag, then send compiled
argument expressions and context to the current evaluator. & \ref{sec:reflection}\\
\code{(reify-args-to-meta tag forms ...)} & Send them to the current
object's metaobject instead. & \ref{sec:reflection}\\
\end{longtable}
\end{minipage}

\code{reify-args} and \code{reify-args-to-meta} are unavailable in a
non-reifying object's script. Extra send/object/assignment parameters
are interpreted by the selected machinery, not universally by their
spelling. Section~\ref{sec:reflection} explains the evaluation context;
Section~\ref{sec:files} separately collects loading and compilation commands.


\section{Troubleshooting Program Protocols}
\label{sec:failures}

\begin{description}
  \item[No return messages.]
  A now-type call has no reply. Check that a script matched and every
  result path either uses \code{!} or forwards the original reply
  destination. Also look earlier in the transcript for a Lisp error.
  \item[An object is not bound.]
  Its definition may have failed while loading the file. Read the first
  error rather than treating later missing bindings as independent failures.
  \item[A future never finishes.]
  The producer may not send the terminal marker, or a reader may have
  consumed it already. Futures do not infer completion.
  \item[A macro expands into an unsupported special form.]
  ABCL/R2 does not support every Common Lisp special form. Inspect the
  macro expansion and compare it with the supported forms following
  Section~\ref{sec:lisp-reading}.
  \item[Script lambdas.]
  Lambdas embedded in scripts cannot capture ABCL lexical variables,
  and their bodies cannot contain ABCL object creation or message sends
  in this implementation~\cite{guide}. Prefer a script object for
  an ABCL callback.
\end{description}

\clearpage
\section{Concept and Protocol Index}
\label{sec:concept-index}

Use this alphabetical index to return to an explanation, not as a list of
every occurrence. Entries give the section and its starting page. For
notation use Appendix~\ref{sec:syntax}; for original filenames use
Appendix~\ref{sec:source-map}.

\begingroup
\renewcommand{\arraystretch}{1.4}
\newcommand{\indexroute}[1]{\ref{#1} (p.~\pageref{#1})}
\begin{longtable}{@{}p{0.52\textwidth}p{0.43\textwidth}@{}}
\textbf{Concept or protocol} & \textbf{Explanation}\\\hline
\endhead
Actors and asynchronous messages & \indexroute{sec:messages}\\
Assignment & \indexroute{sec:state}\\
Atomic regions and their limits & \indexroute{sec:message-extensions}; \indexroute{sec:original-express}\\
Bindings, lexical scope, and Lisp notation & \indexroute{sec:lisp-reading}\\
Breakpoints and environment inspection & \indexroute{sec:break}\\
Built-in object protocols & \indexroute{sec:builtin-reading}\\
Callbacks and script lambdas & \indexroute{sec:objects}; \indexroute{sec:failures}\\
Compilation and cached translations & \indexroute{sec:loading}\\
Completion: application versus evaluation & \indexroute{sec:originals}\\
Continuations, \code{cont}, and explicit reply targets & \indexroute{sec:messages}\\
Copy: identity, state, and shared references & \indexroute{sec:copy}\\
Denotation, \code{den}, and \code{meta} & \indexroute{sec:reflection}\\
Dispatch and script selection & \indexroute{sec:objects}; \indexroute{sec:guards}\\
Environment frames and inspection & \indexroute{sec:reflection}; \indexroute{sec:break}\\
Evaluation context: Env, Id, Gid, Eval, C & \indexroute{sec:reflection}\\
Express messages and preemption & \indexroute{sec:message-extensions}; \indexroute{sec:original-express}\\
Fault tolerance: rejection and retry & \indexroute{sec:retry}\\
Fibonacci: recursive objects and reply joins & \indexroute{sec:pcontrol-fib}\\
Futures, \code{make-future}, and \code{next-value} & \indexroute{sec:futures}; \indexroute{sec:original-future}\\
Groups, \code{group-of}, and object creation & \indexroute{sec:reflection}; \indexroute{sec:original-inherit}\\
Guards, \code{where}, and heavy-meta & \indexroute{sec:guards}\\
Host events and the pending-request slot & \indexroute{sec:host-events}\\
Identity and metaobject towers & \indexroute{sec:meta-meta-reading}\\
Inheritance and forwarding reply destinations & \indexroute{sec:original-inherit}\\
Interrupting evaluation & \indexroute{sec:interrupt}\\
Job-service protocol: run, square, done, report & \indexroute{sec:project}\\
Loading: \code{inc}, \code{comp}, and \code{inc-lib} & \indexroute{sec:loading}\\
Library search path: \code{*library-directories*} & \indexroute{sec:library-path}\\
Macros and supported Lisp forms & \indexroute{sec:first-session}; \indexroute{sec:failures}\\
Matching, dotted patterns, and fallback & \indexroute{sec:objects}\\
MCP and the meta-circular interpreter & \indexroute{sec:mcp-reading}\\
Me and the currently executing object & \indexroute{sec:reflection}\\
Messages: past-, now-, and future-type & \indexroute{sec:messages}; \indexroute{sec:futures}\\
Modules: \code{require} and \code{provide} & \indexroute{sec:loading}\\
Monitor measurements and their limits & \indexroute{sec:pcontrol-monitor}\\
Multicast and positional reply collection & \indexroute{sec:multicast}\\
Non-reifying execution & \indexroute{sec:reflection}; \indexroute{sec:tw-variants}\\
Objects, factories, and scripts & \indexroute{sec:objects}; \indexroute{sec:account}\\
Priority queues and scheduling policy & \indexroute{sec:pcontrol-policy}\\
Quicksort as a message-driven state machine & \indexroute{sec:pcontrol-qsort}\\
Reflection: evaluator and metaobject routes & \indexroute{sec:reflection}; \indexroute{sec:original-if}\\
Reification: \code{reify-args}, \code{reify-args-to-meta} & \indexroute{sec:reflection}; \indexroute{sec:interrupt}\\
Replies, \code{!}, sender, and destination & \indexroute{sec:objects}; \indexroute{sec:originals}\\
Report: receiving-side observation and restoration & \indexroute{sec:report}\\
Rollback, anti-messages, and replay & \indexroute{sec:tw-recovery}; \indexroute{sec:tw-cascade}\\
Script installation and dynamic updates & \indexroute{sec:user-script-reading}\\
State versus temporary bindings & \indexroute{sec:state}; \indexroute{sec:objects}\\
Streams and terminal-message protocols & \indexroute{sec:original-future}; \indexroute{sol:stream}\\
Time Warp: virtual time and simulation & \indexroute{sec:timewarp}\\
Transactions: limits of an account transfer & \indexroute{sol:transfer}\\
Watch: outgoing-send observation and restoration & \indexroute{sec:watch}\\
\end{longtable}
\endgroup



\clearpage
\begin{thebibliography}{9}
\bibitem{guide}
Hidehiko Masuhara.
\emph{ABCL/R2 User's Guide}.
English translation, November 1992, with updates through July 1994.
Included as \filepath{historical/doc/users-guide-e}.
\bibitem{masuhara}
Hidehiko Masuhara.
\emph{Study on Implementation of an Object-Oriented Concurrent Reflective
Language}. Senior thesis, University of Tokyo, February 6, 1992.
Included as \filepath{doc/Masuhara_1992.pdf}.
\bibitem{distribution}
\emph{ABCL/R2 Distribution Readme and Installation Guide}.
Included as \filepath{historical/doc/readme} and
\filepath{historical/doc/install}, with the
current root Makefile providing the SBCL build.
\end{thebibliography}

\end{document}
